% WAFC -- suplemento, versao k = 4 (E5d; copia de k = 3, alteracoes marcadas em colR1; a
% rodada nao foi aceita). Template da Statistica Sinica
% (manuscript/ss-template/supp-temp_20240820.tex). Traz as provas dos
% resultados de ms_4.tex e os lemas auxiliares. As fontes sao
% derivations/01 a 08; o comentario acima de cada secao diz qual.
% Em k = 4 (D44 a D51): S5 a S7 passam ao block lasso balanceado (E1.12, E1.13),
% e o material do lasso coordenado que estava em S5 a S7 vai, com os enunciados
% que estavam no corpo, para a Secao S8 nova (D51(a)). O material deslocado nao
% e repetido em cinza no lugar antigo: uma linha cinza no inicio de cada secao
% diz o que saiu e para onde (ver docs/handoff-E5d.md).
% E5f (E1.15, L10; D54): a Proposition S6.2 (Corolario 15 do 08) no fim de S6, a
% Proposition S8.2 (Proposicao 8 do 05) em S8, com as provas; cada insercao em
% grupo ja azul leva o comentario % E5f.
% E5g (E1.9; D57): o Teorema 5 e o Lema 17 de derivations/09-cota-inferior.tex
% (Theorem S6.1, thm:lower; Lemma S6.3, lem:hypercube), com a prova, no lugar
% da Remark S6.1 (rem:lower), que fica em cinza (D51(b)) e vira uma frase; a
% uniformidade na prova do Theorem 1; cada insercao leva o comentario % E5g.
\documentclass[10pt]{book}
\usepackage[sectionbib]{natbib}
\usepackage{array,epsfig,fancyheadings,rotating}
\usepackage{hyperref}
%%%%%%%%%%%%%%%%%%%%%%%%%%%%%%%%%%%%%%%%%%%%%%%%%%%%%%%%%%%%%%%%%%%%%%%%%%%%%%

\textwidth=31.9pc
\textheight=46.5pc
\oddsidemargin=1pc
\evensidemargin=1pc
\headsep=15pt
\topmargin=.6cm
\parindent=1.7pc
\parskip=0pt

\usepackage{amsmath}
\usepackage{amssymb}
\usepackage{amsfonts}
\usepackage{amsthm}
\usepackage{enumitem}
\usepackage{xcolor}
\usepackage[normalem]{ulem}
\definecolor{colR1}{RGB}{0,90,200}

\setcounter{page}{1}
\newtheorem{theorem}{Theorem}[section]
\newtheorem{lemma}{Lemma}[section]
\newtheorem{corollary}{Corollary}[section]
\newtheorem{proposition}{Proposition}[section]
\theoremstyle{definition}
\newtheorem{definition}{Definition}[section]
\newtheorem{example}{Example}[section]
\newtheorem{remark}{Remark}[section]
\newtheorem*{remarkold}{Remark}   % E5g: a Remark S6.1 removida (D51(b))
\pagestyle{fancy}

% --- macros: as mesmas de ms_3.tex (docs/notacao.md) ---------------------
\newcommand{\R}{\mathbb{R}}
\newcommand{\E}{\mathrm{E}}
\newcommand{\Prob}{\mathrm{P}}
\newcommand{\Var}{\mathrm{Var}}
\newcommand{\T}{^{\top}}
\newcommand{\norm}[1]{\left\lVert #1 \right\rVert}
\newcommand{\normn}[1]{\left\lVert #1 \right\rVert_{n}}
\newcommand{\inner}[2]{\langle #1, #2 \rangle}
\newcommand{\cv}{\ell}
\newcommand{\md}{m}
\newcommand{\lev}{j}
\newcommand{\tsl}{k}
\newcommand{\jzero}{j_{0}}
\newcommand{\NJ}{N_{J}}
\newcommand{\bX}{\mathbf{X}}
\newcommand{\bU}{\mathbf{U}}
\newcommand{\bZ}{\mathbf{Z}}
\newcommand{\fcoef}[1]{\beta_{#1}}
\newcommand{\lvl}[1]{c_{#1}}
\newcommand{\gcomp}[2]{g_{#1#2}}
\newcommand{\wav}[2]{\psi_{#1#2}}
\newcommand{\VJ}{V_{J}}
\newcommand{\WJ}{W_{J}}
\newcommand{\PiJ}{\Pi_{J}}
\newcommand{\coef}[4]{\theta_{#1#2,#3#4}}
\newcommand{\coefs}[4]{\theta^{*}_{#1#2,#3#4}}
\newcommand{\btheta}{\boldsymbol{\theta}}
\newcommand{\thetastar}{\btheta^{*}}
\newcommand{\thetahat}{\widehat{\btheta}}
\newcommand{\bbeta}{\boldsymbol{\beta}}
\newcommand{\betastar}{\bbeta^{*}}
\newcommand{\betahat}{\widehat{\bbeta}}
\newcommand{\bc}{\mathbf{c}}
\newcommand{\chat}{\widehat{\mathbf{c}}}
\newcommand{\Gram}{\Sigma}
\newcommand{\Gramhat}{\widehat{\Sigma}}
\newcommand{\pen}{\lambda}
\newcommand{\penn}{\lambda_{n}}
\newcommand{\Szero}{S_{0}}
\newcommand{\Besov}[3]{B^{#1}_{#2,#3}}
\newcommand{\seff}{s'}
\newcommand{\Amat}{\mathbf{A}}
\newcommand{\Bmat}{\mathbf{B}}
\newcommand{\Btil}{\widetilde{\Bmat}}
\newcommand{\PA}{\Pi_{\Amat}}
\newcommand{\MA}{M_{\Amat}}
\newcommand{\bveps}{\boldsymbol{\varepsilon}}
\newcommand{\bb}{\mathbf{b}}
\newcommand{\bY}{\mathbf{Y}}
\newcommand{\bPsi}{\boldsymbol{\Psi}}
\newcommand{\smax}{\widehat{\sigma}_{\max}}
\newcommand{\gtil}{\widetilde{\gamma}}
\newcommand{\Mstar}{M^{\star}}
\newcommand{\lmin}{\lambda_{\min}}
\newcommand{\lmaxx}{\lambda_{\max}}
\newcommand{\opnorm}[1]{\left\lVert #1 \right\rVert_{\mathrm{op}}}
\newcommand{\compat}{\phi_{0}}
\newcommand{\GBA}{\Gramhat_{\Bmat\mid\Amat}}
\newcommand{\thbar}{\bar{\btheta}}
\newcommand{\bbar}{\bar{\bb}}
\newcommand{\thstarM}[1]{\btheta^{*,(#1)}}
\DeclareMathOperator{\vecop}{vec}
\DeclareMathOperator{\tr}{tr}
\DeclareMathOperator{\supp}{supp}
% --- a penalidade em blocos (docs/notacao.md, S10), as mesmas de ms_4.tex -----
\newcommand{\vbar}{\bar{v}}
\newcommand{\Gcal}{\mathcal{G}}
\newcommand{\Gzero}{\mathcal{G}_{0}}
\newcommand{\Gbar}{\bar{\mathcal{G}}}
\newcommand{\Hcal}{\mathcal{H}}
\newcommand{\Wcal}{\mathcal{W}}
\newcommand{\bn}{b_{n}}
\newcommand{\jst}{j^{*}}
\newcommand{\nGw}[1]{\lVert #1\rVert_{\mathcal{G},w}}
\newcommand{\nGone}[1]{\lVert #1\rVert_{\mathcal{G},1}}
\newcommand{\Psitil}{\widetilde{\Psi}}
\newcommand{\Rid}{R_{\mathcal{G}}}
\newcommand{\lamG}{\lambda^{\mathcal{G}}_{n}}
\newcommand{\lamGone}{\lambda^{\mathcal{G}}_{n,1}}
\newcommand{\lamGp}{\lambda^{\mathcal{G},+}_{n}}   % E1.14 (E5e): \lamG inflado acima do teto
\newcommand{\lamz}[1]{\lambda_{0,#1}}
\newcommand{\lamzbar}[1]{\bar{\lambda}_{0,#1}}
\newcommand{\rhoG}{\rho^{\mathcal{G}}_{n}}
\newcommand{\rhoGt}{\tilde{\rho}^{\mathcal{G}}_{n}}
\newcommand{\Tw}{\mathcal{T}_{\mathcal{G},w}}
\newcommand{\Ecal}{\mathcal{E}}
\newcommand{\ESig}{\mathcal{E}_{\Sigma}}
\newcommand{\Acal}{\mathcal{A}_{n}}
\newcommand{\AG}{A^{\mathcal{G}}_{s,\pi}}
\newcommand{\Dbl}{\Delta_{\cv\md}}
% E1.15 (E5f, D54): os pares em norma l1 e de grupos do risco ideal
% (docs/notacao.md, S10); \mathcal{R}, sem uso anterior.
\newcommand{\Rone}{\mathcal{R}_{1}}
\newcommand{\RGw}{\mathcal{R}_{\mathcal{G},w}}
\newcommand{\RGone}{\mathcal{R}_{\mathcal{G},1}}
% E5g (E1.9; D57): os simbolos da cota inferior que chegam ao manuscrito
% (docs/notacao.md, S10). O K = |K| do notacao.md fica escrito |\Kcal|, porque
% K ja e a dimensao de Psi(u) em S4 e uma constante em S6; os demais simbolos
% de derivations/09-cota-inferior.tex ficam locais, escritos no texto.
\newcommand{\Kcal}{\mathcal{K}}
\newcommand{\rlow}{\underline{\rho}_{n}}
\newcommand{\cA}{c_{\mathrm{A}}}
% E5g (D51(b)): a Remark S6.1 removida inteira fica no lugar, em cinza e sem
% numero, como os enunciados removidos do ms_4.tex.
\newcommand{\removedk}{\par\noindent[removed in k = 4]\par}

%%%%%%%%%%%%%%%%%%%%%%%%%%%%%%%%%%%%%%%%%%%%%%%%%%%%%%%%%%%%%%%%%%%%%%%%%%%%%%
\pagestyle{fancy}
\def\n{\noindent}
\lhead[\fancyplain{} \leftmark]{}
\chead[]{}
\rhead[]{\fancyplain{}\rightmark}
\cfoot{}

%%%%%%%%%%%%%%%%%%%%%%%%%%%%%%%%%%%%%%%%%%%%%%%%%%%%%%%%%%%%%%%%%%%%%%%%%%%%%%
\begin{document}

\renewcommand{\baselinestretch}{2}

\markright{ \hbox{\footnotesize\rm Statistica Sinica: Supplement
}\hfill\\[-13pt]
\hbox{\footnotesize\rm
}\hfill }

\markboth{\hfill{\footnotesize\rm AUTHOR(S)} \hfill}
{\hfill {\footnotesize\rm WAVELET ADDITIVE FUNCTIONAL COEFFICIENTS} \hfill}

\renewcommand{\thefootnote}{}
$\ $\par \fontsize{12}{14pt plus.8pt minus .6pt}\selectfont

%%%%%%%%%%%%%%%%%%%%%%%%%%%%%%%%%%%%%%%%%%%%%%%%%%%%%%%%%%%%%%%%%%%%%%%%%%%%%%

\centerline{\large\bf SPARSE WAVELET ESTIMATION OF}
\vspace{2pt}
\centerline{\large\bf ADDITIVE FUNCTIONAL COEFFICIENTS}
\vspace{.25cm}
\centerline{Author(s)}
\vspace{.4cm}
\centerline{\it Affiliation(s)}
\vspace{.55cm}
\centerline{\bf Supplementary Material}
\vspace{.55cm}
\fontsize{9}{11.5pt plus.8pt minus .6pt}\selectfont
\noindent
This supplement contains the proofs of the results stated in the paper and
the auxiliary lemmas they use. Section S1 fixes the notation and the objects.
Section S2 proves the identification results of Section 2 of the paper,
Section S3 the approximation lemma, Section S4 the bounds for the Gram matrix
of the product design, Section S5 the oracle inequality, \textcolor{gray}{\sout{and}} Section S6 the
% E5g (D57): a S6 passa a ter a cota inferior (Theorem S6.1).
rates \textcolor{colR1}{and a minimax lower bound}\textcolor{colR1}{, \textcolor{gray}{\sout{and}} Section S7 the recovery of the block structure by
thresholding} \textcolor{colR1}{and the risk of the thresholded estimator, and
Section S8 gives the theory of the coordinatewise lasso, the option of the
software that replaces the penalty of equation (2.5) by the $\ell_1$ norm}. References to numbered results without a section prefix are to the
paper; results numbered with the prefix S are stated and proved here.
\par

\setcounter{section}{0}
\setcounter{equation}{0}
\def\theequation{S\arabic{section}.\arabic{equation}}
\def\thesection{S\arabic{section}}

\fontsize{12}{14pt plus.8pt minus .6pt}\selectfont

%%%%%%%%%%%%%%%%%%%%%%%%%%%%%%%%%%%%%%%%%%%%%%%%%%%%%%%%%%%%%%%%%%%%%%%%%%%%%%
\section{Notation and objects}

The model, the basis and the estimator are those of Sections 2.1 to 2.3 of the
paper. We write $\bZ = [\Amat\;\Bmat]$ for the design matrix of equation
(2.4), with $\Amat \in \R^{n\times p}$ the unpenalized block and
$\Bmat \in \R^{n\times d}$, $d = pq\NJ$, the block of products; $\bY$, $f$ and
$\bveps$ for the vectors with entries $Y_i$, $f(\bX_i,\bU_i)$ and
$\varepsilon_i$; $\Gramhat = \bZ\T\bZ/n$ and $\Gram = \E(\bZ\bZ\T)$; and, for
$a, h \in \R^n$, $\normn{a}^2 = n^{-1}\sum_i a_i^2$ and
$\inner{a}{h}_n = n^{-1}\sum_i a_ih_i$.

The \textcolor{colR1}{oracle of the approximation space} is $\betastar = (\bc,\thetastar)$, with
$\coefs{\cv}{\md}{\lev}{\tsl} = \inner{\gcomp{\cv}{\md}}{\wav{\lev}{\tsl}}$
for $\lev < J$, so that $f_J = \bZ\betastar$ is the best approximation of $f$
in the \textcolor{colR1}{approximation space} in the sense of Lemma 2; the deterministic residual is
$\bb = f - f_J$, the oracle support is $\Szero = \supp(\thetastar)$ and
$s_0 = |\Szero|$, and $v = \thetahat - \thetastar$. When $\Amat$ has rank $p$
we write $\PA = \Amat(\Amat\T\Amat)^{-1}\Amat\T$, $\MA = I_n - \PA$,
$\Btil = \MA\Bmat$, $\GBA = \Btil\T\Btil/n$ and $\gtil = \lmin(\GBA)$, and
finally $\smax^2 = \max_a(\Gramhat)_{aa}$, the maximum being over the $d$
penalized columns.

{\color{colR1}
For the penalty of equation (2.5), $\Gcal$ is the partition of the $d$
penalized coordinates into the chunks of Section 2.3 of the paper, $|\Gcal|$
the number of chunks, $\btheta_{G}$ and $|G|$ the subvector and the size of a
chunk $G$, $w_{G} > 0$ its weight and $\varrho = \max_{G}w_{G}/\min_{G}w_{G}$
the ratio of the weights; $\nGw{\btheta} = \sum_{G}w_{G}\norm{\btheta_{G}}_{2}$,
$\nGone{\btheta}$ is the same norm with unit weights, and, for
$\Hcal \subseteq \Gcal$, $\nGw{\btheta_{\Hcal}} =
\sum_{G\in\Hcal}w_{G}\norm{\btheta_{G}}_{2}$ and
$\Wcal(\Hcal) = \sum_{G\in\Hcal}w_{G}^{2}$; $\Gzero(\btheta) =
\{G : \btheta_{G} \ne 0\}$ are the active chunks of $\btheta$. With unit
weights $\Wcal(\Gzero(\btheta))$ is the number of active chunks, and with
$w_{G} = |G|^{1/2}$ it is the number of coordinates in them. A comparator is
any $\bar\bbeta = (\bar\bc,\thbar) \in \R^{p}\times\R^{d}$, with
deterministic residual $\bbar = f - \bZ\bar\bbeta$, active chunks
$\Gbar = \Gzero(\thbar)$, $\bar\Wcal = \Wcal(\Gbar)$ and
$\vbar = \thetahat - \thbar$. For a chunk, $\Bmat_{G}$ and $\Btil_{G}$ are the
columns of $\Bmat$ and $\Btil$ in $G$, $\Psitil_{G} = \Btil_{G}\T\Btil_{G}/n$
and $\Gramhat_{GG} = \Bmat_{G}\T\Bmat_{G}/n$. The ideal risk over chunks is
$\Rid(\btheta;\eta) = \sum_{G}\min(\norm{\btheta_{G}}_{2}^{2},\eta)$, for
$\eta > 0$; it is nondecreasing in $\eta$, and, since
$\min(u,\kappa\eta) \le \kappa\min(u,\eta)$ for $\kappa \ge 1$,
$\Rid(\btheta;\kappa\eta) \le \kappa\Rid(\btheta;\eta)$. Finally
$\gamma = \kappa_{1}c_{U}$, $\Lambda = \kappa_{2}C_{U} + \gamma/2$ is the
constant of equation (3.1), and
\[
  \ESig \;=\; \bigl\{\opnorm{\Gramhat - \Gram} < \gamma/2\bigr\};
\]
on $\ESig$, Weyl's inequality and Proposition 2(i) give
$\lmin(\Gramhat) > \gamma/2$ and $\lmaxx(\Gramhat) < \Lambda$, and the matrix
Bernstein bound in the proof of Proposition 2(ii), Section S4, gives
$\Prob(\ESig) \to 1$ whenever $pq2^{J}\log(pq2^{J})/n \to 0$.
}

For a support $S$ and a symmetric matrix $A$, the compatibility constant is
\[
  \compat^2(S;A) = \min\bigl\{|S|\,\delta\T A \delta :\;
  \norm{\delta_S}_1 = 1,\; \norm{\delta_{S^c}}_1 \le 3 \bigr\},
\]
and the elementary bound
$\delta\T A\delta \ge \lmin(A)\norm{\delta}_2^2 \ge
\lmin(A)\norm{\delta_S}_1^2/|S|$ gives
\begin{equation}
  \label{eq:compat-lmin}
  \compat^2(S;A) \;\ge\; \lmin(A) \qquad \text{for every } S,
\end{equation}
so that a lower bound on the smallest eigenvalue is a compatibility bound on
every support, with the same constant and no cone restriction.

%%%%%%%%%%%%%%%%%%%%%%%%%%%%%%%%%%%%%%%%%%%%%%%%%%%%%%%%%%%%%%%%%%%%%%%%%%%%%%
% derivations/01-identificabilidade.md
\section{Identification and the periodized basis}
\setcounter{equation}{0}

\subsection*{Proof of Proposition 1}

\emph{Levels and components are determined by the coefficients.} Let $\beta$
and $\tilde\beta$ be measurable maps from $[0,1]^q$ to $\R^p$ with
$\sum_{\cv}\fcoef{\cv}(\bU)X_{\cv} = \sum_{\cv}\tilde\beta_{\cv}(\bU)X_{\cv}$
almost surely, and put $\delta(u) = \beta(u)-\tilde\beta(u)$. Then
$\{\delta(\bU)\T\bX\}^2 = 0$ almost surely, and since $\delta(\bU)$ is
measurable with respect to $\bU$ and $\E\norm{\bX}_2^2 < \infty$,
\[
  0 \;=\; \E\bigl[\{\delta(\bU)\T\bX\}^2 \mid \bU\bigr]
  \;=\; \delta(\bU)\T\,\E(\bX\bX\T\mid\bU)\,\delta(\bU)
  \qquad \text{almost surely.}
\]
The conditional second moment matrix is positive definite almost surely, and a
positive definite quadratic form vanishes only at the origin, so
$\delta(\bU) = 0$ almost surely.

\emph{Levels and components are determined within a coefficient.} Because the
density of $\bU$ is positive Lebesgue almost everywhere, a set is null for the
law of $\bU$ if and only if it is Lebesgue null, so the identity
$\lvl{\cv} + \sum_{\md}\gcomp{\cv}{\md}(u_{\md}) =
\tilde{c}_{\cv} + \sum_{\md}\tilde{g}_{\cv\md}(u_{\md})$ holds Lebesgue almost
everywhere on $[0,1]^q$. Write $\Delta = \lvl{\cv}-\tilde{c}_{\cv}$ and
$\delta_{\md} = \gcomp{\cv}{\md}-\tilde{g}_{\cv\md} \in L_2[0,1] \subset
L_1[0,1]$, so that $\Delta + \sum_{\md}\delta_{\md}(u_{\md}) = 0$ almost
everywhere. Each summand is integrable on the cube by Tonelli's theorem.
Integrating in all coordinates and using $\int_0^1\delta_{\md} = 0$, which
holds because both parametrizations satisfy (2.2), gives $\Delta = 0$; fixing
$\md$ and integrating in the remaining coordinates gives
$\delta_{\md}(u_{\md}) = 0$ for almost every $u_{\md}$. \qed

\begin{remark}
\label{rem:centering}
If the centering $\int_0^1 \gcomp{\cv}{\md} = 0$ is replaced by
$\E\{\gcomp{\cv}{\md}(U_{\md})\} = 0$, the same proof applies under the
two sided density bound of Assumption 3, since then $L_2[0,1]$ and
$L_2(P_{U_{\md}})$ coincide as sets: the step that uses Fubini's theorem
yields $\delta_{\md}$ constant almost everywhere, and the constant is its
mean, which is zero. The two parametrizations are related by
$\lvl{\cv} \mapsto \lvl{\cv} + \sum_{\md}\E\{\gcomp{\cv}{\md}(U_{\md})\}$, as
stated in Section 2.1 of the paper. The difference is not cosmetic for the
estimator: the design of Section 2.3, which omits the constant scaling
function, estimates the version centered with respect to Lebesgue measure.
\end{remark}

\subsection*{Proof of Lemma 1}

\emph{The basis.} The periodization of the scaling function is
$\phi^{\rm per}_{00}(u) = \sum_{\iota}\phi(u+\iota)$, a finite sum at each $u$
because the support is compact, and it equals the constant function one, since
the integer translates of an orthonormal scaling function with
$\int\phi = 1$ form a partition of unity \textcolor{colR1}{\citep[Section~7.5.1]{Mallat-2009}}. For a wavelet, exchanging the finite
sum with the integral,
$\int_0^1\wav{\lev}{\tsl} = \int_{\R}2^{\lev/2}\psi(2^{\lev}y-\tsl)\,dy =
2^{-\lev/2}\int_{\R}\psi = 0$ by the vanishing moment of order zero. That
$\{\mathbf{1}\}\cup\{\wav{\lev}{\tsl}\}$ is an orthonormal basis of $L_2[0,1]$
is the standard construction of the periodized multiresolution analysis
\citep[Section 9.3]{Daubechies-1992}. Since $\dim\VJ = 2^J$ and the $2^J-1$
wavelets of level below $J$ are orthogonal to the constant, $\WJ$ is exactly
the set of elements of $\VJ$ with zero integral.

\emph{Injectivity.} Suppose two parameter vectors give the same regression
function almost surely. The functions
$\beta^J_{\cv}(u) = \lvl{\cv}+\sum_{\md}g^J_{\cv\md}(u_{\md})$, with
$g^J_{\cv\md} = \sum_{\lev\tsl}\coef{\cv}{\md}{\lev}{\tsl}\wav{\lev}{\tsl}$,
are measurable and bounded, so the first part of Proposition 1 gives
$\beta^J_{\cv} = \tilde\beta^J_{\cv}$ almost surely; each $g^J_{\cv\md}$ has
zero integral by the previous paragraph, so the second part gives
$\lvl{\cv} = \tilde{c}_{\cv}$ and
$\sum_{\lev\tsl}(\btheta-\tilde\btheta)_{\cv\md,\lev\tsl}\wav{\lev}{\tsl} = 0$
almost everywhere, whence $\btheta = \tilde\btheta$ by orthonormality. For the
Gram matrix, $\E\norm{\bZ}_2^2 < \infty$ because $\E\norm{\bX}_2^2 < \infty$
and each wavelet is bounded, and for $a \ne 0$,
$a\T\Gram a = \E\{(a\T\bZ)^2\} = \norm{f_a}^2_{L_2(P)} > 0$ by injectivity,
where $f_a$ is the regression function with parameter $a$.

\emph{Nullity when the scaling column is kept.} With $\phi_{00}\equiv 1$ the
column $X_{\cv}\phi_{00}(U_{\md})$ coincides with the unpenalized column
$X_{\cv}$, for each of the $pq$ pairs, so the $pq$ vectors that oppose one to
the other lie in the kernel and are linearly independent. There are no others:
a kernel vector defines
$\sum_{\cv}X_{\cv}\{a_{\cv} + \sum_{\md}(\alpha_{\cv\md} +
\sum_{\lev\tsl}a_{\cv\md,\lev\tsl}\wav{\lev}{\tsl}(U_{\md}))\} = 0$ almost
surely; by Proposition 1 applied to the part with zero integral,
$a_{\cv\md,\lev\tsl} = 0$ and $a_{\cv}+\sum_{\md}\alpha_{\cv\md} = 0$, which
is the space of dimension $pq$ spanned by those vectors. \qed

\begin{remark}[The interval basis]
\label{rem:cdv}
With the basis of \citet*{cohen1993wavelets} the wavelets of level
$\lev \ge \jzero$ are orthogonal to $V_{\jzero}$, which contains the
constants, so they still have zero integral and the penalized part of each
block satisfies (2.2) with no restriction. The $2^{\jzero}$ scaling functions
do not: the constant is the combination $1 = \sum_{\tsl}\mu_{\tsl}
\phi^{\rm int}_{\jzero\tsl}$ with
$\mu_{\tsl} = \int_0^1\phi^{\rm int}_{\jzero\tsl}$, and the subspace with zero
integral is $\{\mu\T\alpha = 0\}$. Keeping the $2^{\jzero}$ scaling columns of
each block unrestricted repeats the constant and makes $\Gram$ singular with
nullity $pq$, as in the last part of Lemma 1; reparametrizing
$\alpha_{\cv\md} = Q\gamma_{\cv\md}$ with $Q$ an orthonormal basis of
$\{\mu\T\alpha = 0\}$ restores injectivity, at the price of
$pq(2^{\jzero}-1)$ unpenalized parameters.
\end{remark}

%%%%%%%%%%%%%%%%%%%%%%%%%%%%%%%%%%%%%%%%%%%%%%%%%%%%%%%%%%%%%%%%%%%%%%%%%%%%%%
% derivations/02-aproximacao-besov.tex
\section[\texorpdfstring{Approximation by the wavelet spaces $\WJ$}{Approximation by the wavelet spaces W\_J}]{Approximation by the wavelet \texorpdfstring{\textcolor{colR1}{spaces $\WJ$}}{spaces W\_J}}
\setcounter{equation}{0}

\subsection*{Proof of Lemma 2}

Fix a pair $(\cv,\md)$ and write $g = \gcomp{\cv}{\md}$ and
$\theta_{\lev\tsl} = \coefs{\cv}{\md}{\lev}{\tsl}$.

\emph{Tail sum.} By Step 2 below the coefficients are square summable, so
$g \in L_2[0,1]$ and, the system being complete and $g$ orthogonal to the
constant, $g = \sum_{\lev\ge0}\sum_{\tsl}\theta_{\lev\tsl}\wav{\lev}{\tsl}$ in
$L_2$. The orthogonal projection onto $\WJ$ retains the levels below $J$, and
Parseval's identity gives
$\norm{g-\PiJ g}^2_{L_2[0,1]} = \sum_{\lev\ge J}
\norm{\theta_{\lev\cdot}}_2^2$.

\emph{From $\ell_\pi$ to $\ell_2$ inside a level.} In $\R^{2^{\lev}}$ one has
$\norm{a}_2 \le \norm{a}_{\pi}$ when $\pi \le 2$, and
$\norm{a}_2 \le (2^{\lev})^{1/2-1/\pi}\norm{a}_{\pi}$ by Hölder's inequality
when $\pi > 2$. In both cases $\norm{\theta_{\lev\cdot}}_2 \le
2^{\lev(1/2-1/\pi)_{+}}\norm{\theta_{\lev\cdot}}_{\pi}$, and
Assumption 5 gives
\[
  \norm{\theta_{\lev\cdot}}_2 \;\le\;
  C_g\,2^{-\lev\{s+1/2-1/\pi-(1/2-1/\pi)_{+}\}}
  \;=\; C_g\,2^{-\lev\seff},
\]
since the exponent equals $s$ when $\pi \ge 2$ and $s-(1/\pi-1/2)$ when
$\pi < 2$. This is where the loss $(1/\pi-1/2)_{+}$ appears.

\emph{Geometric sum and density.}
$\sum_{\lev\ge J}C_g^2 2^{-2\lev\seff} = C_g^2 2^{-2J\seff}/(1-2^{-2\seff})$,
finite because $\seff > 0$, which is the first display of Lemma 2. For the
second, write $f-f_J = \sum_{\cv,\md}X_{\cv}h_{\cv\md}(U_{\md})$ with
$h_{\cv\md} = \gcomp{\cv}{\md}-\PiJ\gcomp{\cv}{\md}$, the levels cancelling.
By the Cauchy-Schwarz inequality over the $pq$ terms and
$|X_{\cv}| \le B_X$,
\[
  \E\{(f-f_J)^2\} \;\le\; pq\sum_{\cv,\md}
  \E\{X_{\cv}^2h_{\cv\md}^2(U_{\md})\}
  \;\le\; pq\,B_X^2\sum_{\cv,\md}\E\{h_{\cv\md}^2(U_{\md})\},
\]
and $\E\{h^2(U_{\md})\} = \int_0^1 h^2 p_{U_{\md}} \le C_U\norm{h}^2_{L_2[0,1]}$,
the marginal density being bounded by $C_U$ because integrating the joint
density over the other $q-1$ coordinates of $[0,1]^{q-1}$ does not increase
the bound. One factor $pq$ comes from the Cauchy-Schwarz inequality and the
other from the sum. \qed

The uniform version is used in Section S5 to bound the empirical norm of the
bias without an appeal to Markov's inequality.

\begin{lemma}[Uniform approximation error]
\label{lem:sup}
Under Assumptions 4 and 5 with $s > 1/\pi$, each component has a continuous
version and, for every $J \ge 0$,
\[
  \norm{\gcomp{\cv}{\md}-\PiJ\gcomp{\cv}{\md}}_{\infty}
  \;\le\; \frac{(L-1)\norm{\psi}_{\infty}C_g}{1-2^{-(s-1/\pi)}}\,
  2^{-J(s-1/\pi)} ,
\]
and therefore
$\normn{\bb} \le \norm{f-f_J}_{\infty} \le pq B_X(L-1)\norm{\psi}_{\infty}
C_g\,2^{-J(s-1/\pi)}/\{1-2^{-(s-1/\pi)}\}$ for every sample.
\end{lemma}

\begin{proof}
At a point $u$ and a level $\lev$ at most $L-1$ periodized translates are
nonzero, because $\psi(2^{\lev}u-\tsl)\ne0$ requires
$2^{\lev}u-\tsl \in (0,L-1)$, so
$\sum_{\tsl}|\wav{\lev}{\tsl}(u)| \le (L-1)2^{\lev/2}\norm{\psi}_{\infty}$.
With $\max_{\tsl}|\theta_{\lev\tsl}| \le \norm{\theta_{\lev\cdot}}_{\pi} \le
C_g2^{-\lev(s+1/2-1/\pi)}$ the level contributes at most
$(L-1)\norm{\psi}_{\infty}C_g2^{-\lev(s-1/\pi)}$ in the supremum norm. Since
$s > 1/\pi$ the series converges absolutely and uniformly, so its sum is
continuous and agrees almost everywhere with $g$; the tail from level $J$ on
is the geometric sum of the displayed bounds, and the statement for $f-f_J$ is
the triangle inequality over the $pq$ terms.
\end{proof}

The next proposition is the one referred to in Section 2.2 of the paper: it
quantifies what the periodized basis charges for a component that does not
match at the two ends of the interval.

\begin{proposition}[Cost of periodization]
\label{prop:periodization}
Let the basis be as in Assumption 4 and let
$g \in \Besov{s}{\pi}{r}[0,1]$, the Besov space of the interval and not of the
circle, with $s > 1/\pi$, $N > s$ and $\int_0^1 g = 0$. Then there is a
constant $C$ depending on $g$, on $\psi$ and on $(s,\pi)$ such that, for every
$J$ with $2^J \ge L-1$,
\[
  \norm{g-\PiJ g}^2_{L_2[0,1]} \;\le\; C\,2^{-2J\seff}
  \;+\; 2(L-1)\norm{\psi}_1^2\norm{g}^2_{\infty}\,2^{-J},
\]
so that the $L_2$ rate is $2^{-J\min(\seff,1/2)}$. If $g$ is continuous with
$g(0) \ne g(1)$, the second term is of the exact order $2^{-J}$ and
$\norm{g-\PiJ g}_{\infty}$ does not converge to zero.
\end{proposition}

\begin{proof}
Let $\tilde g$ be the one periodic extension of $g$, so that
$\theta_{\lev\tsl} = \int_{\R}\tilde g(y)2^{\lev/2}\psi(2^{\lev}y-\tsl)\,dy$.
At a level with $2^{\lev} \ge L-1$, separate the interior translates, whose
support lies inside $[0,1]$, from the at most $L-1$ boundary translates that
cross the identification of $0$ with $1$. For an interior translate,
$\theta_{\lev\tsl} = \inner{Eg}{\psi_{\lev\tsl}}_{\R}$ for any extension $Eg$
of $g$ to the line; taking $E$ a bounded extension operator from
$\Besov{s}{\pi}{r}[0,1]$ to $\Besov{s}{\pi}{r}(\R)$
\citep[Chapter 3]{Triebel-1983} and using the wavelet characterization of the
latter, which for the implication from the function to the sequence requires
only $N > s$ vanishing moments, gives
$\norm{(\theta_{\lev\tsl})_{\rm int}}_{\pi} \le C_12^{-\lev(s+1/2-1/\pi)}$,
whose tail sum is bounded as in Lemma 2. For a boundary translate,
$s > 1/\pi$ makes $g$ bounded and
$|\theta_{\lev\tsl}| \le \norm{g}_{\infty}\norm{\psi_{\lev\tsl}}_1 =
\norm{g}_{\infty}\norm{\psi}_12^{-\lev/2}$, and there are at most $L-1$ of
them per level, so their tail sum is at most
$2(L-1)\norm{\psi}_1^2\norm{g}_{\infty}^22^{-J}$. For the exact order, if $g$
is continuous with $g(0) \ne g(1)$ then on the support of a boundary translate
$\tilde g$ equals $g(1)+o(1)$ to the left of the identification point and
$g(0)+o(1)$ to the right, so
$\theta_{\lev\tsl} = 2^{-\lev/2}\{g(0)-g(1)\}\int_{y>a_{\tsl}}\psi(y)\,dy +
o(2^{-\lev/2})$; since $\int\psi = 0$ and $\psi$ vanishes on no subinterval of
its support, that partial integral is nonzero for some translate, and the
boundary contribution of the level is of order $2^{-\lev}$. The failure of
uniform convergence is that of any wavelet series at a jump.
\end{proof}

Under Assumption 5 nothing changes, since that assumption is already a
periodic one; what Proposition S3.1 gives is the price of its failure. With
the basis of \citet*{cohen1993wavelets} the space $\VJ$ reproduces polynomials
of degree below $N$ and its basis characterizes $\Besov{s}{\pi}{r}[0,1]$
without periodicity \citep[Section 3.9]{Cohen-2003}, so Lemma 2 holds with
that space and the same proof.

%%%%%%%%%%%%%%%%%%%%%%%%%%%%%%%%%%%%%%%%%%%%%%%%%%%%%%%%%%%%%%%%%%%%%%%%%%%%%%
% derivations/03-desenho-produtos.tex
\section{The Gram matrix of the product design}
\setcounter{equation}{0}

Recall $\bPsi(u) \in \R^{K}$, $K = 1+q\NJ$, the vector that collects the
constant and the wavelets of the $q$ modulators, and the identity
$\bZ = \Pi\{\bX \otimes \bPsi(\bU)\}$, where $\Pi$ is the fixed permutation
that moves the $p$ entries $X_{\cv}\cdot 1$ to the front. Since $\Pi$ is
orthogonal, every statement about eigenvalues is invariant to the column
order. Assumption 4 also gives a constant $C_\psi$, depending only on the
family, with
\begin{equation}
  \label{eq:pointwise}
  \sup_{u\in[0,1]}\sum_{\lev<J}\sum_{\tsl}\wav{\lev}{\tsl}(u)^2
  \;\le\; C_\psi\,2^{J} \qquad \text{for every } J \ge 1 :
\end{equation}
at each level at most $L$ periodized wavelets are nonzero at a given $u$ and
$\norm{\wav{\lev}{\tsl}}_{\infty} = 2^{\lev/2}\norm{\psi}_{\infty}$, so the
left side is at most
$L\norm{\psi}^2_{\infty}\sum_{\lev<J}2^{\lev} \le L\norm{\psi}^2_{\infty}2^J$.

\subsection*{Proof of Proposition 2}

\emph{Part (i), the modulators.} For $a \in \R^{K}$ write
$h_a(u) = a\T\bPsi(u) = a_0 + \sum_{\md}h_{a,\md}(u_{\md})$ with
$h_{a,\md} = \sum_{\lev\tsl}a_{\md,\lev\tsl}\wav{\lev}{\tsl}$. The functions
$1$ and $\wav{\lev}{\tsl}(u_{\md})$, $\md = 1,\dots,q$, are orthonormal in
$L_2([0,1]^q,du)$: within one coordinate this is the orthonormality of the
periodized system, and across coordinates
$\int\wav{\lev}{\tsl}(u_{\md})\wav{\lev'}{\tsl'}(u_{\md'})\,du =
\int_0^1\wav{\lev}{\tsl}\int_0^1\wav{\lev'}{\tsl'} = 0$ for
$\md \ne \md'$, because every wavelet integrates to zero. Hence
$\int_{[0,1]^q}h_a^2 = \norm{a}_2^2$ and, by Assumption 3,
$c_U\norm{a}_2^2 \le \E\{h_a(\bU)^2\} \le C_U\norm{a}_2^2$, which is the claim
since $a\T\Gram_{\bPsi}a = \E\{h_a(\bU)^2\}$.

\emph{Part (i), the full matrix.} Let $\delta \in \R^{pK}$ and reshape it as
$\Delta \in \R^{p\times K}$, row $\cv$ holding the coefficients attached to
$X_{\cv}$, so that $\{\bX\otimes\bPsi(\bU)\}\T\delta = \bX\T\Delta\bPsi(\bU)$.
By the tower property,
\[
  \delta\T\Pi\T\Gram\Pi\delta
  = \E\bigl[\{\bX\T\Delta\bPsi(\bU)\}^2\bigr]
  = \E\bigl[\bPsi(\bU)\T\Delta\T\,\E(\bX\bX\T\mid\bU)\,\Delta\bPsi(\bU)\bigr].
\]
Assumption 2 bounds the inner quadratic form between
$\kappa_1\norm{\Delta\bPsi(\bU)}_2^2$ and $\kappa_2\norm{\Delta\bPsi(\bU)}_2^2$
almost surely, and
$\E\norm{\Delta\bPsi(\bU)}_2^2 = \delta\T(I_p\otimes\Gram_{\bPsi})\delta$ lies
between $\lmin(\Gram_{\bPsi})\norm{\delta}_2^2$ and
$\lmaxx(\Gram_{\bPsi})\norm{\delta}_2^2$. Combining the two and using the
first part gives $\gamma = \kappa_1c_U \le \lmin(\Gram)$ and
$\lmaxx(\Gram) \le \kappa_2C_U$. If $\bX$ and $\bU$ are independent, the
conditional second moment matrix is constant, the same computation gives
$\delta\T\Pi\T\Gram\Pi\delta = \delta\T\{\E(\bX\bX\T)\otimes\Gram_{\bPsi}\}
\delta$ for every $\delta$, and two symmetric matrices with the same quadratic
form are equal; the eigenvalues of a Kronecker product of symmetric matrices
are the products of the eigenvalues, and
$\lmin\{\E(\bX\bX\T)\} \ge \kappa_1$ by Jensen's inequality, since $\lmin$ is
concave and $\E(\bX\bX\T) = \E\{\E(\bX\bX\T\mid\bU)\}$.

\emph{Part (ii), the empirical matrix.} First the range:
$\norm{\bZ_i}_2^2 = \norm{\bX_i}_2^2\norm{\bPsi(\bU_i)}_2^2 \le
pB_X^2(1+qC_\psi2^J) = R_J$ by \eqref{eq:pointwise} applied to each of the $q$
coordinates. Set $A_i = (\bZ_i\bZ_i\T-\Gram)/n$, independent, centered and
symmetric, with $\Gramhat - \Gram = \sum_i A_i$. Then
$\opnorm{\bZ_i\bZ_i\T} = \norm{\bZ_i}_2^2 \le R_J$ and
$\opnorm{\Gram} \le R_J$, so $\opnorm{A_i} \le 2R_J/n$; and
$\E\{(\bZ\bZ\T-\Gram)^2\} = \E(\norm{\bZ}_2^2\bZ\bZ\T) - \Gram^2 \preceq
R_J\Gram$, so that
$\opnorm{\sum_i\E A_i^2} \le R_J\lmaxx(\Gram)/n \le R_J\kappa_2C_U/n$. The
matrix Bernstein inequality \citep[Theorem 1.4]{tropp2012user} applied to
$\pm\sum_iA_i$ gives, for every $t>0$,
\[
  \Prob\bigl(\opnorm{\Gramhat-\Gram} \ge t\bigr)
  \;\le\; 2D\exp\Bigl\{-\frac{nt^2}{2R_J(\kappa_2C_U+2t/3)}\Bigr\},
\]
and $t = \gamma/2$ with Weyl's inequality gives
$\lmin(\Gramhat) \ge \lmin(\Gram)-\gamma/2 \ge \gamma/2$ on the complementary
event, which is the display of part (ii). The exponent is $-cn/R_J$ with $c$
fixed and $R_J \asymp pq2^J$, while $\log(2D) \asymp \log(pq2^J)$, so the
bound tends to zero exactly when $n/(pq2^J) - \log(pq2^J) \to \infty$, which
is implied by $pq2^J\log(pq2^J)/n \to 0$. On that event
$\compat^2(S;\Gramhat) \ge \gamma/2$ for every support, by
\eqref{eq:compat-lmin}. \qed

\begin{remark}[What the two assumptions exclude]
\label{rem:excl}
Assumption 2 fails exactly when a nontrivial linear combination of the linear
covariates is almost surely a function of the modulators on a set of positive
probability; in the canonical case $X_1\equiv1$ and $X_2 = h(U_1)$ the
conditional second moment matrix has rank one and the population Gram matrix
degenerates as $J$ grows, because $h(u_1)\wav{\lev}{\tsl}(u_1)$ is approached
by the span of the basis. In Assumption 3 it is the joint density that is
bounded below: if $U_2 = U_1$ almost surely, both marginals are uniform but
$\wav{\lev}{\tsl}(U_1)-\wav{\lev}{\tsl}(U_2) \equiv 0$ and
$\Gram_{\bPsi}$ is singular. Independence between $\bX$ and $\bU$ buys the
exact factorization of part (i) and nothing else.
\end{remark}

%%%%%%%%%%%%%%%%%%%%%%%%%%%%%%%%%%%%%%%%%%%%%%%%%%%%%%%%%%%%%%%%%%%%%%%%%%%%%%
% derivations/04-oraculo.tex
\section{The oracle inequality}
\setcounter{equation}{0}

{\color{gray}\noindent[k = 4: the calibration lemma of the lasso, the proof of
Theorem~2 for the lasso, the remark on the cone and the corollary that
instantiates the design condition, which stood in this section, are moved to
Section~S8, unchanged except for the numbering of the results they cite;
Lemmas~S5.1, S5.3 and S5.4 stay, the first extended to the penalty of
equation~(2.5).]\par}

\textcolor{gray}{\sout{Four lemmas precede the proof of Theorem 2.}}
\textcolor{colR1}{A proposition on the partition into chunks and four lemmas
precede the proof of Theorem~2.}

% derivations/08-blocos.tex, Proposicao 7 (numeracao global).
{\color{colR1}
\begin{proposition}[The partition into chunks]
\label{prop:balanced}
Let $\Gcal$ be the partition of Section~2.3 of the paper with $b \ge 2$ and
$J \ge 1$, and let $\jst = \max\{\lev \ge 0 : 2^{\lev} < b\}$.
\begin{enumerate}
\item[(i)] Every chunk of a fine level, $\lev > \jst$, is contained in that
level and has between $b$ and $2b-1$ coefficients. If $J > \jst+1$, the coarse
chunk has $2^{\jst+1}-1$ coefficients, and $b-1 \le 2^{\jst+1}-1 \le 2b-3$;
if $J \le \jst+1$, each block is a single chunk of $2^{J}-1$ coefficients.
\item[(ii)] $|\Gcal| \le pq(1 + 2^{J}/b)$.
\item[(iii)] With $w_{G} = |G|^{1/2}$,
$\varrho^{2} = \max_{G}|G|/\min_{G}|G| \le (2b-1)/(b-1) \le 3$.
\end{enumerate}
\end{proposition}

\begin{proof}
All blocks have the same structure, so it suffices to look at one. The coarse
chunk has $\sum_{\lev=0}^{\min(\jst,J-1)}2^{\lev} = 2^{\min(\jst,J-1)+1}-1$
coefficients. If $J \le \jst+1$, this is $2^{J}-1$ and there is no fine level.
Otherwise $2^{\jst} < b \le 2^{\jst+1}$ and, $b$ being an integer,
$2^{\jst} \le b-1$, so that $b-1 \le 2^{\jst+1}-1 \le 2b-3$. A fine level
$\lev$ has $k_{\lev} = \lfloor 2^{\lev}/b\rfloor \ge 1$ chunks, $k_{\lev}-1$ of
them with $b$ coefficients and the last with $b + r_{\lev}$, where
$r_{\lev} = 2^{\lev} - k_{\lev}b \in \{0,\dots,b-1\}$; all lie between $b$ and
$2b-1$, which is (i). For (ii), a block has
$1 + \sum_{\lev=\jst+1}^{J-1}\lfloor 2^{\lev}/b\rfloor \le
1 + (2^{J}-2^{\jst+1})/b \le 1 + 2^{J}/b$ chunks. For (iii), with
$w_{G} = |G|^{1/2}$, $\varrho^{2} = \max|G|/\min|G|$; with no fine level every
chunk has $2^{J}-1$ coefficients and $\varrho = 1$, and otherwise (i) gives
$\min|G| \ge b-1$ and $\max|G| \le 2b-1$, and
$(2b-1)/(b-1) = 2 + 1/(b-1) \le 3$.
\end{proof}
}


\begin{lemma}[Profiling and the Schur complement]
\label{lem:profile}
Suppose $\Amat$ has rank $p$. Then $\thetahat$ minimizes
\textcolor{gray}{\sout{$\normn{\MA\bY-\Btil\btheta}^2+2\pen\norm{\btheta}_1$}}\textcolor{colR1}{$\normn{\MA\bY-\Btil\btheta}^2+2\pen\nGw{\btheta}$,
or the same with $\norm{\btheta}_1$ for the lasso,} and
$\chat = (\Amat\T\Amat)^{-1}\Amat\T(\bY-\Bmat\thetahat)$; with
$v = \thetahat-\thetastar$,
\begin{equation}
  \label{eq:decomp}
  \widehat{f}-f_J \;=\; \Btil v + \PA(\bb+\bveps)\textcolor{colR1}{,}
\end{equation}
\textcolor{colR1}{and the same identity holds with any comparator
$\bar\bbeta$, $\bbar$ and $\vbar = \thetahat-\thbar$ in place of $\betastar$,
$\bb$ and $v$;}
$\GBA$ is the Schur complement of the block $\Amat$ in $\Gramhat$ and
\begin{equation}
  \label{eq:schur}
  \gtil \;=\; \lmin(\GBA) \;\ge\; \lmin(\Gramhat);
\end{equation}
and if $\lmin(\Gramhat)>0$ the minimizer in (2.5) is unique.
\end{lemma}

\begin{proof}
\textcolor{colR1}{The argument uses only that the penalty does not depend on
$\bc$ and is convex, so it holds for both penalties.} The objective is convex and, for fixed $\btheta$, quadratic and coercive in
$\bc$, with minimum at the stated value and minimized value
$\normn{\MA(\bY-\Bmat\btheta)}^2 = \normn{\MA\bY-\Btil\btheta}^2$, since $\MA$
is a symmetric idempotent; minimizing first in $\bc$ and then in $\btheta$
gives the same set of minimizers as minimizing jointly. For
\eqref{eq:decomp}, $\Amat\chat = \PA(\bY-\Bmat\thetahat)$ and
$\Amat\bc = \PA\Amat\bc$, and $\bY = \Amat\bc+\Bmat\thetastar+\bb+\bveps$, so
$\Amat(\chat-\bc) = \PA(\bb+\bveps)-\PA\Bmat v$ and
$\widehat{f}-f_J = \MA\Bmat v + \PA(\bb+\bveps)$\textcolor{colR1}{; the same
computation with $\bY = \Amat\bar\bc + \Bmat\thbar + \bbar + \bveps$ gives the
identity for a comparator}. For \eqref{eq:schur}, the
variational characterization of the Schur complement gives, for every
$v \in \R^d$,
\[
  v\T\GBA v \;=\; \min_{u\in\R^p}(u\T,v\T)\,\Gramhat\,(u\T,v\T)\T
  \;\ge\; \min_{u}\lmin(\Gramhat)(\norm{u}_2^2+\norm{v}_2^2)
  \;=\; \lmin(\Gramhat)\norm{v}_2^2 .
\]
Finally $\lmin(\Gramhat)>0$ makes the quadratic term strictly convex, and
adding a convex function preserves strict convexity.
\end{proof}

% derivations/08-blocos.tex, Lema 14 (numeracao global).
{\color{colR1}
\begin{lemma}[Calibration of the penalty over chunks]
\label{lem:calib-blocks}
Under Assumption~1, let $\Gcal$ be any partition of the penalized coordinates
and $\alpha \in (0,1)$.
\begin{enumerate}
\item[(i)] With $\lamz{G}$ as in Theorem~2, the event
$\Ecal = \{\norm{(\Btil\T\bveps/n)_{G}}_{2} \le \lamz{G} \text{ for every } G\}$
has $\Prob(\Ecal\mid\bX,\bU) \ge 1-\alpha$. Consequently, for any weights $w$
and every $\pen \ge \pen_{w}$, $\Ecal \subseteq \Tw$ and
$\Prob(\Tw) \ge 1-\alpha$.
\item[(ii)] (Closed form.) $\mathrm{tr}(\Psitil_{G}) \le |G|\smax^{2}$ and
$\opnorm{\Psitil_{G}} \le \lmaxx(\Gramhat)$; hence, on the event
$\{\smax^{2} \le 2B_X^{2}C_U\}\cap\{\lmaxx(\Gramhat) \le \Lambda\}$,
$\lamz{G} \le \lamzbar{G}$ for every $G$, with
\[
  \lamzbar{G} \;=\; \frac{\sigma}{\sqrt{n}}\Bigl\{B_X\sqrt{2C_U|G|}
  + \sqrt{2\Lambda\log(|\Gcal|/\alpha)}\Bigr\},
\]
which is deterministic and is the quantity of equation (3.1) of the paper.
\item[(iii)] (The price of the weights.) With
$\lamG = 2\max_{G}\lamzbar{G}/w_{G}$ and $\lamGone = 2\max_{G}\lamzbar{G}$,
$\lamG w_{G} \le \varrho\lamGone$ for every $G$, and therefore
$(\lamG)^{2}\Wcal(\Hcal) \le \varrho^{2}(\lamGone)^{2}|\Hcal|$ for every
$\Hcal \subseteq \Gcal$. The same holds with $\pen_{w}$ and
$2\max_{G}\lamz{G}$ in place of $\lamG$ and $\lamGone$.
\end{enumerate}
\end{lemma}

The event $\Ecal$ does not depend on the weights: they only choose the chunk
that dictates the penalty, the one with the largest $\lamz{G}/w_{G}$, and the
price of unequal weights is the factor $\varrho$ of (iii), paid in the
estimation term. With unit weights and chunks of at most $b_{\max}$
coefficients, $\lamGone \asymp \sigma\{(b_{\max} + \log|\Gcal|)/n\}^{1/2}$,
which is the penalty level of \citet[eq.~3.14]{Klopp-Pensky-2015} up to
constants and the sub-Gaussian version of the level of
\citet[Lemma~3.1]{Lounici-Pontil-vandeGeer-Tsybakov-2011}.

\begin{proof}
\emph{(i).} Condition on $(\bX,\bU)$. Then $\Btil$ is fixed and, by
Assumption~1, the errors are independent with
$\E\{\exp(a_{i}\varepsilon_{i})\mid\bX,\bU\} \le \exp(\sigma^{2}a_{i}^{2}/2)$,
so that $\E\{\exp(a\T\bveps)\mid\bX,\bU\} \le \exp(\sigma^{2}\norm{a}_{2}^{2}/2)$
for every $a \in \R^{n}$, which is the condition of
\citet[Theorem~2.1]{Hsu-Kakade-Zhang-2012} with mean zero. Fix $G$ and write
$M_{G} = \Btil_{G}\T/n$, so that $(\Btil\T\bveps/n)_{G} = M_{G}\bveps$, and
$\Sigma_{G} = M_{G}\T M_{G} \in \R^{n\times n}$. The theorem is stated for a
square matrix $A$ with $\Sigma = A\T A$; take $A = \Sigma_{G}^{1/2}$, which has
the same $\Sigma_{G}$, and note that
$\norm{\Sigma_{G}^{1/2}\bveps}_{2}^{2} = \bveps\T\Sigma_{G}\bveps =
\norm{M_{G}\bveps}_{2}^{2}$. Since $\mathrm{tr}(\Sigma_{G}) =
\mathrm{tr}(\Psitil_{G})/n$, $\opnorm{\Sigma_{G}} = \opnorm{\Psitil_{G}}/n$
and $\mathrm{tr}(\Sigma_{G}^{2}) \le \opnorm{\Sigma_{G}}\mathrm{tr}(\Sigma_{G})$,
the theorem gives, for every $t > 0$, with conditional probability at least
$1 - e^{-t}$,
\[
  \norm{(\Btil\T\bveps/n)_{G}}_{2}^{2}
  \;\le\; \sigma^{2}\bigl[\mathrm{tr}(\Sigma_{G})
  + 2\{\mathrm{tr}(\Sigma_{G}^{2})t\}^{1/2} + 2\opnorm{\Sigma_{G}}t\bigr]
  \;\le\; \frac{\sigma^{2}}{n}\Bigl[\{\mathrm{tr}(\Psitil_{G})\}^{1/2}
  + \{2\opnorm{\Psitil_{G}}t\}^{1/2}\Bigr]^{2}.
\]
With $t = \log(|\Gcal|/\alpha)$ each chunk fails with conditional probability
at most $\alpha/|\Gcal|$, and a union bound over the chunks and the tower
property give $\Prob(\Ecal) \ge 1-\alpha$. On $\Ecal$,
$\norm{(\Btil\T\bveps/n)_{G}}_{2}/w_{G} \le \lamz{G}/w_{G} \le \pen_{w}/2 \le
\pen/2$ for every $G$, so $\Ecal \subseteq \Tw$.

\emph{(ii).} Since $\MA \preceq I_{n}$,
$\Psitil_{G} = \Bmat_{G}\T\MA\Bmat_{G}/n \preceq \Gramhat_{GG}$, the principal
submatrix of $\Gramhat$ on $G$. Hence $\mathrm{tr}(\Psitil_{G}) \le
\sum_{a\in G}(\Gramhat)_{aa} \le |G|\smax^{2}$, and
$\opnorm{\Psitil_{G}} \le \lmaxx(\Gramhat_{GG}) \le \lmaxx(\Gramhat)$ by Cauchy
interlacing. On the event of the statement,
$\smax \le B_X(2C_U)^{1/2}$ and $\lmaxx(\Gramhat) \le \Lambda$.

\emph{(iii).} For every $G$,
$\lamG w_{G} = 2\max_{G'}(\lamzbar{G'}w_{G}/w_{G'}) \le
2\max_{G'}\lamzbar{G'}\cdot\max_{G'}(w_{G}/w_{G'}) \le \varrho\lamGone$, and
$(\lamG)^{2}\Wcal(\Hcal) = \sum_{G\in\Hcal}(\lamG w_{G})^{2} \le
\varrho^{2}(\lamGone)^{2}|\Hcal|$. The same computation holds with the
$\lamz{G}$.
\end{proof}
}

\begin{lemma}[Empirical column norms]
\label{lem:columns}
Under Assumptions 2 to 4, every penalized column has
$\E\{(\Gramhat)_{aa}\} \le B_X^2C_U$ and, for every $t>0$, with
$R'_J = B_X^2\norm{\psi}^2_{\infty}2^{J-1}$,
\[
  \Prob(\smax^2 > B_X^2C_U+t) \;\le\;
  d\exp\Bigl\{-\frac{nt^2}{2(R'_JB_X^2C_U+R'_Jt/3)}\Bigr\} .
\]
In particular $\Prob(\smax^2 \le 2B_X^2C_U) \to 1$ whenever
$2^J\log d/n \to 0$.
\end{lemma}

\begin{proof}
For the column $a = (\cv\md,\lev\tsl)$ the summand is
$W_i = X_{i\cv}^2\wav{\lev}{\tsl}(U_{i\md})^2 \ge 0$, with
$\E(W) \le B_X^2\int_0^1\wav{\lev}{\tsl}^2p_{U_{\md}} \le B_X^2C_U$ by
orthonormality and the bounded marginal density, and
$0 \le W_i \le B_X^2\norm{\psi}_{\infty}^22^{J-1} = R'_J$ because
$\lev \le J-1$. Then $\Var(W) \le \E(W^2) \le R'_J\E(W) \le R'_JB_X^2C_U$, and
the one sided Bernstein inequality with a union bound over the $d$ columns
gives the display. With $t = B_X^2C_U$ the exponent is
$-3nB_X^2C_U/(8R'_J)$ and $R'_J \asymp 2^J$.
\end{proof}

\begin{lemma}[The $\ell_1$ norm of the \textcolor{colR1}{oracle}]
\label{lem:l1}
Under Assumption 5, for every pair $(\cv,\md)$ and every level,
$\sum_{\tsl}|\coefs{\cv}{\md}{\lev}{\tsl}| \le
2^{\lev(1-1/\pi)}\norm{\theta^{*}_{\cv\md,\lev\cdot}}_{\pi} \le
C_g2^{\lev(1/2-s)}$, with no dependence on $\pi$, and therefore
$\norm{\thetastar}_1 \le pqC_g\sum_{\lev<J}2^{\lev(1/2-s)}$, which is
$O(1)$ if $s>1/2$, of order $J$ if $s = 1/2$ and of order
$2^{J(1/2-s)}$ if $s<1/2$.
\end{lemma}

\begin{proof}
Hölder's inequality in $\R^{2^{\lev}}$ gives
$\norm{a}_1 \le (2^{\lev})^{1-1/\pi}\norm{a}_{\pi}$, and Assumption 5 bounds
the $\ell_\pi$ norm of the level by $C_g2^{-\lev(s+1/2-1/\pi)}$; the resulting
exponent is $\lev(1-1/\pi)-\lev(s+1/2-1/\pi) = \lev(1/2-s)$, in which $\pi$
has disappeared. Summing over the levels below $J$ and over the $pq$ blocks
gives the three regimes of the geometric sum.
\end{proof}

% derivations/08-blocos.tex, prova do Teorema 3 (numeracao global).
{\color{colR1}
\subsection*{Proof of Theorem 2}

The event $\Tw$ has probability at least $1-\alpha$ whenever
$\pen \ge \pen_{w}$, by Lemma~\ref{lem:calib-blocks}(i). Fix a comparator
$\bar\bbeta = (\bar\bc,\thbar)$ and write $\norm{\cdot}$ for $\nGw{\cdot}$ in
this proof.

\emph{Step 1, the basic inequality.} By Lemma~\ref{lem:profile}, $\thetahat$
minimizes the residualized objective; comparing its value at $\thetahat$ with
its value at $\thbar$, and using $\bY = \Amat\bar\bc + \Bmat\thbar + \bbar +
\bveps$ and $\MA\Amat = 0$, so that $\MA\bY = \Btil\thbar + r$ with
$r = \MA(\bbar + \bveps)$,
\begin{equation}
  \label{eq:basic-blocks}
  \normn{\Btil\vbar}^{2} + 2\pen\norm{\thetahat}
  \;\le\; 2\inner{r}{\Btil\vbar}_{n} + 2\pen\norm{\thbar} .
\end{equation}

\emph{Step 2, the two noises.} Since $\MA$ is symmetric and $\MA\Btil = \Btil$,
$\inner{r}{\Btil\vbar}_{n} = \inner{\bbar+\bveps}{\Btil\vbar}_{n}$. The
stochastic part is handled by the dual norm of $\norm{\cdot}$, which is
$z \mapsto \max_{G}\norm{z_{G}}_{2}/w_{G}$: on $\Tw$,
\[
  2\bigl|\inner{\bveps}{\Btil\vbar}_{n}\bigr|
  = 2\Bigl|\sum_{G}\inner{(\Btil\T\bveps/n)_{G}}{\vbar_{G}}\Bigr|
  \le 2\max_{G}\frac{\norm{(\Btil\T\bveps/n)_{G}}_{2}}{w_{G}}
  \sum_{G}w_{G}\norm{\vbar_{G}}_{2}
  \le \pen\norm{\vbar} .
\]
The deterministic part is handled by the Cauchy-Schwarz and Young
inequalities: $2|\inner{\bbar}{\Btil\vbar}_{n}| =
2|\inner{\MA\bbar}{\Btil\vbar}_{n}| \le 2\normn{\bbar}\normn{\Btil\vbar} \le
\normn{\Btil\vbar}^{2}/2 + 2\normn{\bbar}^{2}$. Substituting both in
\eqref{eq:basic-blocks},
\begin{equation}
  \label{eq:basic2-blocks}
  \tfrac12\normn{\Btil\vbar}^{2} + 2\pen\norm{\thetahat}
  \;\le\; \pen\norm{\vbar} + 2\pen\norm{\thbar} + 2\normn{\bbar}^{2} .
\end{equation}

\emph{Step 3, the slow rate.} Bounding $\norm{\vbar} \le \norm{\thetahat} +
\norm{\thbar}$ in \eqref{eq:basic2-blocks} gives part (i). No property of
$\Gramhat$ was used.

\emph{Step 4, decomposability.} For $G \notin \Gbar$, $\thbar_{G} = 0$ and
$\thetahat_{G} = \vbar_{G}$; for $G \in \Gbar$,
$\norm{\thetahat_{G}}_{2} \ge \norm{\thbar_{G}}_{2} - \norm{\vbar_{G}}_{2}$.
Hence $\norm{\thetahat} \ge \norm{\thbar} - \norm{\vbar_{\Gbar}} +
\norm{\vbar_{\Gbar^{c}}}$, while $\norm{\vbar} = \norm{\vbar_{\Gbar}} +
\norm{\vbar_{\Gbar^{c}}}$; substituting in \eqref{eq:basic2-blocks} and
cancelling $2\pen\norm{\thbar}$,
\begin{equation}
  \label{eq:cone-blocks}
  \tfrac12\normn{\Btil\vbar}^{2} + \pen\norm{\vbar_{\Gbar^{c}}}
  \;\le\; 3\pen\norm{\vbar_{\Gbar}} + 2\normn{\bbar}^{2} .
\end{equation}

\emph{Step 5, two cases, with no cone.} By the Cauchy-Schwarz inequality and
the definition of $\gtil$,
\begin{equation}
  \label{eq:cs-blocks}
  \norm{\vbar_{\Gbar}} = \sum_{G\in\Gbar}w_{G}\norm{\vbar_{G}}_{2}
  \le \bar\Wcal^{1/2}\norm{\vbar}_{2}
  \le (\bar\Wcal/\gtil)^{1/2}\normn{\Btil\vbar},
\end{equation}
because $\normn{\Btil\vbar}^{2} = \vbar\T\GBA\vbar \ge \gtil\norm{\vbar}_{2}^{2}$.
If $2\normn{\bbar}^{2} \le \pen\norm{\vbar_{\Gbar}}$, then
\eqref{eq:cone-blocks} gives $\tfrac12\normn{\Btil\vbar}^{2} +
\pen\norm{\vbar_{\Gbar^{c}}} \le 4\pen\norm{\vbar_{\Gbar}}$, and with
\eqref{eq:cs-blocks}, $\tfrac12\normn{\Btil\vbar}^{2} \le
4\pen(\bar\Wcal/\gtil)^{1/2}\normn{\Btil\vbar}$, that is
$\normn{\Btil\vbar}^{2} \le 64\pen^{2}\bar\Wcal/\gtil$ and
$\norm{\vbar}_{2}^{2} \le \normn{\Btil\vbar}^{2}/\gtil \le
64\pen^{2}\bar\Wcal/\gtil^{2}$. If instead
$2\normn{\bbar}^{2} > \pen\norm{\vbar_{\Gbar}}$, then \eqref{eq:cone-blocks}
gives $\tfrac12\normn{\Btil\vbar}^{2} < 8\normn{\bbar}^{2}$, that is
$\normn{\Btil\vbar}^{2} < 16\normn{\bbar}^{2}$, and
$\norm{\vbar}_{2}^{2} < 16\normn{\bbar}^{2}/\gtil$. Each quantity is bounded by
the maximum of the two cases, hence by their sum, which is part (ii).

\emph{Step 6, prediction.} By Lemma~\ref{lem:profile} with the comparator,
$\widehat{f}-f = \Btil\vbar + \PA(\bbar+\bveps) - \bbar$, and
$\normn{\PA\bbar} \le \normn{\bbar}$, which gives the triangle inequality of
part (iii). Conditionally on the covariates the errors are independent,
centered and of variance at most $\sigma^{2}$, so
$\E(\normn{\PA\bveps}^{2}\mid\bX,\bU) \le \sigma^{2}\mathrm{tr}(\PA)/n \le
\sigma^{2}p/n$. Since $\Tw$, $\gtil$ and $\pen$ do not depend on the
comparator, the bounds hold simultaneously for all comparators, and with
$(x+y+z)^{2} \le 3(x^{2}+y^{2}+z^{2})$ and part (ii) they give, on
$\Tw\cap\{\gtil > 0\}$,
\begin{equation}
  \label{eq:oracle-min-blocks}
  \normn{\widehat{f}-f}^{2} \;\le\;
  3\min_{\bar\bbeta}\Bigl\{\frac{64\pen^{2}\bar\Wcal}{\gtil}
  + 20\normn{\bbar}^{2}\Bigr\} + 3\normn{\PA\bveps}^{2} . \qquad\qed
\end{equation}

\begin{remark}[Why there is no cone condition, also over chunks]
\label{rem:nocone-blocks}
The step that dispenses with the cone is the one of the lasso
(Remark~\ref{rem:nocone}): profiling removes the levels from the problem, and
the full smallest eigenvalue of Proposition 2 controls
$\norm{\vbar_{\Gbar}} \le (\bar\Wcal/\gtil)^{1/2}\normn{\Btil\vbar}$ over the
whole space, by \eqref{eq:cs-blocks}.
\citet{Lounici-Pontil-vandeGeer-Tsybakov-2011} need a restricted eigenvalue
condition for groups on a weighted cone (their Assumption~3.1); here it is
dominated by the full smallest eigenvalue, and the constant is the same for
every $\Gbar$.
\end{remark}

% derivations/08-blocos.tex, Observacao 2 (a variante branca); D47(d).
\begin{remark}[The estimator that the software computes]
\label{rem:white}
The software standardizes the columns, orthonormalizes each chunk and
penalizes the orthonormalized coordinates with the weights $|G|^{1/2}$. On the
original scale its penalty is $\sum_{G}w_{G}\norm{Q_{G}^{1/2}\btheta_{G}}_{2}$,
with $Q_{G} = \Bmat_{G}\T M_{1}\Bmat_{G}/n$ the centered Gram matrix of the
chunk, $M_{1} = I_{n} - n^{-1}\mathbf{1}\mathbf{1}\T$.
% L11 (D55, pergunta 49(c)): Simon & Tibshirani (2012), eq. (1.4) e Section 2.
Since
$\norm{M_{1}\Bmat_{G}\btheta_{G}}_{2} = n^{1/2}\norm{Q_{G}^{1/2}\btheta_{G}}_{2}$,
this is the standardized group lasso of
\citet[eq.~1.4]{Simon-Tibshirani-2012}, which penalizes the fit of each
chunk rather than its coefficients and which they show to be equivalent to
orthonormalizing within groups (their Section~2). Theorem~2 holds with any
$Q_{G}$ positive definite and measurable with respect to the covariates, with
three changes: the event becomes
$\{\max_{G}\norm{Q_{G}^{-1/2}(\Btil\T\bveps/n)_{G}}_{2}/w_{G} \le \pen/2\}$, by
the dual norm; $\bar\Wcal$ becomes $q_{\max}\bar\Wcal$ in parts (ii) and
(iii), with $q_{\max} = \max_{G}\lmaxx(Q_{G})$, because Step 5 reads
$\sum_{G\in\Gbar}w_{G}\norm{Q_{G}^{1/2}\vbar_{G}}_{2} \le
(q_{\max}\bar\Wcal)^{1/2}\norm{\vbar}_{2}$; and Lemma~\ref{lem:calib-blocks}(i)
holds with $Q_{G}^{-1/2}\Psitil_{G}Q_{G}^{-1/2}$ in place of $\Psitil_{G}$, by
the same proof with $Q_{G}^{-1/2}\Btil_{G}\T/n$ in place of $\Btil_{G}\T/n$.
When the constant covariate $X_{1}\equiv 1$ is among the columns of $\Amat$,
$\PA \succeq I_{n} - M_{1}$, so that $\MA \preceq M_{1}$ and
$\Psitil_{G} \preceq Q_{G}$: the trace of $Q_{G}^{-1/2}\Psitil_{G}Q_{G}^{-1/2}$
is at most $|G|$ and its norm at most one, and
\[
  \lamz{G} \;\le\; \frac{\sigma}{\sqrt{n}}\Bigl\{|G|^{1/2}
  + \sqrt{2\log(|\Gcal|/\alpha)}\Bigr\},
\]
which is pivotal, with neither $\smax$ nor $\Lambda$; and
$q_{\max} \le \lmaxx(\Gramhat) \le \Lambda$ on $\ESig$, since
$Q_{G} \preceq \Gramhat_{GG}$. The rates of Section~S6 hold for this estimator
with the constant multiplied by $\Lambda$.
\end{remark}

% derivations/08-blocos.tex, Corolario 9 (numeracao global).
\begin{corollary}[The ideal risk form]
\label{cor:ideal}
In the conditions of Theorem~2, on the event $\Tw\cap\{\gtil > 0\}$, let
$\Lambda' \ge \lmaxx(\Gramhat)$, for instance $\Lambda' = \Lambda$ on $\ESig$,
and
\[
  \eta \;=\; \frac{1.6\max_{G}(\pen w_{G})^{2}}{\Lambda'\gtil}.
\]
Then
\[
  \normn{\widehat{f}-f}^{2} \;\le\; 120\Lambda'\Rid(\thetastar;\eta)
  + 120\normn{\bb}^{2} + 3\normn{\PA\bveps}^{2},
  \qquad
  \norm{\thetahat-\thetastar}_{2}^{2} \;\le\;
  \frac{82\Lambda'\Rid(\thetastar;\eta) + 64\normn{\bb}^{2}}{\gtil}.
\]
With $\pen = \lamG$, or $\pen = \pen_{w}$, Lemma~\ref{lem:calib-blocks}(iii)
gives $\eta \le \varrho^{2}\eta_{1}$, with
$\eta_{1} = 1.6(\lamGone)^{2}/(\Lambda'\gtil)$ the level of unit weights, and
$\Rid(\thetastar;\eta) \le \varrho^{2}\Rid(\thetastar;\eta_{1})$.
\end{corollary}

\begin{proof}
Let $T = \{G : \norm{\thetastar_{G}}_{2}^{2} > \eta\}$ and take the comparator
$\bar\bbeta = (\bc,\thbar)$ with $\thbar_{G} = \thetastar_{G}\mathbf{1}\{G\in T\}$.
Then $\Gbar \subseteq T$ and, by the definition of $\eta$,
$\pen^{2}\bar\Wcal \le \sum_{G\in T}(\pen w_{G})^{2} \le
|T|\max_{G}(\pen w_{G})^{2} = |T|\eta\Lambda'\gtil/1.6$. The residual is
$\bbar = \bb + \Bmat(\thetastar-\thbar)$, and since $\Bmat\T\Bmat/n$ is a
principal submatrix of $\Gramhat$,
\[
  \normn{\bbar}^{2} \;\le\; 2\normn{\bb}^{2}
  + 2\lmaxx(\Gramhat)\norm{\thetastar-\thbar}_{2}^{2}
  \;\le\; 2\normn{\bb}^{2} + 2\Lambda'\sum_{G\notin T}\norm{\thetastar_{G}}_{2}^{2}.
\]
By \eqref{eq:oracle-min-blocks},
$\normn{\widehat{f}-f}^{2} \le 192\pen^{2}\bar\Wcal/\gtil + 60\normn{\bbar}^{2}
+ 3\normn{\PA\bveps}^{2} \le 120\Lambda'\{\eta|T| +
\sum_{G\notin T}\norm{\thetastar_{G}}_{2}^{2}\} + 120\normn{\bb}^{2} +
3\normn{\PA\bveps}^{2}$, and the term in braces is exactly
$\Rid(\thetastar;\eta)$. For the coefficients, Theorem~2(ii) gives
$\norm{\vbar}_{2}^{2} \le 64\pen^{2}\bar\Wcal/\gtil^{2} +
16\normn{\bbar}^{2}/\gtil \le (40\Lambda'\eta|T| + 32\normn{\bb}^{2} +
32\Lambda'\sum_{G\notin T}\norm{\thetastar_{G}}_{2}^{2})/\gtil \le
(40\Lambda'\Rid + 32\normn{\bb}^{2})/\gtil$, and
$\norm{\thetahat-\thetastar}_{2}^{2} \le 2\norm{\vbar}_{2}^{2} +
2\norm{\thbar-\thetastar}_{2}^{2} \le (80\Lambda'\Rid +
64\normn{\bb}^{2})/\gtil + 2\Rid$; since $\gtil \le \lmin(\Bmat\T\Bmat/n) \le
\lmaxx(\Gramhat) \le \Lambda'$, the Schur complement being dominated by the
block, $2\Rid \le 2\Lambda'\Rid/\gtil$, which gives the second bound. The last
statement is Lemma~\ref{lem:calib-blocks}(iii) and the two properties of
$\Rid$ in Section~S1.
\end{proof}
}

%%%%%%%%%%%%%%%%%%%%%%%%%%%%%%%%%%%%%%%%%%%%%%%%%%%%%%%%%%%%%%%%%%%%%%%%%%%%%%
% derivations/05-taxas.tex
\section{Rates}
\setcounter{equation}{0}

{\color{gray}\noindent[k = 4: the lemma that read Theorem~2 at an arbitrary
comparator for the lasso, and the proofs of Lemma~3, Proposition~3,
Theorem~3, Corollary~1, Theorem~1 and the remark on compressibility for the
lasso, which stood in this section, are moved to Section~S8, unchanged except
for the numbering of the results they cite.]\par}

Throughout this section $p$, $q$, $\sigma$, $B_X$, $c_U$, $C_U$, $\kappa_1$,
$\kappa_2$, $C_g$, $s$, $\pi$, $r$ and $\alpha$ are fixed as $n$ grows,
$J = J_n \to \infty$, $d_n = pq(2^{J_n}-1)$,
\textcolor{gray}{\sout{$L_n = \log(2d_n/\alpha) \asymp J_n$, and $\penn$ is as in (3.1)}}\textcolor{colR1}{$2^{J_n} \le 2n$,
$\bn = \lceil\log n\rceil$ with $n \ge 3$, so that $\bn \ge 2$, $\Gcal_{n}$ is
the partition of Section~2.3 of the paper with $b = \bn$, the weights have
ratio $\varrho \le \bar\varrho$ for a fixed $\bar\varrho$, and $\pen = \lamG$
is as in (3.1)}. \textcolor{gray}{\sout{Since the
constants of Theorem 2 depend on the confidence levels only through
$1/\alpha'$ and $1/\alpha''$, taking $\alpha' = \alpha'' = \epsilon/3$ gives,
for each $\epsilon>0$, a constant $M_\epsilon$ with
$\Prob(\normn{\widehat{f}-f}^2 > M_\epsilon r_n) < \epsilon$, which is the
meaning of the statements in $O_p$ below. We also write
$\Lambda = \kappa_2C_U+\gamma/2$.}} \textcolor{colR1}{The statements in $O_p$
are proved by exhibiting, for each $\epsilon > 0$, an event of probability at
least $1-\epsilon$ for $n$ large, built with the confidence levels
$\alpha = \alpha' = \alpha'' = \epsilon/4$ for the noise, the bias and the
levels, on which the error is at most a constant $M_\epsilon$ times the rate.
The constant $\Lambda = \kappa_2C_U + \gamma/2$ and the event $\ESig$ are those
of Section~S1.}

% derivations/08-blocos.tex, as duas consequencias da Hipotese B.
{\color{colR1}
\begin{lemma}[The penalty level and the regime]
\label{lem:regime}
In the setting above,
\begin{equation}
  \label{eq:order-lambda}
  (\lamGone)^{2} \;\le\; C_{\pen}\,\frac{\sigma^{2}\bn}{n},
  \qquad
  C_{\pen} = 16\bigl\{2B_X^{2}C_U + \Lambda\bigl(1 + \tfrac12\log(2pq/\alpha)\bigr)\bigr\},
\end{equation}
and, on the event $\Acal = \{\smax^{2} \le 2B_X^{2}C_U\}\cap\ESig$,
$\lamG \ge \pen_{w}$, so that $\Ecal \subseteq \Tw$ for the event $\Ecal$ of
Lemma~\ref{lem:calib-blocks}(i), $\gtil > \gamma/2$, $\Amat$ has rank $p$ and
$\lmaxx(\Gramhat) < \Lambda$. Moreover $\Prob(\Acal) \to 1$ whenever
$pq2^{J_n}\log(pq2^{J_n})/n \to 0$.
\end{lemma}

\begin{proof}
By $(x+y)^{2} \le 2x^{2}+2y^{2}$, $(\lamGone)^{2} = 4\max_{G}\lamzbar{G}^{2} \le
(8\sigma^{2}/n)\{2B_X^{2}C_U b_{\max} + 2\Lambda\log(|\Gcal_{n}|/\alpha)\}$, with
$b_{\max} = \max_{G}|G|$. By Proposition~\ref{prop:balanced},
$b_{\max} \le 2\bn - 1 < 2\bn$ and $|\Gcal_{n}| \le pq(1 + 2^{J_n}/\bn) \le
pq(1+n) \le 2pqn$, because $2^{J_n} \le 2n$ and $\bn \ge 2$; hence
$\log(|\Gcal_{n}|/\alpha) \le \log(2pq/\alpha) + \log n \le \log(2pq/\alpha) +
\bn$, and $(\lamGone)^{2} \le (16\sigma^{2}/n)\{(2B_X^{2}C_U + \Lambda)\bn +
\Lambda\log(2pq/\alpha)\} \le C_{\pen}\sigma^{2}\bn/n$, using $\bn \ge 2$. On
$\Acal$, Lemma~\ref{lem:calib-blocks}(ii) gives $\lamz{G} \le \lamzbar{G}$ for
every $G$, so that $\lamG \ge \pen_{w}$ and, by the same lemma,
$\Ecal \subseteq \Tw$. On $\ESig$, $\lmin(\Gramhat) > \gamma/2$ and
$\lmaxx(\Gramhat) < \Lambda$; \eqref{eq:schur} gives $\gtil > \gamma/2$, and
$\lmin(\Amat\T\Amat/n) \ge \lmin(\Gramhat) > 0$ gives the rank. Finally
$\Prob(\smax^{2} \le 2B_X^{2}C_U) \to 1$ by Lemma~\ref{lem:columns}, which asks
$2^{J}\log d/n \to 0$, and $\Prob(\ESig) \to 1$ when
$pq2^{J_n}\log(pq2^{J_n})/n \to 0$ (Section~S1); the second condition implies
the first.
\end{proof}
}


% derivations/08-blocos.tex, prova do Lema 15 (numeracao global).
{\color{colR1}
\subsection*{Proof of Lemma 3}

Fix a block and write $\theta_{\lev} = \theta^{*}_{\cv\md,\lev\cdot} \in
\R^{2^{\lev}}$, set to zero for $\lev \ge J$, which only decreases the risk;
write $\Gcal_{\lev}$ for the chunks of a fine level $\lev$, those with
$2^{\lev} \ge b$, and $r_{\lev} = \sum_{G\in\Gcal_{\lev}}
\min(\norm{\theta_{G}}_{2}^{2},\eta)$. The coarse chunk, if there is one,
contributes at most $\eta$, so the block contributes at most
$\eta + \sum_{\lev}r_{\lev}$, the sum running over the fine levels. The first
bound on $r_{\lev}$ is the count: each chunk of $\Gcal_{\lev}$ has at least $b$
of the $2^{\lev}$ coefficients of the level, so that
$r_{\lev} \le |\Gcal_{\lev}|\eta \le 2^{\lev}\eta/b$.

\emph{The case $\pi \le 2$.} Write $a = \pi(s+1/2) > 1$. For $u \ge 0$,
$\min(u,\eta) \le \eta^{1-\pi/2}u^{\pi/2}$, and $\norm{x}_{2} \le \norm{x}_{\pi}$
for $\pi \le 2$; hence the second bound,
\[
  r_{\lev} \;\le\; \eta^{1-\pi/2}\sum_{G\in\Gcal_{\lev}}\norm{\theta_{G}}_{\pi}^{\pi}
  \;\le\; \eta^{1-\pi/2}\norm{\theta_{\lev}}_{\pi}^{\pi}
  \;\le\; \eta^{1-\pi/2}C_g^{\pi}2^{-\lev(a-1)},
\]
by Assumption~5, because the chunks of $\Gcal_{\lev}$ are disjoint and
contained in the level. Define $x \in \R$ by
$2^{xa} = bC_g^{\pi}\eta^{-\pi/2}$, so that the two bounds agree at
$\lev = x$, with common value
$h = 2^{x}\eta/b = \eta^{1-\pi/2}C_g^{\pi}2^{-x(a-1)}$. Summing the first bound
over the levels $\lev \le x$ gives at most $2^{\lfloor x\rfloor+1}\eta/b \le 2h$
(an empty sum if $x < 0$), and summing the second over the levels $\lev > x$
gives at most $h\sum_{i\ge0}2^{-i(a-1)} = h/(1-2^{1-a})$, since the first
integer $\lev > x$ has $2^{-\lev(a-1)} \le 2^{-x(a-1)}$. The block therefore
contributes at most $\eta + \{2 + (1-2^{1-a})^{-1}\}h$. Since $1/a = \tau/\pi$,
$2^{x} = b^{\tau/\pi}C_g^{\tau}\eta^{-\tau/2}$ and
$h = b^{\tau/\pi-1}C_g^{\tau}\eta^{1-\tau/2}$, and, with
$1-\tau/2 = 2s/(2s+1)$,
\[
  h \;=\; C_g^{\tau}\Bigl(\frac{\eta}{b}\Bigr)^{2s/(2s+1)}
  b^{\tau/\pi - 1 + 2s/(2s+1)}
  \;=\; C_g^{\tau}\Bigl(\frac{\eta}{b}\Bigr)^{2s/(2s+1)}
  b^{\tau(1/\pi-1/2)},
  \qquad
  \tau\Bigl(\frac{1}{\pi}-\frac12\Bigr) = \frac{2/\pi-1}{2s+1}.
\]

\emph{The case $\pi \ge 2$.} Hölder's inequality in $\R^{2^{\lev}}$ gives
$\norm{\theta_{\lev}}_{2} \le 2^{\lev(1/2-1/\pi)}\norm{\theta_{\lev}}_{\pi} \le
C_g2^{-\lev s}$, with $1/\pi = 0$ if $\pi = \infty$, and the second bound is
$r_{\lev} \le \norm{\theta_{\lev}}_{2}^{2} \le C_g^{2}2^{-2\lev s}$. With
$2^{x(2s+1)} = C_g^{2}b/\eta$ and $h = 2^{x}\eta/b = C_g^{2}2^{-2xs}$, the same
computation gives the contribution $\eta + \{2 + (1-2^{-2s})^{-1}\}h$, and
$h = (C_g^{2}b/\eta)^{1/(2s+1)}\eta/b = C_g^{\tau}(\eta/b)^{2s/(2s+1)}$, because
$2/(2s+1) = \tau$.

Summing over the $pq$ blocks gives the statement; at $\pi = 2$,
$1-2^{1-a} = 1-2^{-2s}$ and the two constants agree. \qed
}

We also record that, on the event of Proposition 2(ii) with $t = \gamma/2$,
\begin{equation}
  \label{eq:lmax}
  \normn{\Bmat\delta}^2 \le \lmaxx(\Gramhat)\norm{\delta}_2^2
  \quad\text{and}\quad \lmaxx(\Gramhat) \le \Lambda,
\end{equation}
the first because $\Bmat\T\Bmat/n$ is a principal submatrix of $\Gramhat$ and
the second by Weyl's inequality and Proposition 2(i).

% derivations/08-blocos.tex, a contagem (eq. contagem-lenta) e a prova do
% Corolario 14 (numeracao global); D47(c), D49(b).
{\color{colR1}
\begin{lemma}[The norm of the oracle over chunks]
\label{lem:count}
Under Assumptions~4 and 5, let $\Gcal$ be the partition of Section~2.3 of the
paper with $b \ge 2$, and let $\vartheta_{\pi} = \min(1-1/\pi,\,1/2)$, with
$\vartheta_{\infty} = 1/2$. Then, for every $J$,
\begin{equation}
  \label{eq:count}
  \nGone{\thetastar} \;\le\; pqC_g\Bigl\{\frac{1}{1-2^{-\seff}}
  + b^{-\vartheta_{\pi}}\sum_{\lev=\jst+1}^{J-1}2^{\lev(1/2-s)}\Bigr\}.
\end{equation}
The bound for each fine level is attained, by a level whose energy is spread
evenly when $\pi \ge 2$, and by a level with one large coefficient in each
chunk when $\pi \le 2$.
\end{lemma}

\begin{proof}
Fix a block and write $\theta_{\lev}$ for its coefficients of level $\lev$.
By Assumption~5, $\norm{\theta_{\lev}}_{2} \le
2^{\lev(1/2-1/\pi)_{+}}\norm{\theta_{\lev}}_{\pi} \le C_g2^{-\lev\seff}$, by
Hölder's inequality in $\R^{2^{\lev}}$ if $\pi \ge 2$ and by
$\norm{\cdot}_{2} \le \norm{\cdot}_{\pi}$ if $\pi \le 2$. The coarse chunk
contributes at most $\sum_{\lev\le\jst}\norm{\theta_{\lev}}_{2} \le
C_g(1-2^{-\seff})^{-1}$. At a fine level $\lev$ the $k_{\lev} \le 2^{\lev}/b$
chunks $G \in \Gcal_{\lev}$ partition the level. If $\pi \ge 2$, by the
Cauchy-Schwarz inequality,
\[
  \sum_{G\in\Gcal_{\lev}}\norm{\theta_{G}}_{2}
  \;\le\; k_{\lev}^{1/2}\norm{\theta_{\lev}}_{2}
  \;\le\; \Bigl(\frac{2^{\lev}}{b}\Bigr)^{1/2}C_g2^{-\lev s}
  \;=\; C_g\,b^{-1/2}\,2^{\lev(1/2-s)} ;
\]
if $\pi \le 2$, by $\norm{\theta_{G}}_{2} \le \norm{\theta_{G}}_{\pi}$ and
Hölder's inequality over the chunks,
\[
  \sum_{G\in\Gcal_{\lev}}\norm{\theta_{G}}_{2}
  \;\le\; k_{\lev}^{1-1/\pi}\Bigl(\sum_{G\in\Gcal_{\lev}}
  \norm{\theta_{G}}_{\pi}^{\pi}\Bigr)^{1/\pi}
  = k_{\lev}^{1-1/\pi}\norm{\theta_{\lev}}_{\pi}
  \;\le\; \Bigl(\frac{2^{\lev}}{b}\Bigr)^{1-1/\pi}C_g2^{-\lev(s+1/2-1/\pi)}
  = C_g\,b^{-(1-1/\pi)}\,2^{\lev(1/2-s)} .
\]
Both are $C_gb^{-\vartheta_{\pi}}2^{\lev(1/2-s)}$, and the equality cases of
the two inequalities are the configurations of the statement. Summing over the
fine levels $\jst < \lev < J$ and over the $pq$ blocks gives \eqref{eq:count}.
\end{proof}

\subsection*{Proof of Proposition 3}

\emph{The nonasymptotic bound.} On $\Tw$, Theorem~2(i) with the comparator
$\bar\bbeta = \betastar$, so that $\bbar = \bb$ and
$\vbar = v = \thetahat-\thetastar$, gives
$\normn{\Btil v}^{2} \le 6\pen\nGw{\thetastar} + 4\normn{\bb}^{2}$, and the
first part of Theorem~2(iii), with $(x+y+z)^{2} \le 3(x^{2}+y^{2}+z^{2})$, gives
\begin{equation}
  \label{eq:slow-nonasymptotic}
  \normn{\widehat{f}-f}^{2} \;\le\; 18\pen\nGw{\thetastar}
  + 24\normn{\bb}^{2} + 3\normn{\PA\bveps}^{2}
  \qquad\text{on } \Tw .
\end{equation}
Steps 1 to 3 and 6 of the proof of Theorem~2, the only ones used, do not
involve $\Gramhat$. With $\pen \ge \pen_{w}$, calibrated on the covariates by
Lemma~\ref{lem:calib-blocks}(i), the bound holds with probability at least
$1-\alpha$ in any design in which $\Amat$ has rank $p$, including designs with
$d > n$, where $\gtil = 0$.

\emph{The penalty level of the theory.} Let
$\Acal^{+} = \{\smax^{2} \le 2B_X^{2}C_U\}\cap\{\lmaxx(\Gramhat) \le \Lambda\}$.
On it Lemma~\ref{lem:calib-blocks}(ii) gives $\lamz{G} \le \lamzbar{G}$ for
every $G$, hence $\lamG \ge \pen_{w}$ and $\Ecal \subseteq \Tw$. The first
event has probability tending to one by Lemma~\ref{lem:columns}, whose proof
uses only $|X_{\cv}| \le B_X$ and $p_{\bU} \le C_U$. For the second,
$\Prob\{\lmaxx(\Gramhat) \le \Lambda\} \ge
\Prob\{\opnorm{\Gramhat-\Gram} < \Lambda - \kappa_2C_U\}$ by Weyl's inequality,
since $\lmaxx(\Gram) \le \kappa_2C_U$, and the matrix Bernstein bound in the
proof of Proposition~2(ii), with $t = \Lambda - \kappa_2C_U$, sends it to one;
that bound uses only $\norm{\bZ_{i}}_{2}^{2} \le R_J$ and
$\lmaxx(\Gram) \le \kappa_2C_U$, the upper half of Proposition~2(i), and
neither $c_U$ nor $\kappa_1$, so that any $\Lambda > \kappa_2C_U$ serves. Both
probabilities tend to one when $pq2^{J_n}\log(pq2^{J_n})/n \to 0$. Fix
$\epsilon > 0$, take $\alpha = \alpha' = \alpha'' = \epsilon/4$ and intersect
$\Acal^{+}$, $\Ecal$, $\{\normn{\bb}^{2} \le \mathcal{B}_{J_n}^{2}/\alpha'\}$
and $\{\normn{\PA\bveps}^{2} \le \sigma^{2}p/(n\alpha'')\}$, the last two by
Markov's inequality, Lemma~2 and Theorem~2(iii); for $n$ large the
intersection has probability at least $1-\epsilon$, and on it
$\normn{\widehat{f}-f}^{2} \le 18\lamG\nGw{\thetastar} +
24\mathcal{B}_{J_n}^{2}/\alpha' + 3\sigma^{2}p/(n\alpha'')$, which is the first
display of Proposition~3. For the chain, Lemma~\ref{lem:calib-blocks}(iii)
gives $\lamG w_{G} \le \varrho\lamGone$ for every $G$, and summing with the
weights $\norm{\thetastar_{G}}_{2}$ gives
$\lamG\nGw{\thetastar} \le \varrho\lamGone\nGone{\thetastar}$, with
$\varrho \le \bar\varrho$; and $\norm{x}_{2} \le \norm{x}_{1}$ in each chunk
gives $\nGone{\thetastar} \le \norm{\thetastar}_{1}$. The order of $\lamGone$
is \eqref{eq:order-lambda}.

\emph{Part (i).} With $s > 1/2$ and $2^{\jst+1} \ge b$,
$\sum_{\lev>\jst}2^{\lev(1/2-s)} \le 2^{(\jst+1)(1/2-s)}/(1-2^{1/2-s}) \le
\bn^{1/2-s}/(1-2^{1/2-s})$, and Lemma~\ref{lem:count} gives
$\nGone{\thetastar} \le pqC_g\{(1-2^{-\seff})^{-1} +
\bn^{1/2-s-\vartheta_{\pi}}(1-2^{1/2-s})^{-1}\} = O(1)$; hence
$\lamG\nGw{\thetastar} = O\{(\bn/n)^{1/2}\} = O\{(\log n/n)^{1/2}\}$. With
$J_n \ge \log_2 n/(4\seff)$, $\mathcal{B}_{J_n}^{2} \asymp 2^{-2J_n\seff} \le
n^{-1/2}$, and $p/n$ is of smaller order. With $\seff > 1/4$ the choice
$J_n = \lceil\log_2 n/(4\seff)\rceil$ has $2^{J_n} < 2n^{1/(4\seff)}$ and
satisfies the conditions of the statement.

\emph{Part (ii).} With $s < 1/2$,
$\sum_{\jst<\lev<J_n}2^{\lev(1/2-s)} \le 2^{J_n(1/2-s)}/(2^{1/2-s}-1)$, so that
$\nGone{\thetastar} = O\{1 + \bn^{-\vartheta_{\pi}}2^{J_n(1/2-s)}\}$ and, since
$1/2 - \vartheta_{\pi} = (1/\pi-1/2)_{+}$,
\[
  \lamG\nGw{\thetastar} \;\lesssim\; \Bigl(\frac{\bn}{n}\Bigr)^{1/2}
  + n^{-1/2}\,\bn^{(1/\pi-1/2)_{+}}\,2^{J_n(1/2-s)} .
\]
Equating the second term with the bias $2^{-2J_n\seff}$ gives
$2^{J_n(1/2-s+2\seff)} \asymp n^{1/2}\bn^{-(1/\pi-1/2)_{+}}$, which is the
choice of the statement, with common value $2^{-2J_n\seff} \asymp
\{\bn^{(2/\pi-1)_{+}}/n\}^{2\seff/(1-2s+4\seff)}$. The term $(\bn/n)^{1/2}$ and
$p/n$ are of smaller order, because $2\seff/(1-2s+4\seff) < 1/2$ is equivalent
to $s < 1/2$. The conditions of the statement hold because $\seff > s/2$ gives
$1/(1-2s+4\seff) < 1$, and $2^{J_n} \ge \bn$ for $n$ large. When $\pi \ge 2$,
$\seff = s$, $(2/\pi-1)_{+} = 0$ and $1-2s+4s = 2s+1$. \qed

% derivations/08-blocos.tex, Corolario 14(iii) e (iv), a equacao mu-blocos e a
% prova dos dois itens (E1.14; E5e, D52).
\subsection*{Proposition 3 above the concentration ceiling}

The conditions $\seff > 1/4$ in Proposition~3(i) and $\seff > s/2$ in
Proposition~3(ii) keep the resolution of the balance below the ceiling
$2^{J_n} \asymp n/\log n$, under which $\smax^{2} \le 2B_X^{2}C_U$ and
$\lmaxx(\Gramhat) \le \Lambda$ hold with probability tending to one. The pairs
that they exclude, all with $\pi < 2$, and the boundary $s = 1/2$ call for
$2^{J_n}$ above that ceiling, where neither bound holds: the wavelets of the
finest level take values up to $2^{(J-1)/2}\norm{\psi}_{\infty}$ and have one
or two sample points in their support, so that $\smax^{2}$ and
$\lmaxx(\Gramhat)$ grow with $2^{J}/n$, and the deterministic level (3.1) is
no longer admissible. With $D = p + d$, $\xi = \Lambda - \kappa_2C_U > 0$,
$C_\psi$ as in Proposition~2(ii) and $\norm{\psi}_{\infty}$ as in
Lemma~\ref{lem:columns}, write
\begin{equation}
  \label{eq:mu-blocks}
  \mu^{\Gcal}_{J} = 1 + c_{\mu}\frac{2^{J}\log(2D)}{n},
  \qquad
  \lamGp = (\mu^{\Gcal}_{J_n})^{1/2}\lamG,
\end{equation}
with
\[
  c_{\psi} = \frac{7\norm{\psi}_{\infty}^{2}}{12C_U},
  \qquad
  c_{\mu} = \max\Bigl\{c_{\psi},\
  \frac{2pB_X^{2}(1+qC_\psi)}{\Lambda}
  \Bigl(\frac{\kappa_2C_U}{2\xi} + \frac43\Bigr)\Bigr\}.
\]
When $pq2^{J_n}\log(pq2^{J_n})/n \to 0$, $\lamGp = \lamG\{1 + o(1)\}$.

\begin{proposition}[Proposition 3 above the concentration ceiling]
\label{prop:slow-ceiling}
Let the assumptions of Proposition~3 hold, with the conditions $2^{J_n} \le 2n$
and $pq2^{J_n}\log(pq2^{J_n})/n \to 0$ replaced by $2^{J_n} \le n^{2}$, and
with $\lamGp$ of \eqref{eq:mu-blocks} in place of $\lamG$ in the estimator.
Then
\[
  \normn{\widehat{f}-f}^{2} = O_p\bigl\{\lamGp\nGw{\thetastar} +
  \mathcal{B}_{J_n}^{2} + p/n\bigr\},
\]
where $\lamGp\nGw{\thetastar} \le \bar\varrho\,(\mu^{\Gcal}_{J_n})^{1/2}
\lamGone\nGone{\thetastar}$,
with $(\lamGone)^{2} \le (C_{\pen} + 16\Lambda)\sigma^{2}\bn/n$ and $C_{\pen}$
as in \eqref{eq:order-lambda}, and:
\begin{enumerate}
\item[(iii)] (the pairs that Proposition~3(i) and (ii) exclude, all with
$\pi < 2$) (a) if $s > 1/2$ and $\seff \le 1/4$, the choice
$2^{J_n} \asymp (n/\log n)^{2/(1+4\seff)}$ gives
\[
  \normn{\widehat{f}-f}^{2} = O_p\bigl\{(\log n/n)^{4\seff/(1+4\seff)}\bigr\};
\]
(b) if $s < 1/2$ and $\seff \le s/2$, the choice
$2^{J_n} \asymp (n/\bn^{1/\pi})^{1/(1-s+2\seff)}$ gives
\[
  \normn{\widehat{f}-f}^{2}
  = O_p\bigl\{n^{-2\seff/(1-s+2\seff)}
  (\log n)^{(1/\pi)2\seff/(1-s+2\seff)}\bigr\};
\]
in both, $2^{J_n} \gtrsim n/\log n$, and the powers of $n$ join those of
Proposition~3(i) and (ii) at $\seff = 1/4$ and $\seff = s/2$;
\item[(iv)] (the boundary $s = 1/2$, with $\pi > 1$)
$\nGone{\thetastar} \le pqC_g\{(1-2^{-\seff})^{-1} + \bn^{-\vartheta_{\pi}}J_n\}$,
and: if $\seff \ge 1/4$, that is $\pi \ge 4/3$, the choice
$2^{J_n} \asymp (n/\bn)^{1/(4\seff)}$ gives
$\normn{\widehat{f}-f}^{2} = O_p\{n^{-1/2}(\log n)^{3/2-\vartheta_{\pi}}\}$,
which is $n^{-1/2}\log n$ when $\pi \ge 2$; if $\seff < 1/4$, that is
$1 < \pi < 4/3$, the choice
$2^{J_n} \asymp \{n/\bn^{1+1/\pi}\}^{2/(1+4\seff)}$ gives
$\normn{\widehat{f}-f}^{2} =
O_p[\{(\log n)^{1+1/\pi}/n\}^{4\seff/(1+4\seff)}]$.
\end{enumerate}
The proof uses no lower bound on $\Gramhat$, and no property of $\Gramhat$
other than the upper bound $\lmaxx(\Gramhat) \le \Lambda\mu^{\Gcal}_{J_n}$, which
holds with probability tending to one under the upper bounds of Assumptions~2
and~3 alone, with no condition on $J_n$.
\end{proposition}

Stopping $J_n$ at the ceiling instead, with $2^{J_n} = n/(K\log n)$ and $K$
fixed and large enough, keeps $\lamG$ admissible and gives
$O_p\{(\log n/n)^{2\seff}\}$ in the pairs of (iii); above the ceiling the
estimation term grows more slowly than the bias decreases, and the exponents
of (iii) are strictly larger than $2\seff$ away from $\seff = 1/4$ and
$\seff = s/2$ (Remark~\ref{rem:ceiling} in Section~S8, for the lasso). Compared
with Proposition~\ref{prop:lasso-slow}(iii) and (iv) for the lasso, the chunks
remove nothing in (iii)(a), where the coarse chunk dominates both norms, as in
Proposition~3(i); the factor $(\log n)^{(1-1/\pi)2\seff/(1-s+2\seff)}$ in
(iii)(b); and, in (iv), $(\log n)^{\vartheta_{\pi}}$ when $\seff \ge 1/4$, half
a logarithm when $\pi \ge 2$, and $(\log n)^{(1-1/\pi)4\seff/(1+4\seff)}$ when
$\seff < 1/4$. This is the mechanism of Proposition~3(ii): the logarithm comes
from the size of the chunks, and the count of Lemma~\ref{lem:count} returns it
to the extent that the signal is spread over them.

\subsection*{Proof of Proposition~\ref{prop:slow-ceiling}}

\emph{The two events above the ceiling.} Let
$\Acal^{++} = \{\smax^{2} \le 2B_X^{2}C_U\mu^{\Gcal}_{J}\}\cap
\{\lmaxx(\Gramhat) \le \Lambda\mu^{\Gcal}_{J}\}$. For the first, the Bernstein
bound of Lemma~\ref{lem:columns} holds for every $J$ and every $t > 0$. With
$x = 2\log d$ and $t = (2xR'_JB_X^{2}C_U/n)^{1/2} + 2xR'_J/(3n)$,
$nt^{2} \ge 2xR'_JB_X^{2}C_U + (2/3)xR'_Jt$, so that the exponent is at most
$-x$; and $(2ab)^{1/2} \le a + b/2$ with $a = B_X^{2}C_U$ and $b = xR'_J/n$
gives $t \le B_X^{2}C_U + (7/6)xR'_J/n$. Since
$2B_X^{2}C_U + (7/3)R'_J\log d/n = 2B_X^{2}C_U(1 + c_{\psi}2^{J}\log d/n)$,
\begin{equation}
  \label{eq:smax-ceiling}
  \Prob\bigl\{\smax^{2} > 2B_X^{2}C_U(1 + c_{\psi}2^{J}\log d/n)\bigr\}
  \;\le\; \frac{1}{d}
  \qquad\text{for every } J,
\end{equation}
and $c_{\psi}2^{J}\log d \le c_{\mu}2^{J}\log(2D)$. For the second, the matrix
Bernstein bound in the proof of Proposition~2(ii),
$\Prob(\opnorm{\Gramhat-\Gram} \ge t) \le
2D\exp\{-nt^{2}/(2R_J(\kappa_2C_U + 2t/3))\}$, holds for every $J$ and every
$t > 0$. With $x = 2\log(2D)$ and
$t = (2xR_J\kappa_2C_U/n)^{1/2} + (4/3)xR_J/n$,
$nt^{2} \ge 2xR_J\kappa_2C_U + (4/3)xR_Jt$, and the probability is at most
$2De^{-x} = 1/(2D)$; and $(2ab)^{1/2} \le \xi + ab/(2\xi)$ with
$a = \kappa_2C_U$ and $b = xR_J/n$ gives
$t \le \xi + \{\kappa_2C_U/(2\xi) + 4/3\}xR_J/n$. By Weyl's inequality and
$\lmaxx(\Gram) \le \kappa_2C_U$, off that event
$\lmaxx(\Gramhat) < \Lambda + \{\kappa_2C_U/(2\xi) + 4/3\}2\log(2D)R_J/n \le
\Lambda\mu^{\Gcal}_{J}$, because
$R_J = pB_X^{2}(1 + qC_\psi2^{J}) \le pB_X^{2}(1+qC_\psi)2^{J}$. Hence
$\Prob(\Acal^{++}) \ge 1 - 1/d - 1/(2D) \to 1$, with no condition on $J_n$.
Both bounds use only $|X_{\cv}| \le B_X$, $p_{\bU} \le C_U$ and
$\lmaxx(\Gram) \le \kappa_2C_U$.

\emph{The general bound.} On $\Acal^{++}$ the two bounds in the proof of
Lemma~\ref{lem:calib-blocks}(ii), $\mathrm{tr}(\Psitil_{G}) \le |G|\smax^{2}$
and $\opnorm{\Psitil_{G}} \le \lmaxx(\Gramhat)$, give
$\lamz{G} \le (\mu^{\Gcal}_{J})^{1/2}\lamzbar{G}$ for every $G$, so that
$\lamGp \ge \pen_{w}$ and $\Ecal \subseteq \Tw$, and the proof of
Proposition~3 goes through with $\lamGp$ in place of $\lamG$ and $\Acal^{++}$
in place of $\Acal^{+}$; the chain gives
$\lamGp\nGw{\thetastar} \le \bar\varrho(\mu^{\Gcal}_{J_n})^{1/2}
\lamGone\nGone{\thetastar}$. With $2^{J_n} \le n^{2}$,
$|\Gcal_{n}| \le pq(1 + n^{2}/\bn) \le pqn^{2}$ and
$\log(|\Gcal_{n}|/\alpha) \le \log(2pq/\alpha) + 2\bn$, and the computation in
the proof of Lemma~\ref{lem:regime} gives
$(\lamGone)^{2} \le (16\sigma^{2}/n)\{(2B_X^{2}C_U + 2\Lambda)\bn +
\Lambda\log(2pq/\alpha)\} \le (C_{\pen} + 16\Lambda)\sigma^{2}\bn/n$. Since
$\log(2D_n) \asymp \log n \asymp \bn$,
\begin{equation}
  \label{eq:lam-plus-blocks}
  (\mu^{\Gcal}_{J_n})^{1/2}\lamGone \;\asymp\;
  \Bigl(\frac{\log n}{n}\Bigr)^{1/2}
  \Bigl(1 + \frac{2^{J_n}\log n}{n}\Bigr)^{1/2},
\end{equation}
which is of order $2^{J_n/2}\log n/n$ when $2^{J_n} \gtrsim n/\log n$.

\emph{Part (iii)(a).} With $s > 1/2$, $\nGone{\thetastar} = O(1)$ as in the
proof of Proposition~3(i). With
$2^{J_n} \asymp (n/\log n)^{2/(1+4\seff)} \gtrsim n/\log n$, because
$\seff \le 1/4$, \eqref{eq:lam-plus-blocks} gives the estimation term
$(\log n/n)(n/\log n)^{1/(1+4\seff)} = (\log n/n)^{4\seff/(1+4\seff)}$, the
order of the bias $2^{-2J_n\seff}$; $p/n$ is of smaller order.

\emph{Part (iii)(b).} With $s < 1/2$ and $\pi < 2$,
$\vartheta_{\pi} = 1 - 1/\pi$ and
$\nGone{\thetastar} = O\{1 + \bn^{-\vartheta_{\pi}}2^{J_n(1/2-s)}\}$ as in the
proof of Proposition~3(ii). With
$2^{J_n} \asymp (n/\bn^{1/\pi})^{1/(1-s+2\seff)}$, which exceeds $n/\bn$ for
$n$ large, because the exponent is at least $1$ and $1/\pi \le 1$, and which is
at most $n^{2}$, because $1-s+2\seff > 1/2$, \eqref{eq:lam-plus-blocks} gives
the estimation term
\[
  2^{J_n/2}\frac{\log n}{n}\bigl\{1 + \bn^{-(1-1/\pi)}2^{J_n(1/2-s)}\bigr\}
  \;\asymp\; 2^{J_n/2}\frac{\log n}{n} +
  2^{J_n(1-s)}\frac{(\log n)^{1/\pi}}{n},
\]
whose second term dominates and, at the choice of the statement, equals
$(\bn^{1/\pi}/n)^{1-(1-s)/(1-s+2\seff)} =
(\bn^{1/\pi}/n)^{2\seff/(1-s+2\seff)}$, the order of the bias
$2^{-2J_n\seff}$.

\emph{Part (iv).} With $s = 1/2$, the sum in \eqref{eq:count} is at most
$J_n$, which gives the bound on $\nGone{\thetastar}$ of the statement, and
$J_n \asymp \log n \asymp \bn$ in both choices, so that
$\nGone{\thetastar} \asymp (\log n)^{1-\vartheta_{\pi}}$, since
$\vartheta_{\pi} \le 1/2$. If $\seff \ge 1/4$,
$2^{J_n} \asymp (n/\bn)^{1/(4\seff)} \lesssim n/\log n$,
$\mu^{\Gcal}_{J_n}$ is bounded, the estimation term is
$(\log n/n)^{1/2}(\log n)^{1-\vartheta_{\pi}}$, and the bias, of order
$(\bn/n)^{1/2}$, is of smaller order. If $\seff < 1/4$, then $\pi < 4/3$ and
$\vartheta_{\pi} = 1 - 1/\pi$; with
$2^{J_n} \asymp (n/\bn^{1+1/\pi})^{2/(1+4\seff)}$, which exceeds $n/\log n$,
the estimation term is $2^{J_n/2}(\log n)^{1+1/\pi}/n \asymp
\{(\log n)^{1+1/\pi}/n\}^{4\seff/(1+4\seff)}$, the order of the bias. \qed

% derivations/08-blocos.tex, prova do Teorema 4 (numeracao global).
\subsection*{Proof of Theorem 3}

By the definition of $J_n$, $n^{1/(2\seff+1)} \le 2^{J_n} <
2n^{1/(2\seff+1)} \le 2n$, so the setting of this section holds, and
$pq2^{J_n}\log(pq2^{J_n})/n \asymp n^{1/(2\seff+1)-1}\log n =
n^{-2\seff/(2\seff+1)}\log n \to 0$, which is the requirement of
Proposition~2(ii); by Lemma~\ref{lem:regime}, $\Prob(\Acal) \to 1$. Fix
$\epsilon > 0$, take $\alpha = \alpha' = \alpha'' = \epsilon/4$, and intersect
$\Acal$, the event $\Ecal$ of Lemma~\ref{lem:calib-blocks}(i),
$\{\normn{\bb}^{2} \le \mathcal{B}_{J_n}^{2}/\alpha'\}$ (Markov's inequality,
since $\E\normn{\bb}^{2} = \norm{f-f_{J_n}}_{L_2(P)}^{2} \le
\mathcal{B}_{J_n}^{2}$ by Lemma~2) and
$\{\normn{\PA\bveps}^{2} \le \sigma^{2}p/(n\alpha'')\}$ (Markov's inequality and
Theorem~2(iii)); for $n$ large the intersection has probability at least
$1-\epsilon$. On it $\Tw$ holds, and \eqref{eq:oracle-min-blocks} with
$\bar\bbeta = \betastar$ gives
\[
  \normn{\widehat{f}-f}^{2} \;\le\;
  \frac{192(\lamG)^{2}\Wcal(\Gzero(\thetastar))}{\gtil}
  + 60\normn{\bb}^{2} + 3\normn{\PA\bveps}^{2}.
\]
By Lemma~\ref{lem:calib-blocks}(iii), Lemma~\ref{lem:regime} and
Proposition~\ref{prop:balanced}(ii), with $\gtil > \gamma/2$,
\[
  \frac{192(\lamG)^{2}\Wcal(\Gcal_{n})}{\gtil}
  \;\le\; \frac{384\bar\varrho^{2}(\lamGone)^{2}|\Gcal_{n}|}{\gamma}
  \;\le\; \frac{384\bar\varrho^{2}C_{\pen}\sigma^{2}pq(\bn + 2^{J_n})}{\gamma n},
\]
which is of order $n^{-2\seff/(2\seff+1)}$, because
$2^{J_n} \asymp n^{1/(2\seff+1)}$ dominates $\bn \asymp \log n$. The bias term
is $60\mathcal{B}_{J_n}^{2}/\alpha' \asymp 2^{-2J_n\seff} \le
n^{-2\seff/(2\seff+1)}$, by the lower bound on $2^{J_n}$, and the last term is
$O(1/n)$, of strictly smaller order. Hence
$\Prob\{\normn{\widehat{f}-f}^{2} > M_\epsilon n^{-2\seff/(2\seff+1)}\} <
\epsilon$ for a constant $M_\epsilon$. \qed

% derivations/08-blocos.tex, prova do Corolario 10 (numeracao global).
\subsection*{Proof of Corollary 1}

By the orthonormality of the wavelets of level below $J_n$ in $L_2[0,1]$ and
Lemma~2,
$\sum_{\cv,\md}\lVert\widehat{g}_{\cv\md}-\gcomp{\cv}{\md}\rVert^{2}_{L_2[0,1]}
\le 2\norm{\thetahat-\thetastar}_{2}^{2} +
2pqC_g^{2}(1-2^{-2\seff})^{-1}2^{-2J_n\seff}$. On the event of the proof of
Theorem~3, Theorem~2(ii) with $\bar\bbeta = \betastar$ gives
$\norm{\thetahat-\thetastar}_{2}^{2} \le
64(\lamG)^{2}\Wcal(\Gcal_{n})/\gtil^{2} + 16\normn{\bb}^{2}/\gtil$, with
$\gtil > \gamma/2$: the two terms have the order of those of Theorem~3, with
$\gamma^{-2}$ in place of $\gamma^{-1}$. For the levels,
Lemma~\ref{lem:profile} with $\bar\bbeta = \betastar$ gives
$\Amat(\chat-\bc) = \PA(\bb+\bveps) - \PA\Bmat v$, so that
$\normn{\Amat(\chat-\bc)} \le \normn{\bb} + \normn{\PA\bveps} +
\Lambda^{1/2}\norm{v}_{2}$, and $\lmin(\Amat\T\Amat/n) > \gamma/2$ gives
$\norm{\chat-\bc}_{2}^{2} \le 2\normn{\Amat(\chat-\bc)}^{2}/\gamma$. Finally the
components of a functional coefficient depend on distinct coordinates of $u$
and the basis is orthonormal, so that
$\lVert\widehat{\fcoef{\cv}}-\fcoef{\cv}\rVert_{L_2([0,1]^q)} \le
|\widehat{c}_{\cv}-\lvl{\cv}| + \sum_{\md}\lVert\widehat{g}_{\cv\md}-
\gcomp{\cv}{\md}\rVert_{L_2[0,1]}$, with $q$ fixed. \qed

% derivations/08-blocos.tex, prova do Corolario 11 (numeracao global); D47(a).
\subsection*{Proof of Theorem 1}

\emph{The window.} Since $\seff \le s$,
$s/\{(2s+1)\seff\} \ge 1/(2s+1)$, and the window is nonempty if and only if
$s/\{(2s+1)\seff\} < 1$, that is $\seff > s/(2s+1)$. With $c < 1$,
$2^{J_n} \le 2n^{c} \le 2n$, so the setting of this section holds; the
requirement of Proposition~2(ii) is $pq2^{J_n}\log(pq2^{J_n})/n \asymp
n^{c-1}\log n \to 0$, and $\Prob(\Acal) \to 1$ by Lemma~\ref{lem:regime}.

\emph{The rate.} Fix $\epsilon > 0$ and the event of the proof of Theorem~3.
On it, Corollary~\ref{cor:ideal} with $\Lambda' = \Lambda$ gives
$\normn{\widehat{f}-f}^{2} \le 120\Lambda\Rid(\thetastar;\eta) +
120\normn{\bb}^{2} + 3\normn{\PA\bveps}^{2}$, and, by
Lemma~\ref{lem:calib-blocks}(iii), \eqref{eq:order-lambda} and
$\gtil > \gamma/2$,
\[
  \eta \;=\; \frac{1.6\max_{G}(\lamG w_{G})^{2}}{\Lambda\gtil}
  \;\le\; \frac{3.2\bar\varrho^{2}(\lamGone)^{2}}{\Lambda\gamma}
  \;\le\; \bar\eta_{n} := \kappa\,\frac{\bn}{n},
  \qquad
  \kappa = \frac{3.2\bar\varrho^{2}C_{\pen}\sigma^{2}}{\Lambda\gamma}.
\]
Since $\Rid$ is nondecreasing in $\eta$, Lemma~3 with $b = \bn$ gives
\[
  \Rid(\thetastar;\eta) \;\le\; \Rid(\thetastar;\bar\eta_{n})
  \;\le\; pq\Bigl\{\AG C_g^{\tau}\kappa^{2s/(2s+1)}\,n^{-2s/(2s+1)}
  \bn^{(2/\pi-1)_{+}/(2s+1)} + \kappa\frac{\bn}{n}\Bigr\},
\]
with $\bn \asymp \log n$; the second term is of smaller order because
$2s/(2s+1) < 1$. The bias is $O(2^{-2J_n\seff}) = O(n^{-2c\seff})$, and
$2c\seff \ge 2s/(2s+1)$ by the window; the term $\sigma^{2}p/n$ is of smaller
order. For the components, the second bound of Corollary~\ref{cor:ideal} and
Lemma~2 give the same rate, as in the proof of Corollary~1.

% E5g (pergunta 51(c), D57): a uniformidade do O_p, lida na prova do
% Corolario 16 de derivations/09-cota-inferior.tex; sem ela, "minimax optimal"
% no corpo fica no sentido ponto a ponto.
\emph{Uniformity over the class.} The events $\Acal$ and $\Ecal$ and the
event for $\normn{\PA\bveps}^{2}$ do not depend on $f$; the event for
$\normn{\bb}^{2}$ does, but its probability is at least $1-\alpha'$ for every
$f$ whose components satisfy Assumption~5, because the bound
$\mathcal{B}_{J_n}^{2}$ on $\E\normn{\bb}^{2}$ of Lemma~2 depends on $f$ only
through $C_g$. On the event, the bound depends on $f$ only through $C_g$, by
Lemmas~2 and~3, and not on the levels $\bc$, since $\bb = f - f_{J_n}$ does
not involve them. Hence
$\sup_{f}\Prob_{f}\{\normn{\widehat{f}-f}^{2} > M_\epsilon r_n\} < \epsilon$
for $n$ large, with $\Prob_{f}$ the law of the sample when the regression
function is $f$ and the supremum over the functions of Assumption~5 with any
levels, and the same holds for the components: the bound of Theorem~1 is
uniform over the class, which is the form that Theorem~\ref{thm:lower} below
is compared with.

\emph{Two readings.} The constant grows with the weights at most as
$\bar\varrho^{\,4s/(2s+1)}$ in the leading term, through
$\kappa^{2s/(2s+1)}$. The gain over the rate $(\log n/n)^{2s/(2s+1)}$ of the
lasso (Theorem~\ref{thm:lasso-main}) is the factor $(\log n)^{2\seff/(2s+1)}$,
because $2s/(2s+1) - (2/\pi-1)_{+}/(2s+1) = 2\seff/(2s+1)$: when $\pi < 2$,
$2s - 2/\pi + 1 = 2(s - 1/\pi + 1/2) = 2\seff$, and when $\pi \ge 2$,
$\seff = s$. The logarithm of the union bound over the columns is absorbed by
the size of the chunks, which is why they must have at least of the order of
$\log|\Gcal_{n}|$ coefficients. \qed

% derivations/08-blocos.tex, Observacao rem:inferior.
% E5g (D57): removida inteira e substituida pelo Theorem S6.1 abaixo; fica no
% lugar em cinza e sem numero (D51(b)), sem o rotulo rem:lower, que nada cita.
{\color{gray}\removedk
\begin{remarkold}[The lower bound of Klopp and Pensky with one modulator]
For $q = 1$, $\bX$ independent of $U_1$ and Gaussian errors,
\citet[Theorem~1, eqs.~3.6 and 3.7]{Klopp-Pensky-2015}, with the number of
covariates fixed, give the minimax lower bound $n^{-2r/(2r+1)}$ over their
class (A3) with the periodized wavelet basis, which their condition (A1)(b)
admits. That class, a weighted $\ell_\nu$ ball
$\sum_{k}|a_{k}|^{\nu}(k+1)^{\nu(r+1/2-1/\nu)} \le C_a^{\nu}$ over all the
coefficients, implies Assumption~5 with $(s,\pi) = (r,\nu)$ and $C_g = C_a$,
level by level, since $k+1 \ge 2^{\lev}$ at level $\lev$; so the lower bound
holds over the class of Assumption~5, and when $\pi \ge 2$ Theorem~1 is minimax
optimal up to a constant in that case. When $\pi < 2$ the factor
$(\log n)^{(2/\pi-1)/(2s+1)}$ remains, the same as in eqs.~(3.19) and (3.20)
of Klopp and Pensky with the number of covariates fixed, and the same as in
the rate of block thresholding in the sequence model
\citep[Theorem~4, eq.~5.4]{Cai-1999}, which holds for $s \ge 1/\pi$; the
condition $\seff > s/(2s+1)$ of Theorem~1 is implied by $s \ge 1/\pi$, since
then $\seff \ge 1/2 > s/(2s+1)$, and is strictly weaker. That this factor is
necessary for the estimator is proved neither here nor by Klopp and Pensky,
and for $q \ge 2$ no lower bound is available.
\end{remarkold}
}

% E5g (E1.9; D57): derivations/09-cota-inferior.tex, Teorema 5 e Lema 17
% (numeracao global), nas tres perdas, com as provas; dentro do grupo azul de
% k = 4. Os simbolos locais do 09 (o hipercubo C, Omega, zeta, d_H, a funcao
% de limiar) ficam escritos no texto; |C| no lugar do M do 09, que colide com
% M_epsilon e com o M da S8; L_n e L_P escritos por extenso, porque L_n e o
% logaritmo da S8.
\subsection*{A minimax lower bound}

The rates of this section are upper bounds. For a lower bound, fix any law
$P$ of $(\bX,\bU)$ that satisfies Assumptions~2 and~3, with $\bX$ allowed to
depend on $\bU$, and let the errors be Gaussian,
\begin{equation}
  \label{eq:gauss}
  \varepsilon_i \sim N(0,\sigma^{2}),\ \text{independent of } (\bX_i,\bU_i),
\end{equation}
which satisfy Assumption~1. Write $\Prob_f$ and $\E_f$ for the law and the
expectation of the sample when the regression function is $f$, and
$\mathcal{F} = \mathcal{F}(s,\pi,C_g)$ for the set of the functions
$f(x,u) = \sum_{\cv}x_{\cv}\{\lvl{\cv} + \sum_{\md}\gcomp{\cv}{\md}(u_{\md})\}$
with $\bc \in \R^{p}$ arbitrary and every component satisfying the centering
and the bound of Assumption~5; by Proposition~1, $f$ determines its levels and
its components. An estimator is any measurable function of the sample: for
prediction, $\widehat{f}$ with values in $L_2(P)$, or only its values
$\widehat{f}(\bX_i,\bU_i)$, and for the components,
$\widehat{g} = (\widehat{g}_{\cv\md})$ with $\widehat{g}_{\cv\md} \in
L_2[0,1]$. For a nonempty set $\Kcal$ of blocks $(\cv,\md)$, the loss for the
components in $\Kcal$ is
\[
  L_{\Kcal}(\widehat{g},f) \;=\; \sum_{(\cv,\md)\in\Kcal}
  \lVert\widehat{g}_{\cv\md} - \gcomp{\cv}{\md}\rVert_{L_2[0,1]}^{2} .
\]
With $\gamma = \kappa_1c_U$ as in Section~S1 and
$\cA = (1-2^{-1/2})/2 \approx 0.146$, put, for
$n \ge n_{0} = \lceil 2\sigma^{2}/(C_g^{2}B_X^{2}C_U)\rceil$,
\begin{equation}
  \label{eq:lower-choices}
  \begin{gathered}
  \zeta_n^{2} = \frac{2\sigma^{2}}{nB_X^{2}C_U},
  \qquad
  \lev_n = \max\Bigl\{\lev \ge 0 : 2^{\lev(2s+1)} \le \frac{C_g^{2}}{\zeta_n^{2}}\Bigr\},
  \\
  \rlow^{2} = C_g^{2/(2s+1)}\Bigl(\frac{\sigma^{2}}{nB_X^{2}C_U}\Bigr)^{2s/(2s+1)} ;
  \end{gathered}
\end{equation}
$\lev_n$ exists because $C_g^{2}/\zeta_n^{2} \ge 1$ when $n \ge n_{0}$, and
$\rlow^{2}$, of order $n^{-2s/(2s+1)}$, is the rate of the lower bound.

\begin{theorem}[Minimax lower bound]
\label{thm:lower}
Let $P$ be any law of $(\bX,\bU)$ that satisfies Assumptions~2 and~3, with
$q \ge 1$ and $\bX$ possibly dependent on $\bU$, let the basis satisfy
Assumption~4, and let the errors be Gaussian, as in \eqref{eq:gauss}. For
every $s > 0$, every $1 \le \pi \le \infty$, every nonempty set $\Kcal$ of
blocks and every $n \ge n_{0}$, with $\rlow$ as in \eqref{eq:lower-choices}:
\begin{enumerate}
\item[(i)] (components)
\begin{gather*}
  \inf_{\widehat{g}}\sup_{f\in\mathcal{F}}\E_{f}L_{\Kcal}(\widehat{g},f)
  \;\ge\; \frac{\cA}{8}\,|\Kcal|\,\rlow^{2},
  \\
  \inf_{\widehat{g}}\sup_{f\in\mathcal{F}}\Prob_{f}\Bigl\{L_{\Kcal}(\widehat{g},f)
  \ge \frac{\cA}{16}\,|\Kcal|\,\rlow^{2}\Bigr\} \;\ge\; \frac{\cA}{2} ;
\end{gather*}
\item[(ii)] (prediction, population norm) the same holds for
$\norm{\widehat{f}-f}_{L_2(P)}^{2}$, with $\gamma|\Kcal|\rlow^{2}$ in place of
$|\Kcal|\rlow^{2}$;
\item[(iii)] (prediction, empirical norm) if moreover $n \ge n_{1}$, the
smallest integer from which on
$2D\exp[-n\gamma^{2}/\{8R_{J}(\kappa_2C_U + \gamma/3)\}] \le \cA/4$, with
$J = \lev_n + 1$ and $D = p + pq\NJ$ and $R_J$ as in Proposition~2(ii), then
\begin{gather*}
  \inf_{\widehat{f}}\sup_{f\in\mathcal{F}}\E_{f}\normn{\widehat{f}-f}^{2}
  \;\ge\; \frac{\cA}{32}\,\gamma|\Kcal|\,\rlow^{2},
  \\
  \inf_{\widehat{f}}\sup_{f\in\mathcal{F}}\Prob_{f}\Bigl\{\normn{\widehat{f}-f}^{2}
  \ge \frac{\cA}{32}\,\gamma|\Kcal|\,\rlow^{2}\Bigr\} \;\ge\; \frac{\cA}{4} .
\end{gather*}
\end{enumerate}
The constants depend neither on $\pi$, $r$, $p$ or $q$, nor on $P$ beyond
$(B_X,C_U)$ in $\rlow$ and $\gamma = \kappa_1c_U$; $n_{0}$ depends only on
$(\sigma,C_g,B_X,C_U)$, and $n_{1}$, which is finite, also on
$(p,q,s,\kappa_1,\kappa_2,c_U)$ and the wavelet family. In particular, with
$\Kcal$ a single block, every component and the prediction have minimax risk
of order at least $n^{-2s/(2s+1)}$, and with $\Kcal$ all the blocks, so that
$|\Kcal| = pq$, the bound is proportional to $pq$.
\end{theorem}

The bound is not asymptotic, and it holds for every pair $(s,\pi)$, in
particular for every pair that Assumption~5 allows. It has the form of the
lower bound of \citet[Theorem~5]{Donoho-Johnstone-1998} for Besov bodies in
the Gaussian sequence model, with radius $C_g$ and noise level
$\sigma(nB_X^{2}C_U)^{-1/2}$ per coordinate. Since Gaussian errors satisfy
Assumption~1, it holds a fortiori over the models of Assumptions~1 to~5, and
the form in probability is the one to compare with the upper bounds, which are
stated in $O_p$ and are uniform over the class (proof of Theorem~1 above).
Hence, when $\pi \ge 2$, $n^{-2s/(2s+1)}$ is the minimax rate in the
prediction norm and for the components, for every $q$ and with covariates
that depend on the modulators, and the estimator of Theorem~1 attains it
uniformly over the class; when $\pi < 2$ and $\seff > s/(2s+1)$ it attains it
up to the factor $(\log n)^{(2/\pi-1)/(2s+1)}$, which the lower bound does not
carry
% D58 (pergunta 52(f)): a condicao de Cai (1999), que saiu com a Remark S6.1.
(the same factor as in the rate of block thresholding in the sequence model,
\citealp[Theorem~4]{Cai-1999}, which holds for $s \ge 1/\pi$, a condition
strictly stronger than $\seff > s/(2s+1)$), and the coordinatewise lasso, at the rate $(\log n/n)^{2s/(2s+1)}$ of
Theorem~\ref{thm:lasso-main}, stays above it by the factor
$(\log n)^{2s/(2s+1)}$ for every $\pi$. When $q = 1$ and $\bX$ is independent
of $U_1$, part (i) with $\Kcal$ all the blocks is, up to the constants, the
nonparametric term of the lower bound of
\citet[Theorem~1, eqs.~3.6 and~3.7]{Klopp-Pensky-2015} with a fixed number of
covariates.

The proof applies Assouad's lemma to a hypercube of wavelets at a single
level. Fix $\Kcal$, a level $\lev \ge 0$ and a height $\zeta > 0$; let
\[
  \mathcal{C} = \bigl\{(\cv,\md,\tsl) : (\cv,\md)\in\Kcal,\ 0 \le \tsl < 2^{\lev}\bigr\},
  \qquad |\mathcal{C}| = |\Kcal|\,2^{\lev},
\]
and $\Omega = \{0,1\}^{\mathcal{C}}$. For $\omega \in \Omega$ the components
are $\gcomp{\cv}{\md}^{\omega} =
\zeta\sum_{\tsl}\omega_{\cv\md,\tsl}\wav{\lev}{\tsl}$ for
$(\cv,\md) \in \Kcal$ and zero otherwise, the levels are zero, and
\[
  f_{\omega}(x,u) \;=\; \zeta\,z_{\mathcal{C}}(x,u)\T\omega,
  \qquad
  z_{\mathcal{C}}(x,u) = \bigl(x_{\cv}\wav{\lev}{\tsl}(u_{\md})\bigr)_{(\cv,\md,\tsl)\in\mathcal{C}} ,
\]
whose entries are columns of the design at resolution $J = \lev + 1$. Write
$\Gram_{\mathcal{C}\mathcal{C}}$ and $\Gramhat_{\mathcal{C}\mathcal{C}}$ for
the principal submatrices of $\Gram$ and $\Gramhat$ at that resolution on the
columns in $\mathcal{C}$, $\bZ_{\mathcal{C}}$ for those columns of $\bZ$,
$d_{H}(\omega,\omega') = \sum_{(\cv,\md,\tsl)\in\mathcal{C}}
\mathbf{1}\{\omega_{\cv\md,\tsl} \ne \omega'_{\cv\md,\tsl}\}$ for the Hamming
distance, and $\Prob_{\omega}$ and $\E_{\omega}$ for $\Prob_{f_{\omega}}$ and
$\E_{f_{\omega}}$. The form of Assouad's lemma that we use is
\citet[Theorem~2.12(iv)]{Tsybakov-2009} with the Kullback--Leibler divergence
bounded by one: if $\operatorname{KL}(\Prob_{\omega},\Prob_{\omega'}) \le 1$
whenever $d_{H}(\omega,\omega') = 1$, then
\begin{equation}
  \label{eq:assouad}
  \inf_{\widehat\omega}\max_{\omega\in\Omega}\E_{\omega}d_{H}(\widehat\omega,\omega)
  \;\ge\; \frac{|\mathcal{C}|}{2}\max\Bigl(\frac{e^{-1}}{2},\,1-2^{-1/2}\Bigr)
  \;=\; \cA|\mathcal{C}| ,
\end{equation}
the infimum being over all the estimators with values in $\Omega$.

\begin{lemma}[A hypercube of wavelets]
\label{lem:hypercube}
In the setting of Theorem~\ref{thm:lower}, for every nonempty $\Kcal$, every
$\lev \ge 0$ and every $\zeta > 0$:
\begin{enumerate}
\item[(i)] $f_{\omega} \in \mathcal{F}$ for every $\omega \in \Omega$ if and
only if $\zeta2^{\lev(s+1/2)} \le C_g$, for every $s > 0$ and every
$1 \le \pi \le \infty$;
\item[(ii)] for every pair,
$\operatorname{KL}(\Prob_{\omega},\Prob_{\omega'}) =
\{n\zeta^{2}/(2\sigma^{2})\}(\omega-\omega')\T\Gram_{\mathcal{C}\mathcal{C}}(\omega-\omega')$,
and, if $\omega$ and $\omega'$ differ only in the coordinate $(\cv,\md,\tsl)$,
\[
  \operatorname{KL}(\Prob_{\omega},\Prob_{\omega'})
  \;=\; \frac{n\zeta^{2}}{2\sigma^{2}}\,\E\bigl\{X_{\cv}^{2}\wav{\lev}{\tsl}(U_{\md})^{2}\bigr\}
  \;\le\; \frac{n\zeta^{2}}{2\sigma^{2}}\,\min(B_X^{2},\kappa_2)\,C_U ;
\]
\item[(iii)] $\lmin(\Gram_{\mathcal{C}\mathcal{C}}) \ge \gamma$, so that
\begin{gather*}
  \norm{f_{\omega}-f_{\omega'}}_{L_2(P)}^{2} \ge \gamma\zeta^{2}d_{H}(\omega,\omega'),
  \\
  \sum_{(\cv,\md)\in\Kcal}\lVert\gcomp{\cv}{\md}^{\omega} -
  \gcomp{\cv}{\md}^{\omega'}\rVert_{L_2[0,1]}^{2} = \zeta^{2}d_{H}(\omega,\omega') ;
\end{gather*}
with a single block $(\cv,\md)$, the first statement uses only
$\E(X_{\cv}^{2}\mid\bU) \ge \kappa_1$ and the lower bound $c_U$ on the
density of $U_{\md}$;
\item[(iv)] for each of the three losses and every estimator there is a
measurable function $\widehat\omega$ of the sample with values in $\Omega$
such that, for every $\omega \in \Omega$,
\begin{enumerate}
\item[(a)] $L_{\Kcal}(\widehat{g},f_{\omega}) \ge (\zeta^{2}/4)\,d_{H}(\widehat\omega,\omega)$,
\item[(b)] $\norm{\widehat{f}-f_{\omega}}_{L_2(P)}^{2} \ge (\gamma\zeta^{2}/4)\,d_{H}(\widehat\omega,\omega)$,
\item[(c)] $\normn{\widehat{f}-f_{\omega}}^{2} \ge (\gamma\zeta^{2}/8)\,
d_{H}(\widehat\omega,\omega)\,\mathbf{1}_{\mathcal{E}_{\mathcal{C}}}$, where
$\mathcal{E}_{\mathcal{C}} = \{\lmin(\Gramhat_{\mathcal{C}\mathcal{C}}) \ge \gamma/2\}$
depends only on the $(\bX_i,\bU_i)$ and has, with $J = \lev+1$ and $D$ and
$R_J$ as in Theorem~\ref{thm:lower}(iii),
\[
  \Prob(\mathcal{E}_{\mathcal{C}}^{c}) \;\le\;
  2D\exp\Bigl\{-\frac{n\gamma^{2}}{8R_{J}(\kappa_2C_U + \gamma/3)}\Bigr\} .
\]
\end{enumerate}
\end{enumerate}
\end{lemma}

The Kullback--Leibler divergence uses only upper bounds on the design, and
through the conditional second moment $\E(X_{\cv}^{2}\mid\bU)$, which is why
the covariates may depend on the modulators at no cost. The separation of the
components uses no property of the design; that of the prediction uses the
smallest eigenvalue, and it is there that the cross terms between modulators
and between covariates enter, through Proposition~2(i).

\subsection*{Proof of Lemma~\ref{lem:hypercube}}

\emph{Part (i).} By the orthonormality of $\{\mathbf{1}\}\cup\{\wav{\lev}{\tsl}\}$
in $L_2[0,1]$ (Section~S2), the coefficients of $\gcomp{\cv}{\md}^{\omega}$
are $\theta_{\cv\md,\lev'\tsl} = \zeta\omega_{\cv\md,\tsl}\mathbf{1}\{\lev' =
\lev\}$, and $\int_0^1\gcomp{\cv}{\md}^{\omega} = 0$ because
$\int_0^1\wav{\lev}{\tsl} = 0$. Only the level $\lev$ is nonzero, so the bound
of Assumption~5 reduces to that level, whatever the index $r$. For
$\pi < \infty$, $\norm{\theta_{\cv\md,\lev\cdot}}_{\pi} =
\zeta(\sum_{\tsl}\omega_{\cv\md,\tsl})^{1/\pi} \le \zeta2^{\lev/\pi}$, with
equality at $\omega \equiv 1$; for $\pi = \infty$ it is $\zeta$ if
$\omega_{\cv\md,\cdot} \ne 0$. In both cases the maximum over $\Omega$ of
$2^{\lev(s+1/2-1/\pi)}\norm{\theta_{\cv\md,\lev\cdot}}_{\pi}$ is
$\zeta2^{\lev(s+1/2)}$. The components outside $\Kcal$ and the levels are
zero.

\emph{Part (ii).} Under $\Prob_{\omega}$ the pairs $(\bX_i,\bU_i)$ have the
law $P$, the same for every $\omega$, and, given them, the $Y_i$ are
independent with $Y_i \sim N\{f_{\omega}(\bX_i,\bU_i),\sigma^{2}\}$. The
divergence between product laws is the sum of the divergences, and that
between two laws with the same marginal is the mean of the divergence between
the conditionals, so that
$\operatorname{KL}(\Prob_{\omega},\Prob_{\omega'}) =
n\E\{(f_{\omega}-f_{\omega'})^{2}(\bX,\bU)\}/(2\sigma^{2})$, and
$f_{\omega}-f_{\omega'} = \zeta z_{\mathcal{C}}\T(\omega-\omega')$ gives the
quadratic form. For a neighbor, $\E\{(f_{\omega}-f_{\omega'})^{2}\} =
\zeta^{2}\E\{\E(X_{\cv}^{2}\mid\bU)\wav{\lev}{\tsl}(U_{\md})^{2}\}$, with
$\E(X_{\cv}^{2}\mid\bU) \le B_X^{2}$ and $\E(X_{\cv}^{2}\mid\bU) =
e_{\cv}\T\E(\bX\bX\T\mid\bU)e_{\cv} \le \kappa_2$ by Assumption~2, and
$\E\{\wav{\lev}{\tsl}(U_{\md})^{2}\} \le C_U\int_0^1\wav{\lev}{\tsl}^{2} = C_U$,
because the density of $U_{\md}$ inherits the bound $C_U$ of Assumption~3.

\emph{Part (iii).} For $v \in \R^{\mathcal{C}}$ and $\tilde v$ its completion
with zeros, $v\T\Gram_{\mathcal{C}\mathcal{C}}v = \tilde v\T\Gram\tilde v \ge
\lmin(\Gram)\norm{v}_{2}^{2} \ge \gamma\norm{v}_{2}^{2}$, by
Proposition~2(i). Since $\omega-\omega' \in \{-1,0,1\}^{\mathcal{C}}$,
$\norm{\omega-\omega'}_{2}^{2} = d_{H}(\omega,\omega')$, and
$\norm{f_{\omega}-f_{\omega'}}_{L_2(P)}^{2} =
\zeta^{2}(\omega-\omega')\T\Gram_{\mathcal{C}\mathcal{C}}(\omega-\omega') \ge
\gamma\zeta^{2}d_{H}(\omega,\omega')$. The identity for the components is
Parseval's. With a single block,
\begin{align*}
  v\T\Gram_{\mathcal{C}\mathcal{C}}v &=
  \E\Bigl[\E(X_{\cv}^{2}\mid\bU)\Bigl\{\sum_{\tsl}v_{\tsl}\wav{\lev}{\tsl}(U_{\md})\Bigr\}^{2}\Bigr]
  \\
  &\ge \kappa_1c_U\int_0^1\Bigl(\sum_{\tsl}v_{\tsl}\wav{\lev}{\tsl}\Bigr)^{2} =
  \gamma\norm{v}_{2}^{2} .
\end{align*}

\emph{Part (iv).} Let $\varpi(y) = \mathbf{1}\{y \ge \zeta/2\}$ for real $y$.
For $\omega_{\cv\md,\tsl} \in \{0,1\}$,
\begin{equation}
  \label{eq:hypercube-threshold}
  |y - \zeta\omega_{\cv\md,\tsl}| \;\ge\; \frac{\zeta}{2}\,
  |\varpi(y) - \omega_{\cv\md,\tsl}| ,
\end{equation}
because if $\varpi(y) \ne \omega_{\cv\md,\tsl}$ then $y$ and
$\zeta\omega_{\cv\md,\tsl}$ lie on opposite sides of $\zeta/2$. Applied
coordinatewise to a vector $a \in \R^{\mathcal{C}}$, it gives
\[
  \norm{a-\zeta\omega}_{2}^{2} \;\ge\; \frac{\zeta^{2}}{4}\,d_{H}(\varpi(a),\omega) .
\]
For
(a), take $a_{\cv\md,\tsl} = \inner{\widehat{g}_{\cv\md}}{\wav{\lev}{\tsl}}$
and $\widehat\omega = \varpi(a)$: by Bessel's inequality,
$\lVert\widehat{g}_{\cv\md}-\gcomp{\cv}{\md}^{\omega}\rVert_{L_2[0,1]}^{2} \ge
\sum_{\tsl}(a_{\cv\md,\tsl}-\zeta\omega_{\cv\md,\tsl})^{2}$ for each
$(\cv,\md) \in \Kcal$, and it remains to sum over $\Kcal$. For (b), let
$\Pi_{\mathcal{C}}$ be the orthogonal projection of $L_2(P)$ onto the span of
the entries of $z_{\mathcal{C}}$, of dimension $|\mathcal{C}|$ because
$\Gram_{\mathcal{C}\mathcal{C}}$ is positive definite by (iii):
$\Pi_{\mathcal{C}}\widehat{f} = z_{\mathcal{C}}\T a$ with
$a = \Gram_{\mathcal{C}\mathcal{C}}^{-1}\E\{z_{\mathcal{C}}(\bX,\bU)\widehat{f}(\bX,\bU)\}$,
the expectation being over a new $(\bX,\bU)$ with $\widehat{f}$ fixed, so that
$a$ is a measurable function of the sample, $P$ being fixed. Since $f_{\omega}$
lies in the span, Pythagoras' theorem gives
$\norm{\widehat{f}-f_{\omega}}_{L_2(P)}^{2} \ge
\norm{\Pi_{\mathcal{C}}\widehat{f}-f_{\omega}}_{L_2(P)}^{2} =
(a-\zeta\omega)\T\Gram_{\mathcal{C}\mathcal{C}}(a-\zeta\omega) \ge
\gamma\norm{a-\zeta\omega}_{2}^{2}$; take $\widehat\omega = \varpi(a)$. For
(c), write $\widehat{\mathbf{F}} = (\widehat{f}(\bX_i,\bU_i))_{i\le n}$. On
$\mathcal{E}_{\mathcal{C}}$, $\bZ_{\mathcal{C}}$ has rank $|\mathcal{C}|$; take
$a = (\bZ_{\mathcal{C}}\T\bZ_{\mathcal{C}})^{-1}\bZ_{\mathcal{C}}\T\widehat{\mathbf{F}}$
and $\widehat\omega = \varpi(a)$ there, and $\widehat\omega = 0$ off it. On
$\mathcal{E}_{\mathcal{C}}$, since $\zeta\bZ_{\mathcal{C}}\omega$ is the vector
of the values of $f_{\omega}$ and lies in the column space of
$\bZ_{\mathcal{C}}$,
$n\normn{\widehat{f}-f_{\omega}}^{2} =
\norm{\widehat{\mathbf{F}}-\zeta\bZ_{\mathcal{C}}\omega}_{2}^{2} \ge
\norm{\bZ_{\mathcal{C}}(a-\zeta\omega)}_{2}^{2} =
n(a-\zeta\omega)\T\Gramhat_{\mathcal{C}\mathcal{C}}(a-\zeta\omega) \ge
n(\gamma/2)\norm{a-\zeta\omega}_{2}^{2}$; off it the right side of (c) is
zero. Finally $\Gramhat_{\mathcal{C}\mathcal{C}}$ is a principal submatrix of
$\Gramhat$ at resolution $J = \lev+1$, so that
$\lmin(\Gramhat_{\mathcal{C}\mathcal{C}}) \ge \lmin(\Gramhat)$, and
$\mathcal{E}_{\mathcal{C}}$ contains the event of Proposition~2(ii), which
gives the bound on its complement. \qed

\subsection*{Proof of Theorem~\ref{thm:lower}}

\emph{The hypercube.} Take $\zeta = \zeta_n$ and $\lev = \lev_n$ of
\eqref{eq:lower-choices}. Since $2^{\lev_n(2s+1)} \le C_g^{2}/\zeta_n^{2}$,
$\zeta_n2^{\lev_n(s+1/2)} \le C_g$, and Lemma~\ref{lem:hypercube}(i) gives
$f_{\omega} \in \mathcal{F}$ for every $\omega$. By
Lemma~\ref{lem:hypercube}(ii), the divergence between neighbors is at most one,
because $n\zeta_n^{2}B_X^{2}C_U = 2\sigma^{2}$, and
\eqref{eq:assouad} gives, for every $\widehat\omega$ with values in $\Omega$,
\begin{equation}
  \label{eq:assouad-n}
  \max_{\omega\in\Omega}\E_{\omega}d_{H}(\widehat\omega,\omega)
  \;\ge\; \cA|\mathcal{C}|,
  \qquad |\mathcal{C}| = |\Kcal|\,2^{\lev_n} .
\end{equation}

\emph{The scale.} By the maximality of $\lev_n$,
$2^{(\lev_n+1)(2s+1)} > C_g^{2}/\zeta_n^{2}$, so that
$2^{\lev_n} > (C_g/\zeta_n)^{2/(2s+1)}/2$ and
\begin{align*}
  2^{\lev_n}\zeta_n^{2} \;&>\; \frac12\,C_g^{2/(2s+1)}(\zeta_n^{2})^{2s/(2s+1)}
  = \frac12\,C_g^{2/(2s+1)}\Bigl(\frac{2\sigma^{2}}{nB_X^{2}C_U}\Bigr)^{2s/(2s+1)}
  \\
  &= 2^{-1/(2s+1)}\rlow^{2} \;\ge\; \frac12\,\rlow^{2},
\end{align*}
that is, $\zeta_n^{2}|\mathcal{C}| \ge |\Kcal|\rlow^{2}/2$.

\emph{Expectation.} The supremum over $\mathcal{F}$ is at least the maximum
over the hypercube. With the $\widehat\omega$ of
Lemma~\ref{lem:hypercube}(iv)(a) and \eqref{eq:assouad-n},
\[
  \max_{\omega}\E_{\omega}L_{\Kcal}(\widehat{g},f_{\omega}) \;\ge\;
  \frac{\zeta_n^{2}}{4}\,\cA|\mathcal{C}| \;\ge\; \frac{\cA}{8}\,|\Kcal|\,\rlow^{2} ;
\]
with that of (iv)(b), the same with the factor $\gamma$. For (iii), since
$d_{H} \le |\mathcal{C}|$ and the law of the $(\bX_i,\bU_i)$ does not depend
on $\omega$, $\E_{\omega}\{d_{H}(\widehat\omega,\omega)\mathbf{1}_{\mathcal{E}_{\mathcal{C}}}\}
\ge \E_{\omega}d_{H}(\widehat\omega,\omega) -
|\mathcal{C}|\Prob(\mathcal{E}_{\mathcal{C}}^{c})$, and for $n \ge n_{1}$
\[
  \max_{\omega}\E_{\omega}\normn{\widehat{f}-f_{\omega}}^{2}
  \;\ge\; \frac{\gamma\zeta_n^{2}}{8}\Bigl(\cA|\mathcal{C}| -
  \frac{\cA}{4}|\mathcal{C}|\Bigr)
  = \frac{3\cA}{32}\,\gamma\zeta_n^{2}|\mathcal{C}|
  \;\ge\; \frac{\cA}{32}\,\gamma|\Kcal|\,\rlow^{2} .
\]

\emph{Probability.} Let $\omega^{*}$ attain the maximum in
\eqref{eq:assouad-n}, which exists because $\Omega$ is finite. Since
$0 \le d_{H} \le |\mathcal{C}|$,
$\cA|\mathcal{C}| \le \E_{\omega^{*}}d_{H} \le \cA|\mathcal{C}|/2 +
|\mathcal{C}|\Prob_{\omega^{*}}(d_{H} \ge \cA|\mathcal{C}|/2)$, so that
$\Prob_{\omega^{*}}(d_{H} \ge \cA|\mathcal{C}|/2) \ge \cA/2$. On that event,
(iv)(a) gives $L_{\Kcal} \ge (\zeta_n^{2}/4)\cA|\mathcal{C}|/2 \ge
(\cA/16)|\Kcal|\rlow^{2}$, and (iv)(b) the same with the factor $\gamma$; on
its intersection with $\mathcal{E}_{\mathcal{C}}$, of probability at least
$\cA/2 - \cA/4$ when $n \ge n_{1}$, (iv)(c) gives
$\normn{\widehat{f}-f_{\omega^{*}}}^{2} \ge
(\gamma\zeta_n^{2}/8)\cA|\mathcal{C}|/2 \ge (\cA/32)\gamma|\Kcal|\rlow^{2}$.
The estimator being arbitrary, the infima follow.

\emph{The integer $n_{1}$ is finite.} Since $2^{\lev_n+1} \le
2(C_g/\zeta_n)^{2/(2s+1)}$, of order $n^{1/(2s+1)}$, $D$ and $R_J$ with
$J = \lev_n + 1$ are of order $n^{1/(2s+1)}$, and the bound of
Lemma~\ref{lem:hypercube}(iv)(c) is of order
$n^{1/(2s+1)}\exp(-cn^{2s/(2s+1)})$ for a constant $c > 0$, which tends to
zero. The two particular cases are $\Kcal$ a single block and $\Kcal$ all the
blocks. \qed

% E5f (E1.15; D54): derivations/08-blocos.tex, Corolario 15 (numeracao global;
% cor:truncada-blocos) e a sua prova, dentro do grupo azul de k = 4.
\subsection*{The slow rate through a truncated comparator}

Proposition~3 reads Theorem~2(i) at the comparator $\betastar$ only, and pays
$\pen w_{G}\norm{\thetastar_{G}}_{2}$ for every chunk, including the chunks
whose energy would cost less as bias. Theorem~2 holds for every comparator on
the same event, and the comparator truncated chunk by chunk,
\begin{equation}
  \label{eq:truncated-blocks}
  \thbar_{G} \;=\; \thetastar_{G}\,\mathbf{1}\bigl\{\norm{\thetastar_{G}}_{2}
  > \pen w_{G}\bigr\},
  \qquad G \in \Gcal,
\end{equation}
charges the penalty only for the chunks above their own level and sends the
others to the bias. What remains is
\begin{equation}
  \label{eq:RG}
  \RGw(\btheta;\pen) = \sum_{G\in\Gcal}
  \min\bigl(\pen w_{G}\norm{\btheta_{G}}_{2},\,\norm{\btheta_{G}}_{2}^{2}\bigr),
  \qquad
  \RGone(\btheta;t) = \sum_{G\in\Gcal}
  \min\bigl(t\norm{\btheta_{G}}_{2},\,\norm{\btheta_{G}}_{2}^{2}\bigr),
\end{equation}
the counterparts in the group norm of the ideal risk $\Rid$, with
$\RGone(\btheta;t) \le t\nGone{\btheta}$.

\begin{proposition}[Slow rate through the truncated comparator]
\label{prop:truncated}
Let the assumptions of Proposition~3 hold; in the statements with $\lamGp$
in the estimator, the conditions on $J_n$ are replaced by $2^{J_n} \le n^{2}$,
as in Proposition~\ref{prop:slow-ceiling}.
\begin{enumerate}
\item[(i)] On the event $\Tw$, with $\thbar$ of \eqref{eq:truncated-blocks},
\[
  \normn{\widehat{f}-f}^{2}
  \;\le\; 18\pen\nGw{\thbar} + 48\normn{\Bmat(\thetastar-\thbar)}^{2}
  + 48\normn{\bb}^{2} + 3\normn{\PA\bveps}^{2} .
\]
If $\pen$ and $w$ are deterministic,
$\E\normn{\Bmat(\thetastar-\thbar)}^{2} \le
\kappa_2C_U\norm{\thetastar-\thbar}_{2}^{2}$,
\[
  18\pen\nGw{\thbar} + 48\kappa_2C_U\norm{\thetastar-\thbar}_{2}^{2}
  \;\le\; \max(18,48\kappa_2C_U)\,\RGw(\thetastar;\pen),
\]
and $\RGw(\thetastar;\lamG) \le \bar\varrho\,\RGone(\thetastar;\lamGone)$.
Consequently, with deterministic weights,
\[
  \normn{\widehat{f}-f}^{2} \;=\; O_p\bigl\{\RGone(\thetastar;\lamGone)
  + \mathcal{B}_{J_n}^{2} + p/n\bigr\},
\]
and, with $\lamGp$ in the estimator, the same holds with
$(\mu^{\Gcal}_{J_n})^{1/2}\lamGone$ in place of $\lamGone$, for every sequence
$J_n$.
\item[(ii)] For every $J$, every $b \ge 2$ and every $t > 0$, with
$\tau = (s+1/2)^{-1}$: if $s < 1/2$,
\begin{equation}
  \label{eq:count-truncated}
  \RGone(\thetastar;t) \;\le\; pq\Bigl\{\frac{C_g\,t}{1-2^{-\seff}}
  + C_g^{\tau}\,t^{4s/(2s+1)}\,b^{\{(2/\pi-1)_{+}-2s\}/(2s+1)}
  \Bigl(\frac{1}{1-2^{s-1/2}} + \frac{1}{1-2^{-\min(\pi,2)\seff}}\Bigr)\Bigr\};
\end{equation}
if $s \ge 1/2$, nothing is gained in order: $\RGone(\thetastar;t) \le
t\nGone{\thetastar}$, and for $t \le C_g$ the class of Assumption~5 contains a
$\thetastar$ with $\RGone(\thetastar;t) \ge tC_g$ if $s > 1/2$, and, if
$s = 1/2$, one in which every fine level $\lev < J$ with
$2^{\lev} \le C_gb^{1-\vartheta_{\pi}}/t$ contributes at least
$tC_gb^{-\vartheta_{\pi}}/2$ to $\RGone(\thetastar;t)$, the order of
Lemma~\ref{lem:count}.
\item[(iii)] If $s < 1/2$ and $\seff > s/(2s+1)$, then with
$J_n = \lceil c\log_2 n\rceil$ and $c$ in the window (3.2) of Theorem~1,
\[
  \normn{\widehat{f}-f}^{2}
  \;=\; O_p\bigl\{n^{-2s/(2s+1)}(\log n)^{(2/\pi-1)_{+}/(2s+1)}\bigr\},
\]
the prediction rate of Theorem~1.
\item[(iv)] If $s < 1/2$ and $\seff \le s/(2s+1)$, which forces $\pi < 2$,
then with $\lamGp$ in the estimator and the choice
$2^{J_n} \asymp \{n^{2s}\bn^{-(2s+2/\pi-1)/2}\}^{1/\{s+\seff(2s+1)\}}$, which
exceeds $n/\bn$,
\[
  \normn{\widehat{f}-f}^{2}
  \;=\; O_p\bigl\{n^{-e}(\log n)^{e(2s+2/\pi-1)/(4s)}\bigr\},
  \qquad e = \frac{4s\seff}{s+\seff(2s+1)} ;
\]
at $\seff = s/(2s+1)$ this is
$n^{-2s/(2s+1)}(\log n)^{(2s+2/\pi-1)/\{2(2s+1)\}}$.
\end{enumerate}
\end{proposition}

Proposition~\ref{prop:truncated} is a statement on prediction. Like
Proposition~3, it uses no lower bound on $\Gramhat$: besides the upper bound
on $\lmaxx(\Gramhat)$ that makes the level admissible, it uses only
$\lmaxx(\Gram) \le \kappa_2C_U$, in the new bias, and the components still
need the design condition (Corollary~1). When $s < 1/2$ and $\pi < 2$ it
improves on Proposition~3(ii) and on Proposition~\ref{prop:slow-ceiling}(iii)(b)
in every pair: (iii) gives the rate of Theorem~1 in the window of that
theorem, and (iv) replaces Proposition~\ref{prop:slow-ceiling}(iii)(b) and the
part of Proposition~3(ii) with $s/2 < \seff \le s/(2s+1)$. When $\pi \ge 2$
the rate is that of Proposition~3(ii), $n^{-2s/(2s+1)}$, now for every
$c \in [1/(2s+1),1)$. When $s \ge 1/2$ nothing changes: a coarse chunk of
norm $C_g$, which the class allows, costs $\pen C_g$ in the penalty or $C_g^{2}$ in the bias, and
the slow route does not go below $\lamGone \asymp (\log n/n)^{1/2}$;
Proposition~3(i) and Proposition~\ref{prop:slow-ceiling}(iii)(a) and (iv)
remain the statements with no design condition in those pairs. Against the
lasso read at its own truncated comparator
(Proposition~\ref{prop:lasso-truncated}), the chunks remove the factor
$(\log n)^{2\seff/(2s+1)}$ in (iii), the gain of Theorem~1 over
Theorem~\ref{thm:lasso-main}, and $(\log n)^{e\seff/(2s)}$ in (iv). The
proposition is stated for the estimator (2.5) and not for the one that the
software computes (Remark~\ref{rem:white}); for the latter the comparator
would truncate at $\pen w_{G}\Lambda^{1/2}$ on the event $q_{\max} \le \Lambda$,
and we do not pursue it.

\subsection*{Proof of Proposition~\ref{prop:truncated}}

\emph{Part (i).} Theorem~2(i) and (iii) with $\bar\bbeta = (\bc,\thbar)$ give,
as in the proof of \eqref{eq:slow-nonasymptotic},
$\normn{\widehat{f}-f}^{2} \le 18\pen\nGw{\thbar} + 24\normn{\bbar}^{2} +
3\normn{\PA\bveps}^{2}$, and $\bbar = f - \Amat\bc - \Bmat\thbar = \bb +
\Bmat(\thetastar-\thbar)$ gives $\normn{\bbar}^{2} \le 2\normn{\bb}^{2} +
2\normn{\Bmat(\thetastar-\thbar)}^{2}$. If $\pen$ and $w$ are deterministic,
so is $u = \thetastar-\thbar$, and $\E\normn{\Bmat u}^{2} =
u\T\E(\Bmat\T\Bmat/n)u \le \lmaxx(\Gram)\norm{u}_{2}^{2} \le
\kappa_2C_U\norm{u}_{2}^{2}$, because $\E(\Bmat\T\Bmat/n)$ is a principal
submatrix of $\Gram$, by the upper half of Proposition~2(i). On the chunks with
$\norm{\thetastar_{G}}_{2} > \pen w_{G}$ the minimum in \eqref{eq:RG} is
$\pen w_{G}\norm{\thetastar_{G}}_{2}$, and on the others it is
$\norm{\thetastar_{G}}_{2}^{2}$; summing gives the bound by
$\max(18,48\kappa_2C_U)\RGw(\thetastar;\pen)$. By
Lemma~\ref{lem:calib-blocks}(iii), $\lamG w_{G} \le \varrho\lamGone$, and
$\min(\varrho x,y) \le \varrho\min(x,y)$ for $\varrho \ge 1$ gives
$\RGw(\thetastar;\lamG) \le \varrho\,\RGone(\thetastar;\lamGone)$; with
$\lamGp$, the same with $(\mu^{\Gcal}_{J})^{1/2}\lamGone$. The form in $O_p$
is obtained as in the proof of Proposition~3, on the intersection of
$\Acal^{+}$ below the ceiling, or of $\Acal^{++}$ of the proof of
Proposition~\ref{prop:slow-ceiling} for every $J_n$, with $\Ecal$ and with
the events given by Markov's inequality for $\normn{\bb}^{2}$,
$\normn{\Bmat(\thetastar-\thbar)}^{2}$ and $\normn{\PA\bveps}^{2}$.

\emph{Part (ii), $s < 1/2$.} Fix a block. The coarse chunk contributes at most
$t\norm{\theta_{\mathrm{coarse}}}_{2} \le tC_g(1-2^{-\seff})^{-1}$, by the
proof of Lemma~\ref{lem:count}. At a fine level $\lev$, with $\Gcal_{\lev}$
its chunks and $v_{\lev} = \sum_{G\in\Gcal_{\lev}}\min(t\norm{\theta_{G}}_{2},
\norm{\theta_{G}}_{2}^{2})$, there are two bounds. The first is in the group
norm: $v_{\lev} \le t\sum_{G}\norm{\theta_{G}}_{2} \le
tC_gb^{-\vartheta_{\pi}}2^{\lev(1/2-s)} =: A_{\lev}$, by the proof of
Lemma~\ref{lem:count}. The second is the tail. If $\pi \le 2$,
$\min(ty,y^{2}) \le t^{2-\pi}y^{\pi}$ for $y \ge 0$ (if $y \le t$,
$y^{2} = y^{\pi}y^{2-\pi} \le y^{\pi}t^{2-\pi}$; if $y > t$,
$ty = t^{2-\pi}t^{\pi-1}y \le t^{2-\pi}y^{\pi}$, because $\pi \ge 1$), and
$\norm{\theta_{G}}_{2} \le \norm{\theta_{G}}_{\pi}$ gives
$v_{\lev} \le t^{2-\pi}\sum_{G}\norm{\theta_{G}}_{\pi}^{\pi} =
t^{2-\pi}\norm{\theta_{\lev}}_{\pi}^{\pi} \le
t^{2-\pi}C_g^{\pi}2^{-\lev\pi\seff}$, since $\pi(s+1/2-1/\pi) = \pi\seff$; if
$\pi \ge 2$, $v_{\lev} \le \norm{\theta_{\lev}}_{2}^{2} \le C_g^{2}2^{-2\lev s}$.
Call this second bound $B_{\lev}$. With
\[
  2^{x} = \bigl\{(C_g/t)\,b^{1-\vartheta_{\pi}}\bigr\}^{\tau},
  \qquad
  h = C_g^{\tau}\,t^{2-\tau}\,b^{\tau(1-\vartheta_{\pi})-1},
\]
one has $A_{\lev} = h2^{(\lev-x)(1/2-s)}$ and
$B_{\lev} = h2^{-(\lev-x)\min(\pi,2)\seff}$. Indeed, at $\lev = x$,
$A_{x} = C_g^{\tau}t^{2-\tau}b^{-\vartheta_{\pi}+(1-\vartheta_{\pi})(\tau-1)}$,
because $\tau(1/2-s) = \tau-1$, and the exponent of $b$ is
$\tau(1-\vartheta_{\pi})-1$; if $\pi \le 2$, $1-\vartheta_{\pi} = 1/\pi$ and
$B_{x} = t^{2-\pi}C_g^{\pi}\{(C_g/t)b^{1/\pi}\}^{-(\pi-\tau)} =
C_g^{\tau}t^{2-\tau}b^{\tau/\pi-1}$, because $\tau\pi\seff = \pi-\tau$; if
$\pi \ge 2$, $1-\vartheta_{\pi} = 1/2$ and
$B_{x} = C_g^{2}\{(C_g/t)b^{1/2}\}^{-(2-\tau)} =
C_g^{\tau}t^{2-\tau}b^{\tau/2-1}$, because $2s\tau = 2-\tau$. For every $\pi$,
$\tau(1-\vartheta_{\pi})-1 = \{(2/\pi-1)_{+}-2s\}/(2s+1)$ and
$2-\tau = 4s/(2s+1)$. Summing $A_{\lev}$ over the fine levels with
$\lev \le x$ and $B_{\lev}$ over those with $\lev > x$, two geometric series
of ratios $2^{-(1/2-s)}$ and $2^{-\min(\pi,2)\seff}$ that start from at most
$h$, gives at most $h\{(1-2^{s-1/2})^{-1} + (1-2^{-\min(\pi,2)\seff})^{-1}\}$
per block. Truncating at $\lev < J$ only removes terms, and summing over the
$pq$ blocks gives \eqref{eq:count-truncated}.

\emph{Part (ii), $s \ge 1/2$.} The first statement is
$\min(ty,y^{2}) \le ty$. If $s > 1/2$, the block with
$\theta^{*}_{\cv\md,00} = C_g$ and all its other coefficients zero satisfies
Assumption~5, and its coarse chunk, of norm $C_g$, contributes
$\min(tC_g,C_g^{2}) = tC_g$. If $s = 1/2$, take a fine level $\lev$ with
$2^{\lev} \le C_gb^{1-\vartheta_{\pi}}/t$, and let $k_{\lev}$ be the number of
its chunks, so that $2^{\lev}/(2b) \le k_{\lev} \le 2^{\lev}/b$. If
$\pi \le 2$, put a spike of height
$y_{\lev} = C_g2^{-\lev(1-1/\pi)}k_{\lev}^{-1/\pi}$ in each chunk, which
saturates Assumption~5 at that level; since $k_{\lev} \le 2^{\lev}/b$,
$y_{\lev} \ge C_g2^{-\lev}b^{1/\pi} \ge t$, and the level contributes
$k_{\lev}ty_{\lev} = tC_g2^{-\lev(1-1/\pi)}k_{\lev}^{1-1/\pi} \ge
tC_g(2b)^{-(1-1/\pi)}$. If $\pi \ge 2$, set all the coefficients of the level
equal to $C_g2^{-\lev}$, which again saturates Assumption~5: each chunk has
norm $|G|^{1/2}C_g2^{-\lev} \ge b^{1/2}C_g2^{-\lev} \ge t$, and the level
contributes $tC_g2^{-\lev}\sum_{G}|G|^{1/2} \ge tC_g2^{-\lev}k_{\lev}b^{1/2}
\ge tC_gb^{-1/2}/2$. In both cases the contribution is at least
$tC_gb^{-\vartheta_{\pi}}/2$.

\emph{Part (iii).} With $\pen = \lamG$ and
$(\lamGone)^{2} \le C_{\pen}\sigma^{2}\bn/n$ by \eqref{eq:order-lambda},
\eqref{eq:count-truncated} with $b = \bn$ gives
$\RGone(\thetastar;\lamGone) \lesssim (\bn/n)^{1/2} +
(\bn/n)^{2s/(2s+1)}\bn^{\{(2/\pi-1)_{+}-2s\}/(2s+1)} = (\bn/n)^{1/2} +
n^{-2s/(2s+1)}\bn^{(2/\pi-1)_{+}/(2s+1)}$, and the first term is of smaller
order because $2s/(2s+1) < 1/2$. The bias and the conditions on $J_n$ are
those of the proof of Theorem~1: $c \ge s/\{(2s+1)\seff\}$ gives
$\mathcal{B}_{J_n}^{2} \lesssim n^{-2s/(2s+1)}$, and $c < 1$ gives
$2^{J_n} \le 2n$ and $pq2^{J_n}\log(pq2^{J_n})/n \to 0$.

\emph{Part (iv).} Write $\eta = 2s/\{s+\seff(2s+1)\} \in [1,2)$ and
$\beta = (2s+2/\pi-1)/(2s+1)$, which lies in $(0,\,4s/(2s+1))$ because
$\pi < 2$ and $2\seff = 2s+1-2/\pi > 0$. The choice of the statement is
$2^{J_n} \asymp n^{\eta}\bn^{-\eta\beta(2s+1)/(4s)}$, which exceeds $n/\bn$
because $\beta(2s+1)/(4s) < 1$, and is at most $n^{2}$. By
\eqref{eq:lam-plus-blocks}, $t = (\mu^{\Gcal}_{J_n})^{1/2}\lamGone \asymp
2^{J_n/2}\bn/n$, and \eqref{eq:count-truncated} gives
\[
  t^{4s/(2s+1)}\bn^{\beta-4s/(2s+1)}
  \;\asymp\; 2^{J_n\cdot 2s/(2s+1)}\,n^{-4s/(2s+1)}\,\bn^{\beta},
\]
which, equated with the bias $2^{-2J_n\seff}$, gives
$2^{J_n\{2s/(2s+1)+2\seff\}} = n^{4s/(2s+1)}\bn^{-\beta}$; this is the choice
of the statement, because $\{2s/(2s+1)+2\seff\}(2s+1) = 2\{s+\seff(2s+1)\}$.
The common value is $2^{-2J_n\seff} =
\{\bn^{\beta}/n^{4s/(2s+1)}\}^{2\seff/\{2s/(2s+1)+2\seff\}} =
n^{-e}\bn^{e\beta(2s+1)/(4s)}$, and $\beta(2s+1)/(4s) = (2s+2/\pi-1)/(4s)$.
The term of the coarse chunk, of order $t$, is polynomially smaller, because
$t \to 0$ and $4s/(2s+1) < 1$. \qed
}

%%%%%%%%%%%%%%%%%%%%%%%%%%%%%%%%%%%%%%%%%%%%%%%%%%%%%%%%%%%%%%%%%%%%%%%%%%%%%%
% derivations/06-selecao-limiar.tex (Lema 13 e Corolario 8 na numeracao
% global; a Hipotese S e a Assumption 6 do artigo). Entrou em k = 2 (E5c).
{\color{colR1}
\section{Recovery of the block structure by thresholding}
\setcounter{equation}{0}

Throughout this section $\nu_{\cv\md}$, $\widehat{\nu}_{\cv\md}$, $\mathcal{S}$,
$\widehat{\mathcal{S}}(t)$, \textcolor{gray}{\sout{$\rho_n$}} $\rhoG$ and $\nu_{\min,n}$ are as in Section 3.6 of the paper,
$\nu_{\min} = \min_{(\cv,\md)\in \mathcal{S}}\nu_{\cv\md}$ for a fixed sample size, and
\[
  \Delta \;=\; \max_{\cv,\md}\,
  \bigl\lVert \widehat{g}_{\cv\md} - \gcomp{\cv}{\md} \bigr\rVert_{L_2[0,1]}
\]
is the largest error over the blocks. The first statement is deterministic
and knows nothing about the \textcolor{gray}{\sout{lasso}} estimator.

\begin{lemma}[Thresholding under a uniform block error]
\label{lem:sandwich}
Let $\{\gcomp{\cv}{\md}\}$ and $\{\widehat{g}_{\cv\md}\}$ be any two families
of functions in $L_2[0,1]$ indexed by the blocks, with $\nu_{\cv\md}$,
$\widehat{\nu}_{\cv\md}$, $\mathcal{S}$, $\widehat{\mathcal{S}}(t)$, $\Delta$ and $\nu_{\min}$ defined
from them as above, and let $\bar{\Delta} \ge \Delta$. Then, for every $t \ge 0$,
\begin{equation}
  \label{eq:sandwich}
  \bigl\{(\cv,\md) :\; \nu_{\cv\md} > t + \bar{\Delta}\bigr\}
  \;\subseteq\; \widehat{\mathcal{S}}(t) \;\subseteq\;
  \bigl\{(\cv,\md) :\; \nu_{\cv\md} > t - \bar{\Delta}\bigr\} .
\end{equation}
In particular: (i) if $t \ge \bar{\Delta}$, then $\widehat{\mathcal{S}}(t) \subseteq \mathcal{S}$, with no
use of $\nu_{\min}$; (ii) if $t < \nu_{\min} - \bar{\Delta}$, then $\mathcal{S} \subseteq \widehat{\mathcal{S}}(t)$;
and (iii) if $\bar{\Delta} \le t < \nu_{\min} - \bar{\Delta}$, then $\widehat{\mathcal{S}}(t) = \mathcal{S}$, and such a $t$
exists if and only if $\nu_{\min} > 2\bar{\Delta}$.
\end{lemma}

\begin{proof}
By the triangle inequality in $L_2[0,1]$, for every block,
$|\widehat{\nu}_{\cv\md} - \nu_{\cv\md}| \le
\lVert\widehat{g}_{\cv\md} - \gcomp{\cv}{\md}\rVert_{L_2[0,1]} \le \Delta \le \bar{\Delta}$.
If $\nu_{\cv\md} > t + \bar{\Delta}$, then
$\widehat{\nu}_{\cv\md} \ge \nu_{\cv\md} - \bar{\Delta} > t$, which is the left inclusion of
\eqref{eq:sandwich}; if $\widehat{\nu}_{\cv\md} > t$, then
$\nu_{\cv\md} \ge \widehat{\nu}_{\cv\md} - \bar{\Delta} > t - \bar{\Delta}$, which is the right one. For
(i), $t - \bar{\Delta} \ge 0$ and the right inclusion give $\widehat{\mathcal{S}}(t) \subseteq
\{\nu_{\cv\md} > 0\} = \mathcal{S}$. For (ii), every block of $\mathcal{S}$ has
$\nu_{\cv\md} \ge \nu_{\min} > t + \bar{\Delta}$, and the left inclusion applies. Part (iii) is
the conjunction of the two, and $[\bar{\Delta}, \nu_{\min} - \bar{\Delta})$ is nonempty exactly when
$\bar{\Delta} < \nu_{\min} - \bar{\Delta}$.
\end{proof}

\begin{remark}[The window is exact]
\label{rem:window}
The condition $\nu_{\min} > 2\bar{\Delta}$ is not an artifact of the proof. Take two blocks,
one active with $\nu_{11} = 2\bar{\Delta}$ and one null with $\nu_{21} = 0$, and fitted
norms $\widehat{\nu}_{11} = \widehat{\nu}_{21} = \bar{\Delta}$, which is compatible with
$\Delta \le \bar{\Delta}$: the two statistics tie and no threshold separates them. This
says that the argument does not extend below $\nu_{\min} = 2\bar{\Delta}$, not that no
procedure does, which would be a lower bound and is not proved here.
\end{remark}

% derivations/06-selecao-limiar.tex, Lema 16 (numeracao global; adendo de
% E1.12); D45.
\begin{lemma}[The risk of the thresholded estimator]
\label{lem:thrisk}
Let $\{\gcomp{\cv}{\md}\}$ and $\{\widehat{g}_{\cv\md}\}$ be any two families of
functions in $L_2[0,1]$ indexed by the blocks and $t \ge 0$, with
$\nu_{\cv\md}$, $\widehat{\nu}_{\cv\md}$, $\mathcal{S}$ and
$\widehat{\mathcal{S}}(t)$ as above, $\Dbl = \lVert\widehat{g}_{\cv\md} -
\gcomp{\cv}{\md}\rVert_{L_2[0,1]}$, and
$\widehat{g}^{(t)}_{\cv\md} = \widehat{g}_{\cv\md}
\mathbf{1}\{(\cv,\md)\in\widehat{\mathcal{S}}(t)\}$.
\begin{enumerate}
\item[(i)] With no assumption,
\[
  \sum_{\cv,\md}\bigl\lVert\widehat{g}^{(t)}_{\cv\md}-\gcomp{\cv}{\md}
  \bigr\rVert_{L_2[0,1]}^{2}
  = \sum_{(\cv,\md)\in\widehat{\mathcal{S}}(t)}\Dbl^{2}
  + \sum_{(\cv,\md)\in\mathcal{S}\setminus\widehat{\mathcal{S}}(t)}\nu_{\cv\md}^{2},
\]
and every $(\cv,\md) \in \mathcal{S}\setminus\widehat{\mathcal{S}}(t)$ has
$\nu_{\cv\md} \le t + \Dbl$; hence the left side is at most
$2\sum_{\cv,\md}\Dbl^{2} + 2t^{2}|\mathcal{S}\setminus\widehat{\mathcal{S}}(t)|
\le 2\sum_{\cv,\md}\Dbl^{2} + 2t^{2}|\mathcal{S}|$.
\item[(ii)] If $\mathcal{S} \subseteq \widehat{\mathcal{S}}(t)$, the left side
of (i) is $\sum_{\widehat{\mathcal{S}}(t)}\Dbl^{2} \le \sum_{\cv,\md}\Dbl^{2}$;
if $\widehat{\mathcal{S}}(t) = \mathcal{S}$, it is $\sum_{\mathcal{S}}\Dbl^{2}$,
and the null blocks no longer contribute.
\item[(iii)] Let $\lmin(\Gramhat) > 0$, $\betastar = (\bc,\thetastar)$,
$f_J = \bZ\betastar$, $\betahat = (\chat,\thetahat)$ any fit,
$\widehat{\mathcal{S}}(t) = \{(\cv,\md) : \norm{\thetahat_{\cv\md,\cdot}}_{2} > t\}$,
$\mathcal{S}_J = \{(\cv,\md) : \thetastar_{\cv\md,\cdot} \ne 0\}$,
$\thetahat^{(t)}$ the vector $\thetahat$ with the blocks outside
$\widehat{\mathcal{S}}(t)$ set to zero, and
$\widehat{f}^{(t)} = \Amat\chat + \Bmat\thetahat^{(t)}$. Then
\[
  \normn{\widehat{f}^{(t)}-f_J}^{2} \;\le\;
  \frac{2\lmaxx(\Gramhat)}{\lmin(\Gramhat)}\normn{\widehat{f}-f_J}^{2}
  + 2\lmaxx(\Gramhat)\,t^{2}\,|\mathcal{S}_J\setminus\widehat{\mathcal{S}}(t)|,
\]
and, if $\mathcal{S}_J \subseteq \widehat{\mathcal{S}}(t)$,
$\normn{\widehat{f}^{(t)}-f_J}^{2} \le
\{\lmaxx(\Gramhat)/\lmin(\Gramhat)\}\normn{\widehat{f}-f_J}^{2}$.
\end{enumerate}
\end{lemma}

\begin{proof}
\emph{(i).} If $(\cv,\md) \in \widehat{\mathcal{S}}(t)$, then
$\widehat{g}^{(t)}_{\cv\md} = \widehat{g}_{\cv\md}$ and the error of the block
is $\Dbl^{2}$; otherwise $\widehat{g}^{(t)}_{\cv\md} = 0$ and the error is
$\nu_{\cv\md}^{2}$, which vanishes outside $\mathcal{S}$. Summing gives the
identity. For $(\cv,\md) \notin \widehat{\mathcal{S}}(t)$,
$\widehat{\nu}_{\cv\md} \le t$, and the triangle inequality gives
$\nu_{\cv\md} \le \widehat{\nu}_{\cv\md} + \Dbl \le t + \Dbl$, hence
$\nu_{\cv\md}^{2} \le 2t^{2} + 2\Dbl^{2}$; summing over
$\mathcal{S}\setminus\widehat{\mathcal{S}}(t)$ and adding the sum over
$\widehat{\mathcal{S}}(t)$, a disjoint set, gives the bound.
\emph{(ii).} If $\mathcal{S} \subseteq \widehat{\mathcal{S}}(t)$, the second sum
of (i) is empty.
\emph{(iii).} The proof of (i) and (ii) uses only the triangle inequality in a
norm; applied to the vectors $\thetahat_{\cv\md,\cdot}$ and
$\thetastar_{\cv\md,\cdot}$ of $\R^{\NJ}$, with $\mathcal{S}_J$ in place of
$\mathcal{S}$, it gives $\norm{\thetahat^{(t)}-\thetastar}_{2}^{2} \le
2\norm{\thetahat-\thetastar}_{2}^{2} +
2t^{2}|\mathcal{S}_J\setminus\widehat{\mathcal{S}}(t)|$, and
$\norm{\thetahat^{(t)}-\thetastar}_{2}^{2} \le \norm{\thetahat-\thetastar}_{2}^{2}$
when $\mathcal{S}_J \subseteq \widehat{\mathcal{S}}(t)$. Since $\chat$ does not
change, the same bounds hold for $\betahat^{(t)} = (\chat,\thetahat^{(t)})$
against $\betastar$. Finally
$\normn{\widehat{f}^{(t)}-f_J}^{2} = (\betahat^{(t)}-\betastar)\T\Gramhat
(\betahat^{(t)}-\betastar) \le
\lmaxx(\Gramhat)\norm{\betahat^{(t)}-\betastar}_{2}^{2}$ and
$\norm{\betahat-\betastar}_{2}^{2} \le
\normn{\bZ(\betahat-\betastar)}^{2}/\lmin(\Gramhat) =
\normn{\widehat{f}-f_J}^{2}/\lmin(\Gramhat)$. By the orthonormality of the
basis, $\norm{\thetahat_{\cv\md,\cdot}}_{2} = \widehat{\nu}_{\cv\md}$, so that the
$\widehat{\mathcal{S}}(t)$ of (iii) is that of (i).
\end{proof}

The factor $2$ in (i) is attained: an active block with
$\widehat{g}_{\cv\md} = t\,\gcomp{\cv}{\md}/\nu_{\cv\md}$ and
$\nu_{\cv\md} = 2t$ has $\widehat{\nu}_{\cv\md} = t$, is set to zero, and pays
$\nu_{\cv\md}^{2} = 4t^{2} = 2\Dbl^{2} + 2t^{2}$. In (iii), the passage between
the norm of the coefficients and the prediction norm uses the smallest
eigenvalue of $\Gramhat$ to go from prediction to coefficients and the largest
to come back, and the factor is the condition number of $\Gramhat$, which on
$\ESig$ is smaller than $2\Lambda/\gamma$.


{\color{gray}\noindent[k = 4: the proof of Corollary~2 for the lasso, which stood
here, is moved to Section~S8, unchanged except for the numbering of the results
it cites.]\par}

% derivations/06-selecao-limiar.tex, provas do Corolario 12, da Observacao
% rem:adaptativa-blocos e do Corolario 13 (adendo de E1.12).
\subsection*{Proof of Corollary 2}

\emph{A nonasymptotic bound on $\Delta$.} Work at $J = J_n$ and
$\pen = \lamG$, and intersect $\Acal$, the event $\Ecal$ of
Lemma~\ref{lem:calib-blocks}(i) and
$\{\normn{\bb}^{2} \le \mathcal{B}_{J_n}^{2}/\alpha'\}$; the intersection has
probability at least $1 - \alpha - \alpha' - \Prob(\Acal^{c})$, by
Lemma~\ref{lem:calib-blocks}, Markov's inequality and Lemma~2. On $\Acal$,
$\Ecal \subseteq \Tw$ and $\gtil > \gamma/2$ (Lemma~\ref{lem:regime}), and
Theorem~2(ii) with the comparator $\betastar$ and
$\Wcal(\Gzero(\thetastar)) \le \Wcal(\Gcal_{n})$ gives, with
$v = \thetahat - \thetastar$,
\[
  \norm{v}_{2}^{2} \;\le\; \frac{64(\lamG)^{2}\Wcal(\Gcal_{n})}{\gtil^{2}}
  + \frac{16\normn{\bb}^{2}}{\gtil}
  \;\le\; \frac{256(\lamG)^{2}\Wcal(\Gcal_{n})}{\gamma^{2}}
  + \frac{32\mathcal{B}^{2}_{J_n}}{\alpha'\gamma} .
\]
For a fixed block, the orthonormality of the wavelets of level below $J_n$
gives $\lVert\widehat{g}_{\cv\md} - \PiJ\gcomp{\cv}{\md}\rVert_{L_2[0,1]} =
\norm{v_{\cv\md,\cdot}}_2 \le \norm{v}_2$, and the first display of Lemma~2
bounds $\lVert\PiJ\gcomp{\cv}{\md} - \gcomp{\cv}{\md}\rVert^2_{L_2[0,1]}$ by
$C_g^2(1-2^{-2\seff})^{-1}2^{-2J_n\seff}$. By $(x+y)^2 \le 2x^2 + 2y^2$,
\begin{equation}
  \label{eq:Dn-blocks}
  \Delta^{2} \;\le\; (\bar{\Delta}^{\Gcal}_n)^{2} \;=\;
  \frac{512(\lamG)^{2}\Wcal(\Gcal_{n})}{\gamma^{2}}
  + \frac{64\,\mathcal{B}^2_{J_n}}{\alpha'\gamma}
  + \frac{2C_g^2}{1-2^{-2\seff}}\,2^{-2J_n\seff},
\end{equation}
with $(\lamG)^{2}\Wcal(\Gcal_{n}) \le \varrho^{2}(\lamGone)^{2}|\Gcal_{n}|$ by
Lemma~\ref{lem:calib-blocks}(iii). The residual probability
$\Prob(\Acal^{c})$ tends to zero by Lemma~\ref{lem:regime} and does not depend
on the components; it is the $\omega_n$ of Corollary~2(iii), and by
Lemma~S7.1(iii) this paragraph is the proof of that part. The bound
$\bar{\Delta}^{\Gcal}_n$ depends on the components only through the constants
of Assumption~5, so everything above holds uniformly over the class, which is
what allows the target to vary with $n$ in Assumption~6.

\emph{The order of $\bar{\Delta}^{\Gcal}_n$.} By the proof of Theorem~3,
$(\lamG)^{2}\Wcal(\Gcal_{n}) \le \bar\varrho^{2}C_{\pen}\sigma^{2}pq(\bn +
2^{J_n})/n \asymp n^{-2\seff/(2\seff+1)}$, and $2^{-2J_n\seff} \le
n^{-2\seff/(2\seff+1)}$ by the definition of $J_n$. Hence
$\bar{\Delta}^{\Gcal}_n \le c(\alpha,\alpha')\rhoG$ with
$c(\alpha,\alpha') < \infty$. Given $\epsilon > 0$, take
$\alpha = \alpha' = \epsilon/4$ and $n$ large enough that
$\Prob(\Acal^{c}) < \epsilon/2$; then
$\Prob\{\Delta > c(\epsilon/4,\epsilon/4)\rhoG\} < \epsilon$, that is
$\Delta = O_p(\rhoG)$, uniformly over the class.

\emph{Part (i).} If $t_n/\rhoG \to \infty$, then $\Prob(\Delta \le t_n) \to 1$,
and on that event Lemma~S7.1(i) with $\bar{\Delta} = \Delta$ gives
$\widehat{\mathcal{S}}(t_n) \subseteq \mathcal{S}$. No property of
$\nu_{\min,n}$ has been used.

\emph{Part (ii).} If $\mathcal{S}$ is empty, part (ii) is part (i). Otherwise,
the condition $\limsup_n t_n/\nu_{\min,n} < 1$ gives $\eta > 0$ and $n_0$ with
$\nu_{\min,n} - t_n \ge \eta\,\nu_{\min,n}$ for $n \ge n_0$, and Assumption~6
gives $(\nu_{\min,n} - t_n)/\rhoG \ge \eta\,\nu_{\min,n}/\rhoG \to \infty$.
Hence $\Prob(\Delta \le t_n) \to 1$ and
$\Prob(\Delta < \nu_{\min,n} - t_n) \to 1$, and on the intersection of the two
events Lemma~S7.1(iii) with $\bar{\Delta} = \Delta$ gives
$\widehat{\mathcal{S}}(t_n) = \mathcal{S}$. For the existence,
$t_n = (\rhoG\nu_{\min,n})^{1/2}$ has
$t_n/\rhoG = (\nu_{\min,n}/\rhoG)^{1/2} \to \infty$ and
$t_n/\nu_{\min,n} = (\rhoG/\nu_{\min,n})^{1/2} \to 0$. \qed

\emph{At the resolution of Theorem~1.} Theorem~1 gives
$\sum_{\cv,\md}\lVert\widehat{g}_{\cv\md}-\gcomp{\cv}{\md}\rVert^2_{L_2[0,1]}
= O_p\{(\rhoGt)^{2}\}$, with $\rhoGt$ as in Section~3.6 of the paper, hence
$\Delta = O_p(\rhoGt)$, and the bound in its proof depends on the components
only through the constants of Assumption~5 and of the design, so that this
holds uniformly over the class. The proofs of parts (i) and (ii) above use
nothing about $\Delta$ beyond $\Delta = O_p(\rhoG)$ uniformly over the class,
together with Lemma~S7.1 and Assumption~6, and they go through verbatim with
$\rhoGt$ in place of $\rhoG$.

\subsection*{Proof of Corollary 3}

Write $r_n$ for the rate of the statement and $\Dbl$ as in
Lemma~\ref{lem:thrisk}. For the fit of Theorem~3,
$\normn{\widehat{f}-f}^{2} = O_p(r_n^{2})$ and
$\sum_{\cv,\md}\Dbl^{2} = O_p(r_n^{2})$ by Theorem~3 and Corollary~1; for the
lasso, by Theorem~\ref{thm:lasso-linear} and
Corollary~\ref{cor:lasso-components}. In both cases
$\Prob(\ESig) \to 1$, by the proofs of Theorems~3 and \ref{thm:lasso-linear},
and $\normn{\bb}^{2} = O_p(\mathcal{B}_{J_n}^{2}) = O_p(r_n^{2})$, by the choice
of $J_n$.

\emph{Part (i).} The bound for the components is Lemma~\ref{lem:thrisk}(i)
with $|\mathcal{S}| \le pq$. For prediction,
$\normn{\widehat{f}^{(t)}-f}^{2} \le 2\normn{\widehat{f}^{(t)}-f_{J_n}}^{2} +
2\normn{\bb}^{2}$, and on $\ESig$, $\lmaxx(\Gramhat) < \Lambda$ and
$\lmin(\Gramhat) > \gamma/2$, so that Lemma~\ref{lem:thrisk}(iii) gives, for
every $t$,
\[
  \normn{\widehat{f}^{(t)}-f}^{2} \;\le\;
  \frac{8\Lambda}{\gamma}\normn{\widehat{f}-f_{J_n}}^{2}
  + 4\Lambda t^{2}|\mathcal{S}_{J_n}\setminus\widehat{\mathcal{S}}(t)|
  + 2\normn{\bb}^{2},
\]
with $\normn{\widehat{f}-f_{J_n}}^{2} \le 2\normn{\widehat{f}-f}^{2} +
2\normn{\bb}^{2} = O_p(r_n^{2})$.

\emph{Part (ii).} By Lemma~S7.1(ii), $\mathcal{S} \subseteq
\widehat{\mathcal{S}}(t_n)$ whenever $t_n < \nu_{\min,n} - \Delta$. With
$\limsup_n t_n/\nu_{\min,n} < 1$ there are $\eta > 0$ and $n_0$ with
$\nu_{\min,n} - t_n \ge \eta\nu_{\min,n}$ for $n \ge n_0$, and since
$\Delta^{2} \le \sum_{\cv,\md}\Dbl^{2} = O_p(r_n^{2})$ and
$\nu_{\min,n}/r_n \to \infty$, $\Prob(\Delta < \nu_{\min,n} - t_n) \to 1$. On
that event Lemma~\ref{lem:thrisk}(ii) holds and, since
$\mathcal{S}_{J_n} \subseteq \mathcal{S}$, so does the second case of
Lemma~\ref{lem:thrisk}(iii), which on $\ESig$ gives
$\normn{\widehat{f}^{(t_n)}-f}^{2} \le
(4\Lambda/\gamma)\normn{\widehat{f}-f_{J_n}}^{2} + 2\normn{\bb}^{2}$. If
moreover $t_n/r_n \to \infty$, part (ii) of Corollary~2, or of
Corollary~\ref{cor:lasso-threshold} for the lasso, gives
$\Prob\{\widehat{\mathcal{S}}(t_n) = \mathcal{S}\} \to 1$, and
Lemma~\ref{lem:thrisk}(ii) concludes. \qed

}

%%%%%%%%%%%%%%%%%%%%%%%%%%%%%%%%%%%%%%%%%%%%%%%%%%%%%%%%%%%%%%%%%%%%%%%%%%%%%%
% derivations/04-oraculo.tex, 05-taxas.tex e 06-selecao-limiar.tex (o lasso
% coordenado). Secao nova em k = 4 (D51(a)): os enunciados estavam nas Secoes
% 3.2 a 3.6 do ms_3 e as provas nas Secoes S5 a S7 do supp_3; a Proposicao 4
% entra emendada por D49(a) (as condicoes s' > 1/4 em (i) e s' > s/2 em (ii)).
{\color{colR1}
\section{The coordinatewise lasso}
\setcounter{equation}{0}

\noindent[k = 4: the statements of this section stood in Sections 3.2 to 3.6 of
the paper, and their proofs in Sections S5 to S7 of this supplement.]

This section gives the theory of the estimator obtained by replacing
$\nGw{\btheta}$ by $\norm{\btheta}_{1}$ in equation (2.5) of the paper, the
lasso of \citet{tibshirani1996regression} with the levels left unpenalized,
which the software offers as an option. Assumptions 1 to 5, Lemmas 1 and 2,
Proposition 2 and Lemmas~\ref{lem:profile}, \ref{lem:columns} and \ref{lem:l1}
are used as they stand, with the notation $\Szero$, $s_0$ and $v$ of
Section~S1. The penalty level is
\begin{equation}
  \label{eq:lasso-lambda}
  \penn \;=\; 2\sigma B_X\sqrt{2C_U}\,
  \sqrt{\frac{2\log(2d_n/\alpha)}{n}} \;\asymp\;
  \sigma B_X \sqrt{C_U}\,\sqrt{\frac{J_n}{n}} ,
\end{equation}
with $d_n$ the number of penalized columns at resolution $J_n$.

\subsection*{Statements}

\begin{theorem}[Oracle inequality for the lasso]
\label{thm:lasso-oracle}
Let Assumption 1 hold, suppose $\Amat$ has rank $p$, and take
$\pen \ge 2\pen_0$ with
$\pen_0 = \sigma\smax\{2\log(2d/\alpha)/n\}^{1/2}$. On the event
$\mathcal{T} = \{\norm{\Btil\T\bveps/n}_{\infty} \le \pen/2\}$, which has
probability at least $1 - \alpha$:
\begin{enumerate}
\item[(i)] with no condition on the design,
$\normn{\Btil v}^{2} \le 6\pen\norm{\thetastar}_{1} + 4\normn{\bb}^{2}$;
\item[(ii)] if $\gtil > 0$,
\[
  \normn{\Btil v}^{2} \le \frac{64\pen^{2}s_0}{\gtil} + 16\normn{\bb}^{2},
  \qquad
  \norm{v}_{2}^{2} \le \frac{64\pen^{2}s_0}{\gtil^{2}}
  + \frac{16\normn{\bb}^{2}}{\gtil} ;
\]
\item[(iii)] in either regime,
$\normn{\widehat{f}-f} \le \normn{\Btil v} + 2\normn{\bb}
+ \normn{\PA\bveps}$, and
$\E(\normn{\PA\bveps}^{2}\mid\bX,\bU) \le \sigma^{2}p/n$.
\end{enumerate}
\end{theorem}

\begin{theorem}[Rate of linear approximation for the lasso]
\label{thm:lasso-linear}
Let Assumptions 1 to 5 hold with $p$, $q$ and the constants fixed, and take
\begin{equation}
  \label{eq:lasso-Jn}
  J_n = \min\bigl\{J \in \mathbb{N} :\;
  2^{J} \ge (n/\log n)^{1/(2\seff+1)}\bigr\},
\end{equation}
so that $J_n = \log_2 n/(2\seff+1) + O(\log\log n)$. Then the requirement of
Proposition 2(ii) holds along $J_n$, being itself of the exact order
$(\log n/n)^{2\seff/(2\seff+1)}$, and, with $\penn$ as in
\eqref{eq:lasso-lambda},
\[
  \normn{\widehat{f}-f}^{2}
  \;=\; O_p\Bigl\{\Bigl(\frac{\log n}{n}\Bigr)^{2\seff/(2\seff+1)}\Bigr\} .
\]
At this resolution the estimation term and the bias term are of the same
order, and the term $\sigma^{2}p/n$ of Theorem~\ref{thm:lasso-oracle}(iii) is
of strictly smaller order.
\end{theorem}

\begin{corollary}[Components and coefficients for the lasso]
\label{cor:lasso-components}
In the setting of Theorem~\ref{thm:lasso-linear}, and with $r_n$ denoting the
rate of that theorem,
\[
  \sum_{\cv,\md}\bigl\lVert\widehat{g}_{\cv\md}-\gcomp{\cv}{\md}
  \bigr\rVert_{L_2[0,1]}^{2} = O_p(r_n),
  \qquad
  \bigl\lVert\widehat{\fcoef{\cv}}-\fcoef{\cv}
  \bigr\rVert_{L_2([0,1]^{q})} = O_p(r_n^{1/2})
\]
for $\cv = 1,\dots,p$, and the same statements hold with the rate of
Theorem~\ref{thm:lasso-main} in place of $r_n$ under the conditions of that
theorem.
\end{corollary}

% E5e (E1.14; D52): a inflacao de penn acima do teto, eq. mu de 05-taxas.tex.
The level $\penn$ is admissible only while $2^{J_n}\log d_n/n \to 0$, which
Lemma~\ref{lem:columns} needs for $\smax^{2} \le 2B_X^{2}C_U$; above that
ceiling $\smax^{2}$ grows with $2^{J}/n$. With $c_{\psi} =
7\norm{\psi}_{\infty}^{2}/(12C_U)$ as in Section~S6, write
\begin{equation}
  \label{eq:lasso-lamplus}
  \mu_{J} = 1 + c_{\psi}\frac{2^{J}\log d}{n},
  \qquad
  \penn^{+} = \mu_{J_n}^{1/2}\penn ;
\end{equation}
when $2^{J_n}\log d_n/n \to 0$, $\penn^{+} = \penn\{1 + o(1)\}$.

% derivations/05-taxas.tex, Proposicao 4, emendada por D49(a), com os itens
% (iii) e (iv) de E1.14 (E5e).
\begin{proposition}[Slow rate for the lasso, no design condition]
\label{prop:lasso-slow}
% E5f (D54): as cotas superiores das Assumptions 2 e 3, que o penn e a prova usam.
Let Assumptions 1, 4 and 5 hold, \textcolor{colR1}{together with the upper
bounds in Assumptions 2 and 3,} with $\Amat$ of rank $p$ and $\penn$ as in
\eqref{eq:lasso-lambda}. Then
$\normn{\widehat{f}-f}^{2} = O_p\{\penn\norm{\thetastar}_{1} +
\mathcal{B}_{J_n}^{2} + p/n\}$, where
$\norm{\thetastar}_{1} \le pqC_g\sum_{\lev<J_n}2^{\lev(1/2-s)}$ does not
depend on $\pi$, whenever $2^{J_n}\log d_n/n \to 0$; with $\penn^{+}$ of
\eqref{eq:lasso-lamplus} in place of $\penn$, in the estimator and in the
bound, the same holds for every sequence $J_n$. In particular, (i) if $s > 1/2$ and $\seff > 1/4$ then
$\normn{\widehat{f}-f}^{2} = O_p\{(J_n/n)^{1/2}\}$ for every
$J_n \ge \log_2 n/(4\seff)$; and (ii) if $s < 1/2$ and $\seff > s/2$, the choice
$2^{J_n} \asymp (n/J_n)^{1/(1-2s+4\seff)}$ gives
$\normn{\widehat{f}-f}^{2} = O_p\{(J_n/n)^{2\seff/(1-2s+4\seff)}\}$, an
exponent that equals $2s/(2s+1)$ when $\pi \ge 2$. With $\penn^{+}$:
\begin{enumerate}
\item[(iii)] (the pairs that (i) and (ii) exclude, all with $\pi < 2$)
(a) if $s > 1/2$ and $\seff \le 1/4$, the choice
$2^{J_n} \asymp (n/\log n)^{2/(1+4\seff)}$ gives
$\normn{\widehat{f}-f}^{2} = O_p\{(\log n/n)^{4\seff/(1+4\seff)}\}$;
(b) if $s < 1/2$ and $\seff \le s/2$, the choice
$2^{J_n} \asymp (n/\log n)^{1/(1-s+2\seff)}$ gives
$\normn{\widehat{f}-f}^{2} = O_p\{(\log n/n)^{2\seff/(1-s+2\seff)}\}$;
in both, $2^{J_n} \gtrsim n/\log n$, above the ceiling of
Lemma~\ref{lem:columns}, and the rates join those of (i) and (ii) at
$\seff = 1/4$ and $\seff = s/2$;
\item[(iv)] (the boundary $s = 1/2$, with $\pi > 1$)
$\norm{\thetastar}_{1} \le pqC_gJ_n$, and: if $\seff \ge 1/4$, that is
$\pi \ge 4/3$, the choice $2^{J_n} \asymp (n/\log n)^{1/(4\seff)}$ gives
$\normn{\widehat{f}-f}^{2} = O_p\{n^{-1/2}(\log n)^{3/2}\}$; if
$\seff < 1/4$, that is $1 < \pi < 4/3$, the choice
$2^{J_n} \asymp \{n/(\log n)^{2}\}^{2/(1+4\seff)}$ gives
$\normn{\widehat{f}-f}^{2} = O_p[\{(\log n)^{2}/n\}^{4\seff/(1+4\seff)}]$.
\end{enumerate}
\end{proposition}

% E5f (E1.15; D54): derivations/05-taxas.tex, Proposicao 8 (numeracao global;
% prop:truncada); a prova vem depois da Remark S8.2.
Proposition~\ref{prop:lasso-slow} reads Theorem~\ref{thm:lasso-oracle}(i) at
$\thetastar$ and pays $\pen|\theta^{*}_{a}|$ also for the coefficients
smaller than $\pen$, which would cost less as bias. Lemma~\ref{lem:comparator}
below holds for every comparator on the same event, and the comparator
truncated at $t > 0$,
$\thbar_{a} = \theta^{*}_{a}\mathbf{1}\{|\theta^{*}_{a}| > t\}$ for
$a = 1,\dots,d$, charges the penalty only for the coefficients above $t$. What
remains is
\begin{equation}
  \label{eq:R1}
  \Rone(\btheta;t) \;=\; \sum_{a=1}^{d}\min\bigl(t|\theta_{a}|,\,\theta_{a}^{2}\bigr)
  \;\le\; t\norm{\btheta}_{1},
\end{equation}
the counterpart in the $\ell_1$ norm of the ideal risk
$\sum_{a}\min(\theta_{a}^{2},t^{2})$.

\begin{proposition}[Slow rate for the lasso through the truncated comparator]
\label{prop:lasso-truncated}
Let the assumptions of Proposition~\ref{prop:lasso-slow} hold.
\begin{enumerate}
\item[(i)] On the event $\mathcal{T}$, for every $\pen \ge 2\pen_0$ and
every $t > 0$, with $\thbar$ truncated at $t$,
\[
  \normn{\widehat{f}-f}^{2}
  \;\le\; 18\pen\norm{\thbar}_{1} + 48\normn{\Bmat(\thetastar-\thbar)}^{2}
  + 48\normn{\bb}^{2} + 3\normn{\PA\bveps}^{2} .
\]
If $t$ is deterministic,
$\E\normn{\Bmat(\thetastar-\thbar)}^{2} \le
\kappa_2C_U\norm{\thetastar-\thbar}_{2}^{2}$. With $t = \pen$ deterministic,
$18\pen\norm{\thbar}_{1} + 48\kappa_2C_U\norm{\thetastar-\thbar}_{2}^{2} \le
\max(18,48\kappa_2C_U)\Rone(\thetastar;\pen)$, and
\[
  \normn{\widehat{f}-f}^{2} \;=\; O_p\bigl\{\Rone(\thetastar;\penn)
  + \mathcal{B}_{J_n}^{2} + p/n\bigr\}
\]
whenever $2^{J_n}\log d_n/n \to 0$; with $\penn^{+}$ of
\eqref{eq:lasso-lamplus} in place of $\penn$, in the estimator and in the
bound, the same holds for every sequence $J_n$ with $2^{J_n} \le n^{2}$.
\item[(ii)] For every $J$ and every $t > 0$, with $\tau = (s+1/2)^{-1}$: if
$s < 1/2$,
\begin{equation}
  \label{eq:count-lasso-truncated}
  \Rone(\thetastar;t) \;\le\; pq\,C_g^{\tau}\,t^{4s/(2s+1)}
  \Bigl\{\frac{1}{1-2^{s-1/2}} + \frac{1}{1-2^{-\min(\pi,2)\seff}}\Bigr\};
\end{equation}
if $s \ge 1/2$, nothing is gained in order: $\Rone(\thetastar;t) \le
t\norm{\thetastar}_{1}$, and for $t \le C_g$ the class of Assumption~5
contains a $\thetastar$ with $\Rone(\thetastar;t) \ge tC_g$ if $s > 1/2$,
and one with $\Rone(\thetastar;t) \ge
tC_g\min\{J,\,\lfloor\log_2(C_g/t)\rfloor + 1\}$ if $s = 1/2$, the order of
$t\norm{\thetastar}_{1}$ in Proposition~\ref{prop:lasso-slow}.
\item[(iii)] If $s < 1/2$ and $\seff > s/(2s+1)$, then with
$J_n = \lceil c\log_2 n\rceil$ and $c$ in the window \eqref{eq:lasso-window},
\[
  \normn{\widehat{f}-f}^{2}
  \;=\; O_p\Bigl\{\Bigl(\frac{\log n}{n}\Bigr)^{2s/(2s+1)}\Bigr\},
\]
the prediction rate of Theorem~\ref{thm:lasso-main}.
\item[(iv)] If $s < 1/2$ and $\seff \le s/(2s+1)$, which forces $\pi < 2$,
then with $\penn^{+}$ and the choice
$2^{J_n} \asymp (n/\log n)^{2s/\{s+\seff(2s+1)\}}$, which has
$2^{J_n} \gtrsim n/\log n$,
\[
  \normn{\widehat{f}-f}^{2}
  \;=\; O_p\Bigl\{\Bigl(\frac{\log n}{n}\Bigr)^{4s\seff/\{s+\seff(2s+1)\}}\Bigr\};
\]
at $\seff = s/(2s+1)$ the exponent is $2s/(2s+1)$, that of (iii).
\end{enumerate}
\end{proposition}

Proposition~\ref{prop:lasso-truncated} is a statement on prediction: it uses
the design as Proposition~\ref{prop:lasso-slow} does, and, in the new bias,
the upper bound $\lmaxx(\Gram) \le \kappa_2C_U$; the components still need
the design condition (Corollary~\ref{cor:lasso-components}). When $s < 1/2$
and $\pi < 2$ it improves on Proposition~\ref{prop:lasso-slow} in every pair:
(iii) gives the rate of Theorem~\ref{thm:lasso-main} in the window of that
theorem, with a gain of $2(s-\seff)(1-2s)/\{(2s+1)(1-2s+4\seff)\}$ in the
exponent over Proposition~\ref{prop:lasso-slow}(ii), and (iv) replaces
Proposition~\ref{prop:lasso-slow}(iii)(b) and the part of
Proposition~\ref{prop:lasso-slow}(ii) with $s/2 < \seff \le s/(2s+1)$. When
$\pi \ge 2$ the rate is that of Proposition~\ref{prop:lasso-slow}(ii), now
for every $c \in [1/(2s+1),1)$. When $s \ge 1/2$ nothing changes, for the
reason given after Proposition~\ref{prop:truncated}.

For a vector $\btheta$ let $\theta_{(1)} \ge \theta_{(2)} \ge \cdots$ be the
nonincreasing rearrangement of the absolute values of its entries and let
$\btheta^{(M)}$ retain the $M$ largest of them.

\begin{lemma}[Besov regularity implies compressibility]
\label{lem:lasso-weak}
Under Assumptions 4 and 5 with $\pi(s+1/2) > 1$, the oracle satisfies, for
every $J$ and every $k \ge 1$,
\[
  \theta^{*}_{(k)} \;\le\; C_\tau\,k^{-1/\tau},
  \qquad \tau = \frac{1}{s+1/2},
  \qquad C_\tau = C_g(pq\,A_{s,\pi})^{s+1/2},
  \qquad A_{s,\pi} = 1 + \frac{1}{1-2^{1-\pi(s+1/2)}},
\]
with constants that do not depend on $J$ and with $\tau \in (0,2)$ for every
$s > 0$. Consequently
$\sum_{k>M}\theta^{*2}_{(k)} \le \tau(2-\tau)^{-1}C_\tau^{2}M^{1-2/\tau}$ for
every $M \ge 1$.
\end{lemma}

The inclusion of Besov balls in weak $\ell_\tau$ balls is classical
\citep[Proposition~9.15]{johnstone2019gaussian}; what the lemma adds is the
constant $C_\tau$, uniform in $J$, for the oracle of the product design, which
is the form in which the proof of Theorem~\ref{thm:lasso-main} uses it.

\begin{theorem}[Rate of nonlinear approximation for the lasso]
\label{thm:lasso-main}
Let Assumptions 1 to 5 hold with $p$, $q$ and all the constants fixed as $n$
grows, and set $\tau = (s+1/2)^{-1}$. Suppose $\seff > s/(2s+1)$ and take
$J_n = \lceil c \log_2 n\rceil$ with
\begin{equation}
  \label{eq:lasso-window}
  \max\Bigl\{\frac{\tau}{2},\; \frac{2-\tau}{4\seff}\Bigr\}
  \;\le\; c \;<\; 1 ,
\end{equation}
a window that is nonempty precisely under that condition. Then, with $\penn$
as in \eqref{eq:lasso-lambda},
\[
  \normn{\widehat{f}-f}^{2}
  \;=\; O_p\Bigl\{\Bigl(\frac{\log n}{n}\Bigr)^{2s/(2s+1)}\Bigr\},
  \qquad
  \sum_{\cv,\md}\bigl\lVert \widehat{g}_{\cv\md}-\gcomp{\cv}{\md}
  \bigr\rVert_{L_2[0,1]}^{2}
  \;=\; O_p\Bigl\{\Bigl(\frac{\log n}{n}\Bigr)^{2s/(2s+1)}\Bigr\} .
\]
\end{theorem}

The logarithmic factor is the price of the single penalty level
\eqref{eq:lasso-lambda}, which takes a union bound over the $d_n$ columns;
Theorem~1 of the paper removes the factor $(\log n)^{2\seff/(2s+1)}$ of it.
Write $\rho_n = (\log n/n)^{\seff/(2\seff+1)}$ for the square root of the rate
of Theorem~\ref{thm:lasso-linear}.

\begin{corollary}[Block structure by thresholding the lasso]
\label{cor:lasso-threshold}
Let the assumptions of Theorem~\ref{thm:lasso-linear} hold, with $J_n$ and
$\penn$ as there.
\begin{enumerate}
\item[(i)] (Screening.) If the thresholds satisfy $t_n/\rho_n \to \infty$,
then $\Prob\{\widehat{\mathcal{S}}(t_n) \subseteq \mathcal{S}\} \to 1$.
\item[(ii)] (Recovery.) If, in addition, Assumption 6 holds with $\rho_n$ in
place of $\rhoG$ and $\limsup_n t_n/\nu_{\min,n} < 1$, then
$\Prob\{\widehat{\mathcal{S}}(t_n) = \mathcal{S}\} \to 1$. Such thresholds
always exist; one is $t_n = (\rho_n\nu_{\min,n})^{1/2}$.
\item[(iii)] (A window at fixed $n$.) With $\alpha$ the level in
\eqref{eq:lasso-lambda}, for every $n$ and every $\alpha' \in (0,1)$ there is
an event, of probability at least $1 - \alpha - \alpha' - \omega_n$ with
$\omega_n \to 0$ not depending on the components, on which
$\widehat{\mathcal{S}}(t) = \mathcal{S}$ for every $t$ with
$\bar{\Delta}_n \le t < \nu_{\min,n} - \bar{\Delta}_n$; here
$\bar{\Delta}_n \le c(\alpha,\alpha')\rho_n$ is the explicit bound
\eqref{eq:Dn} on $\max_{(\cv,\md)}\lVert \widehat{g}_{\cv\md} -
\gcomp{\cv}{\md}\rVert_{L_2[0,1]}$.
\end{enumerate}
\end{corollary}

Parts (i) and (ii) hold also at the resolution of
Theorem~\ref{thm:lasso-main}, with $\rho_n$ replaced by
$\tilde\rho_n = (\log n/n)^{s/(2s+1)}$, the square root of the rate of that
theorem, both in the condition on $t_n$ and in Assumption~6. Since
$\tilde\rho_n \le \rho_n$, with a gap in the exponent when $\pi < 2$, the
separation required is weaker exactly where the components are spatially
inhomogeneous.

\subsection*{Proofs}

\begin{lemma}[Calibration of the penalty]
\label{lem:calib}
Under Assumption 1, for every $\alpha \in (0,1)$,
$\Prob(\norm{\Btil\T\bveps/n}_{\infty} > \pen_0) \le \alpha$ with
$\pen_0 = \sigma\smax\{2\log(2d/\alpha)/n\}^{1/2}$.
\end{lemma}

\begin{proof}
Condition on the covariates, so that $\Btil$ and $\smax$ are fixed and the
errors are independent and sub-Gaussian with parameter $\sigma$. For a column
$a$, $(\Btil\T\bveps)_a/n$ is sub-Gaussian with parameter
$\sigma\lVert\Btil^{(a)}\rVert_2/n$, and, $\MA$ being an orthogonal
projection, $\lVert\Btil^{(a)}\rVert_2 \le \lVert\Bmat^{(a)}\rVert_2 =
\sqrt{n}\,(\Gramhat)_{aa}^{1/2} \le \sqrt{n}\,\smax$. With
$t = \pen_0$ the sub-Gaussian tail bound is at most $\alpha/d$ per column; a
union bound over the $d$ columns and the tower property finish the proof.
\end{proof}

\subsection*{Proof of Theorem~\ref{thm:lasso-oracle}}

\emph{Step 1, the basic inequality.} By Lemma S5.1, $\thetahat$ minimizes the
residualized objective; comparing its value at $\thetahat$ with its value at
$\thetastar$, and using $\MA\Amat = 0$, so that
$\MA\bY = \Btil\thetastar + \MA(\bb+\bveps)$, gives, with
$r = \MA(\bb+\bveps)$,
\begin{equation}
  \label{eq:basic}
  \normn{\Btil v}^2 + 2\pen\norm{\thetahat}_1
  \;\le\; 2\inner{r}{\Btil v}_n + 2\pen\norm{\thetastar}_1 .
\end{equation}

\emph{Step 2, the two noises.} Since $\MA\Btil = \Btil$ and $\MA$ is
symmetric, $\inner{r}{\Btil v}_n = \inner{\bb+\bveps}{\Btil v}_n$. The
stochastic part is handled by duality: on $\mathcal{T}$,
$2|\inner{\bveps}{\Btil v}_n| \le
2\norm{\Btil\T\bveps/n}_{\infty}\norm{v}_1 \le \pen\norm{v}_1$. The
deterministic part is handled by the Cauchy-Schwarz and Young inequalities,
and not by duality, because $\bb$ is not small coordinatewise:
$2|\inner{\bb}{\Btil v}_n| \le 2\normn{\bb}\normn{\Btil v} \le
\normn{\Btil v}^2/2 + 2\normn{\bb}^2$. Substituting both in
\eqref{eq:basic},
\begin{equation}
  \label{eq:basic2}
  \tfrac12\normn{\Btil v}^2 + 2\pen\norm{\thetahat}_1
  \;\le\; \pen\norm{v}_1 + 2\pen\norm{\thetastar}_1 + 2\normn{\bb}^2 .
\end{equation}

\emph{Step 3, the slow rate.} Bounding
$\norm{v}_1 \le \norm{\thetahat}_1+\norm{\thetastar}_1$ in
\eqref{eq:basic2} gives
$\tfrac12\normn{\Btil v}^2 + \pen\norm{\thetahat}_1 \le
3\pen\norm{\thetastar}_1 + 2\normn{\bb}^2$, which is part (i). No property of
$\Gramhat$ was used.

\emph{Step 4, the cone.} Since $\thetastar$ vanishes off $\Szero$,
$\norm{\thetahat}_1 \ge \norm{\thetastar}_1 - \norm{v_{\Szero}}_1 +
\norm{v_{\Szero^c}}_1$ and
$\norm{v}_1 = \norm{v_{\Szero}}_1+\norm{v_{\Szero^c}}_1$; substituting in
\eqref{eq:basic2} and cancelling $2\pen\norm{\thetastar}_1$,
\begin{equation}
  \label{eq:cone}
  \tfrac12\normn{\Btil v}^2 + \pen\norm{v_{\Szero^c}}_1
  \;\le\; 3\pen\norm{v_{\Szero}}_1 + 2\normn{\bb}^2 .
\end{equation}

\emph{Step 5, two cases.} If $2\normn{\bb}^2 \le \pen\norm{v_{\Szero}}_1$,
then \eqref{eq:cone} gives
$\tfrac12\normn{\Btil v}^2 + \pen\norm{v_{\Szero^c}}_1 \le
4\pen\norm{v_{\Szero}}_1$; by the Cauchy-Schwarz inequality and
\eqref{eq:schur}, $\norm{v_{\Szero}}_1 \le \sqrt{s_0}\norm{v}_2 \le
\sqrt{s_0}\normn{\Btil v}/\sqrt{\gtil}$, so
$\normn{\Btil v} \le 8\pen\sqrt{s_0/\gtil}$ and
$\norm{v}_2^2 \le \normn{\Btil v}^2/\gtil \le 64\pen^2s_0/\gtil^2$; the same
display gives $\norm{v_{\Szero^c}}_1 \le 4\norm{v_{\Szero}}_1$ and hence
$\norm{v}_1 \le 5\norm{v_{\Szero}}_1 \le 40\pen s_0/\gtil$. If instead
$2\normn{\bb}^2 > \pen\norm{v_{\Szero}}_1$, then \eqref{eq:cone} gives
$\tfrac12\normn{\Btil v}^2 < 8\normn{\bb}^2$, that is
$\normn{\Btil v}^2 < 16\normn{\bb}^2$, together with
$\pen\norm{v}_1 < 10\normn{\bb}^2$ and
$\norm{v}_2^2 < 16\normn{\bb}^2/\gtil$. Each quantity is bounded by the
maximum of the two cases, hence by their sum, which is part (ii).

\emph{Step 6, prediction.} By \eqref{eq:decomp},
$\widehat{f}-f = \Btil v + \PA(\bb+\bveps) - \bb$, and $\PA$ being a
projection, $\normn{\PA\bb} \le \normn{\bb}$, which gives the triangle
inequality of part (iii). Conditionally on the covariates the errors are
independent, centered and of variance at most $\sigma^2$, so
\[
  \E(\normn{\PA\bveps}^2\mid\bX,\bU)
  = \frac1n\tr\{\PA\,\mathrm{Cov}(\bveps\mid\bX,\bU)\}
  \le \frac{\sigma^2}{n}\tr(\PA) \le \frac{\sigma^2p}{n},
\]
since the diagonal entries of $\PA$ are nonnegative and its trace is the rank
of $\Amat$. \qed

\begin{remark}[Why there is no cone condition]
\label{rem:nocone}
The usual treatment of unpenalized coordinates \citep[Section
6.9]{buhlmann2011statistics} places them inside the support and requires
compatibility on the cone $\norm{\delta_{S^c}}_1 \le 3\norm{\delta_S}_1$.
Profiling removes them from the problem, and \eqref{eq:schur} transfers the
full smallest eigenvalue of $\Gramhat$ to the residualized block, so that
Proposition 2(ii) is used as it stands and the constant does not depend on
$\Szero$. If only the restricted version of the design condition were
available, the cone would return, and with the constant $4$ rather than $3$,
since Step 5 produces $\norm{v_{\Szero^c}}_1 \le 4\norm{v_{\Szero}}_1$: the
bias consumes part of the slack.
\end{remark}

\begin{corollary}[The two rates with the design condition instantiated]
\label{cor:instant}
Add Assumptions 2 to 5 and suppose $2^J\log d/n \to 0$, which is implied by
the requirement of Proposition 2(ii). Take
$\pen = 2\sigma B_X\sqrt{2C_U}\{2\log(2d/\alpha)/n\}^{1/2}$. Then, with
probability at least $1-\alpha-\alpha'-\alpha''$ up to terms that tend to
zero,
\[
  \normn{\widehat{f}-f}^2 \;\lesssim\;
  \frac{\sigma^2B_X^2C_U}{\kappa_1c_U}\,\frac{s_0\log(pq\NJ)}{n}
  \;+\; \frac{\mathcal{B}^2_J}{\alpha'} \;+\; \frac{\sigma^2p}{n\alpha''},
\]
and, with no condition on the design, the same bound with
$\pen\norm{\thetastar}_1$ of Lemma S5.4 in place of the first term.
\end{corollary}

\begin{proof}
Intersect four events: $\mathcal{T}$, of probability at least $1-\alpha$, by
Lemma~\ref{lem:calib}; $\{\smax^2 \le 2B_X^2C_U\}$, by Lemma S5.3, which makes the stated
$\pen$ admissible, that is $\pen \ge 2\pen_0$;
$\{\lmin(\Gramhat)\ge\gamma/2\}$, by Proposition 2(ii), which with
\eqref{eq:schur} gives $\gtil \ge \gamma/2$; and
$\{\normn{\bb}^2 \le \mathcal{B}_J^2/\alpha'\}$, which is Markov's inequality
applied to $\E\normn{\bb}^2 = \norm{f-f_J}^2_{L_2(P)} \le \mathcal{B}_J^2$,
the identity holding because the observations are identically distributed.
Substituting in parts (ii) and (iii) of Theorem~\ref{thm:lasso-oracle} gives the statement, and
part (i) with Lemma S5.4 gives the unconditional version.
\end{proof}

The first lemma is Theorem~\ref{thm:lasso-oracle} read at an arbitrary comparator. It consumes no
new step: the proof of that theorem uses the oracle only through the
decomposition of $\bY$ and through the size of its support.

\begin{lemma}[Theorem~\ref{thm:lasso-oracle} with an arbitrary comparator]
\label{lem:comparator}
Under Assumption 1, with $\Amat$ of rank $p$ and $\pen \ge 2\pen_0$, on the
event $\mathcal{T}$ the following hold simultaneously for every
$\bar\bbeta = (\bar\bc,\thbar)$, with $\bbar = f - \bZ\bar\bbeta$,
$\bar M = |\supp(\thbar)|$ and $\bar v = \thetahat-\thbar$:
$\tfrac12\normn{\Btil\bar v}^2 + \pen\norm{\thetahat}_1 \le
3\pen\norm{\thbar}_1 + 2\normn{\bbar}^2$; if $\gtil>0$,
$\normn{\Btil\bar v}^2 \le 64\pen^2\bar M/\gtil + 16\normn{\bbar}^2$ and
$\norm{\bar v}_2^2 \le 64\pen^2\bar M/\gtil^2 +
16\normn{\bbar}^2/\gtil$; and
\begin{equation}
  \label{eq:oracle-min}
  \normn{\widehat{f}-f}^2 \;\le\;
  3\min_{\bar\bbeta}\Bigl\{\frac{64\pen^2\bar M}{\gtil}
  + 20\normn{\bbar}^2\Bigr\} + 3\normn{\PA\bveps}^2 .
\end{equation}
\end{lemma}

\begin{proof}
Follow the proof of Theorem~\ref{thm:lasso-oracle} with
$(\bar\bc,\thbar,\bbar,\supp(\thbar),\bar M,\bar v)$ in place of
$(\bc,\thetastar,\bb,\Szero,s_0,v)$. Step 1 uses only
$\bY = \Amat\bar\bc+\Bmat\thbar+\bbar+\bveps$, which is the definition of
$\bbar$, and $\MA\Amat = 0$. Step 2 uses $\mathcal{T}$, which does not depend
on the comparator, and $\normn{\MA\bbar} \le \normn{\bbar}$. Steps 3 to 5 are
algebra with $\supp(\thbar)$ in place of $\Szero$, using
$\thbar_{\supp(\thbar)^c} = 0$, the Cauchy-Schwarz inequality and
\eqref{eq:schur}. Step 6 gives
$\widehat{f}-\bZ\bar\bbeta = \Btil\bar v + \PA(\bbar+\bveps)$ and hence
$\widehat{f}-f = \Btil\bar v+\PA(\bbar+\bveps)-\bbar$. Since $\mathcal{T}$,
$\gtil$ and the condition on $\pen$ do not involve $\bar\bbeta$, the
inequalities hold simultaneously on the same event and the minimum may be
taken; \eqref{eq:oracle-min} follows from
$(x+y+z)^2 \le 3(x^2+y^2+z^2)$.
\end{proof}

\subsection*{Proof of Lemma~\ref{lem:lasso-weak}}

\emph{The envelope.} Fix $\epsilon>0$ and count
$N(\epsilon) = \#\{a : |\theta^{*}_a| > \epsilon\}$. Inside a block and a
level, Markov's inequality in the $\ell_\pi$ norm gives
$\#\{\tsl : |\coefs{\cv}{\md}{\lev}{\tsl}|>\epsilon\} \le
\epsilon^{-\pi}\norm{\theta^{*}_{\cv\md,\lev\cdot}}^{\pi}_{\pi}$, and
Assumption 5 gives
$\norm{\theta^{*}_{\cv\md,\lev\cdot}}^{\pi}_{\pi} \le
C_g^{\pi}2^{-\lev\pi(s+1/2-1/\pi)} = C_g^{\pi}2^{\lev(1-a)}$ with
$a = \pi(s+1/2)>1$. Since a level has $2^{\lev}$ coordinates,
\[
  N_{\cv\md}(\epsilon) \;\le\; \sum_{\lev\ge0}
  \min\{2^{\lev},\, C_g^{\pi}\epsilon^{-\pi}2^{\lev(1-a)}\} .
\]
The two arguments agree at $2^{\lev_\epsilon a} = (C_g/\epsilon)^{\pi}$, that
is at $2^{\lev_\epsilon} = (C_g/\epsilon)^{\tau}$, since
$\pi/a = (s+1/2)^{-1} = \tau$. Below that level the minimum is $2^{\lev}$ and
the sum is at most $2^{\lev_\epsilon}$; above it the minimum is the second
term, which decreases geometrically with ratio $2^{1-a}<1$, and the sum is
$2^{\lev_\epsilon}/(1-2^{1-a})$. Adding the two parts and the $pq$ blocks
gives $N(\epsilon) \le pqA_{s,\pi}(C_g/\epsilon)^{\tau}$; truncating at
$\lev<J$ only decreases the count, so the bound holds for every $J$. Finally,
$\epsilon < \theta^{*}_{(k)}$ implies $N(\epsilon) \ge k$, hence
$\epsilon \le C_g(pqA_{s,\pi}/k)^{1/\tau}$, and letting $\epsilon$ increase to
$\theta^{*}_{(k)}$ gives the envelope; and $\tau = (s+1/2)^{-1} < 2$ for every
$s>0$.

\emph{The tail.} With $2/\tau>1$,
$\sum_{k>M}\theta^2_{(k)} \le C_\tau^2\sum_{k>M}k^{-2/\tau} \le
C_\tau^2\int_M^{\infty}x^{-2/\tau}dx = C_\tau^2M^{1-2/\tau}\tau/(2-\tau)$.
\qed

\subsection*{Proof of Proposition~\ref{prop:lasso-slow}}

By parts (i) and (iii) of Theorem~\ref{thm:lasso-oracle}, on $\mathcal{T}$,
$\normn{\Btil v}^2 \le 6\pen\norm{\thetastar}_1+4\normn{\bb}^2$ and
$\normn{\widehat{f}-f}^2 \le 3(\normn{\Btil v}^2+4\normn{\bb}^2+
\normn{\PA\bveps}^2)$, and no property of $\Gramhat$ enters. Lemma S5.3 makes
$\penn$ admissible with probability tending to one, since
$2^{J_n}\log d_n/n \to 0$ in both regimes, and it is for this that the conditions
$\seff > 1/4$ in (i) and $\seff > s/2$ in (ii) are needed: without them the choice
in (i) gives $2^{J_n} \ge n^{1/(4\seff)} \ge n$, and the balance in (ii) gives
$2^{J_n} \gtrsim n/\log n$, and parts (iii) and (iv) cover those pairs with
$\penn^{+}$; Markov's inequality
applied to $\E\normn{\bb}^2 \le \mathcal{B}^2_{J_n}$, from Lemma 2, and to
$\E(\normn{\PA\bveps}^2\mid\bX,\bU) \le \sigma^2p/n$ gives the general
statement. If $s>1/2$, Lemma S5.4 gives $\norm{\thetastar}_1 = O(1)$ and
$\penn\norm{\thetastar}_1 \asymp (J_n/n)^{1/2}$, and it suffices that the bias
not dominate, that is $2^{-2J_n\seff} \le (J_n/n)^{1/2}$, which holds for
$J_n \ge \log_2(n/J_n)/(4\seff)$. If $s<1/2$, then
$\norm{\thetastar}_1 \asymp 2^{J_n(1/2-s)}$, the two terms are
$(J_n/n)^{1/2}2^{J_n(1/2-s)}$ and $2^{-2J_n\seff}$, and equating them gives
$2^{J_n} \asymp (n/J_n)^{1/(1-2s+4\seff)}$ with common value
$(J_n/n)^{2\seff/(1-2s+4\seff)}$; when $\seff = s$ the denominator is
$1+2s$.

% E5e (E1.14; D52): a prova dos itens (iii) e (iv) da Proposicao 4 de 05-taxas.tex.
\emph{The level above the ceiling, for (iii) and (iv).} By
\eqref{eq:smax-ceiling}, $\Prob(\smax^{2} > 2B_X^{2}C_U\mu_{J}) \le 1/d$ for
every $J$. Off that event,
\[
  2\pen_0 \;\le\; 2\sigma(2B_X^{2}C_U\mu_{J})^{1/2}
  \Bigl\{\frac{2\log(2d/\alpha)}{n}\Bigr\}^{1/2} \;=\; \penn^{+},
\]
and the argument of the general statement goes through with
$\penn^{+}$ and no condition on $J_n$: on $\mathcal{T}$,
$\normn{\widehat{f}-f}^{2} \le 18\penn^{+}\norm{\thetastar}_{1} +
24\normn{\bb}^{2} + 3\normn{\PA\bveps}^{2}$, and Markov's inequality
concludes. In every choice of (iii) and (iv), $2^{J_n} \le n^{2}$, so that
$\log d_n \asymp J_n \asymp \log n$ and
\begin{equation}
  \label{eq:lam-plus}
  \penn^{+} \;\asymp\; \Bigl(\frac{\log n}{n}\Bigr)^{1/2}
  \Bigl(1 + \frac{2^{J_n}\log n}{n}\Bigr)^{1/2},
\end{equation}
in particular $\penn^{+} \asymp 2^{J_n/2}\log n/n$ if
$2^{J_n} \gtrsim n/\log n$.

\emph{Part (iii)(a).} With $s > 1/2$, $\norm{\thetastar}_{1} = O(1)$ as in
(i). With $2^{J_n} \asymp (n/\log n)^{2/(1+4\seff)} \gtrsim n/\log n$,
because $\seff \le 1/4$, \eqref{eq:lam-plus} gives
$\penn^{+}\norm{\thetastar}_{1} \asymp (\log n/n)(n/\log n)^{1/(1+4\seff)} =
(\log n/n)^{4\seff/(1+4\seff)}$, and the bias is
$2^{-2J_n\seff} \asymp (\log n/n)^{4\seff/(1+4\seff)}$; $p/n$ is of smaller
order.

\emph{Part (iii)(b).} With $s < 1/2$, $\norm{\thetastar}_{1} \asymp
2^{J_n(1/2-s)}$ as in (ii). With
$2^{J_n} \asymp (n/\log n)^{1/(1-s+2\seff)} \gtrsim n/\log n$, because
$\seff \le s/2$, \eqref{eq:lam-plus} gives
$\penn^{+}\norm{\thetastar}_{1} \asymp 2^{J_n(1-s)}\log n/n \asymp
(\log n/n)^{1-(1-s)/(1-s+2\seff)} = (\log n/n)^{2\seff/(1-s+2\seff)}$, which
is also the order of the bias $2^{-2J_n\seff}$. At the boundaries,
$\seff = 1/4$ in (a) and $\seff = s/2$ in (b), $2^{J_n} \asymp n/\log n$,
$\mu_{J_n}$ is bounded, and the rates are $(\log n/n)^{1/2}$ and
$(\log n/n)^{s}$, those of (i) and (ii) at these points.

\emph{Part (iv).} With $s = 1/2$, Lemma S5.4 gives
$\norm{\thetastar}_{1} \le pqC_gJ_n$, and $J_n \asymp \log n$ in both
choices. If $\seff \ge 1/4$,
$2^{J_n} \asymp (n/\log n)^{1/(4\seff)} \lesssim n/\log n$, $\mu_{J_n}$ is
bounded, $\penn^{+}\norm{\thetastar}_{1} \asymp (\log n/n)^{1/2}\log n$,
and the bias, $2^{-2J_n\seff} \asymp (\log n/n)^{1/2}$, is of smaller order.
If $\seff < 1/4$, $2^{J_n} \asymp \{n/(\log n)^{2}\}^{2/(1+4\seff)}$, which
exceeds $n/\log n$ for $n$ large, and $\penn^{+}\norm{\thetastar}_{1} \asymp
2^{J_n/2}(\log n)^{2}/n \asymp \{(\log n)^{2}/n\}^{4\seff/(1+4\seff)}$, the
order of the bias. \qed

% derivations/05-taxas.tex, Observacao rem:teto (E1.14; E5e).
\begin{remark}[Why the resolution passes the ceiling]
\label{rem:ceiling}
In the pairs of (iii) the bound $\penn\norm{\thetastar}_{1} +
\mathcal{B}_{J}^{2}$ is dominated by the bias at every $J$ below the ceiling,
and the best that can be done there is to stop at it: $2^{J_n} = n/(K\log n)$
with $K > 4\norm{\psi}_{\infty}^{2}/(3C_U)$ fixed keeps $\penn$ admissible,
because the exponent of Lemma~\ref{lem:columns} at $t = B_X^{2}C_U$ is then
$-\{3C_UK/(4\norm{\psi}_{\infty}^{2})\}\log n$ and $d_n \le pqn$, and gives
$O_p\{(\log n/n)^{2\seff}\}$. Above the ceiling $\penn^{+}$ pays for the
growth of $\smax$, and the estimation term,
$2^{J/2}(\log n/n)\norm{\thetastar}_{1}$, grows more slowly than the bias
decreases until $J$ reaches the $J_n$ of (iii). The exponents of (iii),
$4\seff/(1+4\seff)$ and $2\seff/(1-s+2\seff)$, are strictly larger than
$2\seff$ away from $\seff = 1/4$ and $\seff = s/2$, and the $J_n$ of (iii)
minimizes the bound over $J$ up to a constant, so that this is the best bound
that the argument gives. In (iv) with $\seff \ge 1/4$, the factor
$(\log n)^{3/2}$ comes from $\norm{\thetastar}_{1} \asymp J_n \asymp \log n$,
which no choice avoids.
\end{remark}

% E5f (E1.15; D54): a prova da Proposicao 8 de derivations/05-taxas.tex.
\subsection*{Proof of Proposition~\ref{prop:lasso-truncated}}

\emph{Part (i).} Lemma~\ref{lem:comparator} with $\bar\bbeta = (\bc,\thbar)$
gives, discarding $\pen\norm{\thetahat}_{1} \ge 0$,
$\normn{\Btil\bar v}^{2} \le 6\pen\norm{\thbar}_{1} + 4\normn{\bbar}^{2}$,
and Step~6 in its proof gives $\normn{\widehat{f}-f} \le \normn{\Btil\bar v}
+ 2\normn{\bbar} + \normn{\PA\bveps}$; with
$(x+y+z)^{2} \le 3(x^{2}+y^{2}+z^{2})$, $\normn{\widehat{f}-f}^{2} \le
18\pen\norm{\thbar}_{1} + 24\normn{\bbar}^{2} + 3\normn{\PA\bveps}^{2}$.
Since $\bbar = \bb + \Bmat(\thetastar-\thbar)$, $\normn{\bbar}^{2} \le
2\normn{\bb}^{2} + 2\normn{\Bmat(\thetastar-\thbar)}^{2}$. If $t$ is
deterministic, so is $u = \thetastar-\thbar$, and
$\E\normn{\Bmat u}^{2} \le \lmaxx(\Gram)\norm{u}_{2}^{2} \le
\kappa_2C_U\norm{u}_{2}^{2}$, as in the proof of
Proposition~\ref{prop:truncated}. With $t = \pen$, on the coordinates with
$|\theta^{*}_{a}| > \pen$ the minimum in \eqref{eq:R1} is
$\pen|\theta^{*}_{a}|$, and on the others it is $\theta^{*2}_{a}$; summing
gives the bound by $\max(18,48\kappa_2C_U)\Rone(\thetastar;\pen)$. The form
in $O_p$ is that of the proof of Proposition~\ref{prop:lasso-slow}: the
event $\mathcal{T}$ with $\penn \ge 2\pen_0$, by Lemma~\ref{lem:columns}
below the ceiling, or with $\penn^{+} \ge 2\pen_0$, by
\eqref{eq:smax-ceiling} for every $J$, and Markov's inequality for
$\normn{\bb}^{2}$, $\normn{\Bmat(\thetastar-\thbar)}^{2}$ and
$\normn{\PA\bveps}^{2}$.

\emph{Part (ii), $s < 1/2$.} Fix a block and a level $\lev$, write
$\theta_{\lev} \in \R^{2^{\lev}}$ for its coefficients, with
$\norm{\theta_{\lev}}_{\pi} \le C_g2^{-\lev(s+1/2-1/\pi)}$ by Assumption~5,
and let $v_{\lev} = \sum_{\tsl}\min(t|\theta_{\lev\tsl}|,\theta_{\lev\tsl}^{2})$.
The first bound is in the $\ell_1$ norm: by H\"older's inequality in
$\R^{2^{\lev}}$, $v_{\lev} \le t\norm{\theta_{\lev}}_{1} \le
t2^{\lev(1-1/\pi)}\norm{\theta_{\lev}}_{\pi} \le tC_g2^{\lev(1/2-s)} =:
A_{\lev}$. The second is the tail, as in the proof of
Proposition~\ref{prop:truncated}(ii) with coordinates in place of chunks:
$v_{\lev} \le t^{2-\pi}\norm{\theta_{\lev}}_{\pi}^{\pi} \le
t^{2-\pi}C_g^{\pi}2^{-\lev\pi\seff}$ if $\pi \le 2$, and
$v_{\lev} \le \norm{\theta_{\lev}}_{2}^{2} \le C_g^{2}2^{-2\lev s}$ if
$\pi \ge 2$; call it $B_{\lev}$. With
\[
  2^{x} = (C_g/t)^{\tau},
  \qquad
  h = C_g^{\tau}t^{2-\tau} = C_g^{\tau}t^{4s/(2s+1)},
\]
one has $A_{\lev} = h2^{(\lev-x)(1/2-s)}$, because $1 + \tau(1/2-s) = \tau$,
and $B_{\lev} = h2^{-(\lev-x)\min(\pi,2)\seff}$, because
$\tau\pi\seff = \pi-\tau$ and $2s\tau = 2-\tau$. Summing $A_{\lev}$ over
$\lev \le x$ and $B_{\lev}$ over $\lev > x$, two geometric series of ratios
$2^{-(1/2-s)}$ and $2^{-\min(\pi,2)\seff}$ that start from at most $h$, gives
at most $h\{(1-2^{s-1/2})^{-1} + (1-2^{-\min(\pi,2)\seff})^{-1}\}$ per block;
truncating at $\lev < J$ only removes terms, and summing over the $pq$
blocks gives \eqref{eq:count-lasso-truncated}. This is the computation of
the proof of Lemma~3, with $\min(ty,y^{2})$ in place of $\min(y^{2},\eta)$.

\emph{Part (ii), $s \ge 1/2$.} The first statement is
$\min(t|y|,y^{2}) \le t|y|$. For the second, the block with
$\theta^{*}_{\cv\md,00} = C_g$ and all its other coefficients zero satisfies
Assumption~5, and gives $\Rone(\thetastar;t) \ge \min(tC_g,C_g^{2}) = tC_g$.
If $s = 1/2$, set, at each level $\lev < J$ with $2^{\lev} \le C_g/t$, all
the $2^{\lev}$ coefficients equal to $C_g2^{-\lev} \ge t$; then
$\norm{\theta_{\lev}}_{\pi} = C_g2^{-\lev(1-1/\pi)}$, the bound of
Assumption~5 at $s = 1/2$, and each such level contributes
$2^{\lev}tC_g2^{-\lev} = tC_g$.

\emph{Part (iii).} With $J_n = \lceil c\log_2 n\rceil$ and $c < 1$,
$2^{J_n}\log d_n/n \asymp n^{c-1}\log n \to 0$, and (i) holds with
$\penn \asymp (\log n/n)^{1/2}$. By \eqref{eq:count-lasso-truncated},
$\Rone(\thetastar;\penn) \lesssim (\log n/n)^{2s/(2s+1)}$. The lower end of
the window \eqref{eq:lasso-window} is $s/\{(2s+1)\seff\}$, because
$(2-\tau)/(4\seff) = s/\{(2s+1)\seff\}$ and $\seff \le s$, so that the bias
is $\mathcal{B}_{J_n}^{2} \lesssim 2^{-2J_n\seff} \le n^{-2c\seff} \le
n^{-2s/(2s+1)}$; and $p/n$ is of smaller order.

\emph{Part (iv).} Write $\eta = 2s/\{s+\seff(2s+1)\}$, which lies in $[1,2)$
exactly when $\seff \le s/(2s+1)$; then $2^{J_n} \asymp (n/\log n)^{\eta}
\gtrsim n/\log n$, $2^{J_n} \le n^{2}$ and, by \eqref{eq:lam-plus},
$\penn^{+} \asymp 2^{J_n/2}\log n/n \asymp (\log n/n)^{1-\eta/2}$. By
\eqref{eq:count-lasso-truncated}, $\Rone(\thetastar;\penn^{+}) \lesssim
(\log n/n)^{(1-\eta/2)4s/(2s+1)}$, and $(1-\eta/2)4s/(2s+1) =
4s\seff/\{s+\seff(2s+1)\}$; the bias is $2^{-2J_n\seff} \asymp
(\log n/n)^{2\seff\eta}$, with the same exponent. At $\seff = s/(2s+1)$,
$\eta = 1$ and the exponent is $2s/(2s+1)$. Outside the window of (iii),
$\pi < 2$, because $\pi \ge 2$ gives $\seff = s > s/(2s+1)$. \qed

\subsection*{Proof of Theorem~\ref{thm:lasso-linear}}

By the definition of $J_n$,
$(n/\log n)^{1/(2\seff+1)} \le 2^{J_n} < 2(n/\log n)^{1/(2\seff+1)}$ and
$J_n \asymp \log n$, so
\[
  \frac{pq2^{J_n}\log(pq2^{J_n})}{n}
  \;\asymp\; \frac{(n/\log n)^{1/(2\seff+1)}\log n}{n}
  \;=\; \Bigl(\frac{\log n}{n}\Bigr)^{2\seff/(2\seff+1)}
  \;\longrightarrow\; 0,
\]
which is the requirement of Proposition 2(ii); \eqref{eq:schur} then gives
$\gtil \ge \gamma/2$ with probability tending to one. On that event
Corollary~\ref{cor:instant} with $s_0 \le d_n$ gives
\[
  \normn{\widehat{f}-f}^2 \;\lesssim\;
  \frac{\sigma^2B_X^2C_U}{\kappa_1c_U}\frac{d_nL_n}{n}
  + \frac{\mathcal{B}^2_{J_n}}{\alpha'} + \frac{\sigma^2p}{n\alpha''} .
\]
With $d_n \asymp pq2^{J_n}$ and $L_n \asymp \log n$ the first term is the
display above; the second is
$\mathcal{B}^2_{J_n} \asymp 2^{-2J_n\seff} \le
(\log n/n)^{2\seff/(2\seff+1)}$ by the lower bound on $2^{J_n}$; and the third
is $O(1/n)$, of strictly smaller order because $2\seff/(2\seff+1)<1$. \qed

\subsection*{Proof of Corollary~\ref{cor:lasso-components}}

Inside a block the wavelets of level below $J$ are orthonormal in $L_2[0,1]$,
so
$\lVert\widehat{g}_{\cv\md}-\PiJ\gcomp{\cv}{\md}\rVert^2_{L_2[0,1]} =
\sum_{\lev\tsl}v^2_{\cv\md,\lev\tsl}$, and summing over the $pq$ blocks gives
$\norm{v}_2^2$; adding the bias of Lemma 2 in each block,
\[
  \sum_{\cv,\md}\lVert\widehat{g}_{\cv\md}-\gcomp{\cv}{\md}
  \rVert^2_{L_2[0,1]} \;\le\; 2\norm{v}_2^2
  + \frac{2pqC_g^2}{1-2^{-2\seff}}2^{-2J\seff} .
\]
Theorem~\ref{thm:lasso-oracle}(ii) bounds $\norm{v}_2^2$ by
$64\pen^2s_0/\gtil^2+16\normn{\bb}^2/\gtil$, and on the event of Theorem~\ref{thm:lasso-linear} the
quantity $\gtil$ is bounded below by the constant $\gamma/2$, so the two terms
are of the same order as those of Theorem~\ref{thm:lasso-linear}, the only difference being a
constant with $\gamma^{-2}$ in place of $\gamma^{-1}$. This is where the
design condition is used a second time. For the levels, Lemma S5.1 gives
$\Amat(\chat-\bc) = \PA(\bb+\bveps)-\PA\Bmat v$, so that
$\normn{\Amat(\chat-\bc)} \le \normn{\bb}+\normn{\PA\bveps}+
\lmaxx(\Gramhat)^{1/2}\norm{v}_2$, while $\Amat\T\Amat/n$ is a principal
submatrix of $\Gramhat$, whence
$\norm{\chat-\bc}_2^2 \le \normn{\Amat(\chat-\bc)}^2/(\gamma/2)$. Finally the
components of a functional coefficient depend on distinct coordinates of $u$
and the basis is orthonormal, so
$\lVert\widehat{\fcoef{\cv}}-\fcoef{\cv}\rVert_{L_2([0,1]^q)} \le
|\widehat{c}_{\cv}-\lvl{\cv}| + \sum_{\md}\lVert\widehat{g}_{\cv\md}-
\gcomp{\cv}{\md}\rVert_{L_2[0,1]}$, with $q$ fixed. The same argument applies
verbatim under the conditions of Theorem~\ref{thm:lasso-main}. \qed

\subsection*{Proof of Theorem~\ref{thm:lasso-main}}

\emph{A nonasymptotic bound.} Apply Lemma~\ref{lem:comparator} with
$\bar\bbeta = (\bc,\thstarM{M})$, the best $M$ term approximation of the
oracle. Then $\bar M \le M$ and
$\bbar = \bb + \Bmat(\thetastar-\thstarM{M})$, so that, by the
Cauchy-Schwarz inequality, Lemma~\ref{lem:lasso-weak} and \eqref{eq:lmax},
\[
  \normn{\bbar}^2 \;\le\; 2\normn{\bb}^2
  + 2\lmaxx(\Gramhat)\norm{\thetastar-\thstarM{M}}_2^2
  \;\le\; 2\normn{\bb}^2 + \frac{2\Lambda C_\tau^2\tau}{2-\tau}M^{1-2/\tau} .
\]
Substituting in \eqref{eq:oracle-min}, on the intersection of $\mathcal{T}$
with $\{\lmin(\Gramhat)\ge\gamma/2\}$,
\[
  \normn{\widehat{f}-f}^2 \;\le\; 3\min_{1\le M\le d}h(M)
  + 120\normn{\bb}^2 + 3\normn{\PA\bveps}^2,
  \qquad
  h(M) = \frac{64\pen^2}{\gtil}M
  + \frac{40\Lambda C_\tau^2\tau}{2-\tau}M^{1-2/\tau} .
\]
The function $h$ is strictly convex on the positive half line, since
$1-2/\tau<0$, and $h'$ vanishes at
$\Mstar = \{5\Lambda C_\tau^2\gtil/(8\pen^2)\}^{\tau/2}$, where
$h(\Mstar) = 2(2-\tau)^{-1}(64\pen^2/\gtil)\Mstar$; moreover
$h(\lceil\Mstar\rceil) \le h(\Mstar)+64\pen^2/\gtil \le 2h(\Mstar)$ whenever
$\Mstar \ge 1$, because $(2-\tau)/2<1$.

\emph{The resolution.} With $\penn^2 \asymp \sigma^2B_X^2C_UJ_n/n$,
$\gtil \in [\gamma/2,\Lambda]$ and the other constants fixed,
$\Mstar_n \asymp \penn^{-\tau} \asymp (n/J_n)^{\tau/2}$ and
$h(\Mstar_n) \asymp \penn^2\Mstar_n \asymp (J_n/n)^{(2-\tau)/2}$. Writing
$J_n = \lceil c\log_2 n\rceil$, so that $2^{J_n} \asymp n^{c}$, the three
inequalities of the window \eqref{eq:lasso-window} are, in order: $\Mstar_n \le d_n$, which asks
$(n/J_n)^{\tau/2} \lesssim pq\,n^{c}$, that is $c \ge \tau/2$ up to
logarithms; that the bias not dominate, $n^{-2c\seff} \lesssim
(J_n/n)^{(2-\tau)/2}$, that is $c \ge (2-\tau)/(4\seff)$; and the requirement
of Proposition 2(ii), $n^{c-1}\log n \to 0$, that is $c<1$. The term
$\sigma^2p/n$ is of smaller order because $(2-\tau)/2<1$. The window is
nonempty if and only if $\tau/2<1$, which always holds, and
$(2-\tau)/(4\seff)<1$, that is $\seff>(2-\tau)/4$.

\emph{Under Assumption 5 the hypothesis is automatic.} By Lemma~\ref{lem:lasso-weak} the
envelope holds with $\tau = (s+1/2)^{-1}$, whence $2-\tau = 4s/(2s+1)$,
$(2-\tau)/2 = 2s/(2s+1)$ and $(2-\tau)/4 = s/(2s+1)$, which is the condition
$\seff>s/(2s+1)$ of the statement; for $\pi \ge 2$ it holds for every $s>0$,
and for $\pi<2$ it reads $2s^2/(2s+1)>1/\pi-1/2$. The rate for the components
is the argument of Corollary~\ref{cor:lasso-components} with
$\norm{\thetahat-\thetastar}_2 \le \norm{\bar v}_2 +
\norm{\thetastar-\thstarM{M}}_2$ and Lemma~\ref{lem:comparator}, the only difference being an
extra factor $\gtil^{-1}$ in the first term, which does not change the order.
Finally $2s/(2s+1) \ge 2\seff/(2\seff+1)$, with equality only when
$\seff = s$, by monotonicity of $x \mapsto 2x/(2x+1)$ and $\seff \le s$, so
the rate is never worse than that of Theorem~\ref{thm:lasso-linear} and is strictly better when
$\pi<2$. \qed

\begin{remark}[What compressibility buys]
\label{rem:s0}
Read at the oracle, Theorem~\ref{thm:lasso-oracle}(ii) has $s_0 = d$ and delivers the rate of
\textcolor{colR1}{linear approximation}; read at a sparse comparator, through Lemma~\ref{lem:comparator}, it delivers the rate
of Theorem~\ref{thm:lasso-main}. It is the same inequality in both cases, and neither says
anything about $\supp(\thetahat)$: no minimal signal and no
irrepresentability condition is used, and what is gained is a rate governed by
the effective dimension $\Mstar_n$ rather than by $d_n$.
\end{remark}

\subsection*{Proof of Corollary~\ref{cor:lasso-threshold}}

\emph{A nonasymptotic bound on $\Delta$.} Work at $J = J_n$ and $\pen = \penn$,
and intersect the four events of the proof of Corollary~\ref{cor:instant}: $\mathcal{T}$,
of probability at least $1-\alpha$, by Lemma~\ref{lem:calib};
$\{\smax^2 \le 2B_X^2C_U\}$, by Lemma S5.3, which makes $\penn$ admissible;
$\{\lmin(\Gramhat) \ge \gamma/2\}$, by Proposition 2(ii), which with
\eqref{eq:schur} gives $\gtil \ge \gamma/2$; and
$\{\normn{\bb}^2 \le \mathcal{B}^2_{J_n}/\alpha'\}$, by Markov's inequality and
Lemma 2. On this event Theorem~\ref{thm:lasso-oracle}(ii) with $s_0 \le d_n$ gives
\[
  \norm{v}_2^2 \;\le\; \frac{64\penn^2d_n}{\gtil^2}
  + \frac{16\normn{\bb}^2}{\gtil}
  \;\le\; \frac{256\penn^2d_n}{\gamma^2}
  + \frac{32\mathcal{B}^2_{J_n}}{\alpha'\gamma} .
\]
For a fixed block, the orthonormality of the wavelets of level below $J_n$
gives $\lVert\widehat{g}_{\cv\md} - \PiJ\gcomp{\cv}{\md}\rVert_{L_2[0,1]} =
\norm{v_{\cv\md,\cdot}}_2 \le \norm{v}_2$, and the first display of Lemma 2
bounds $\lVert\PiJ\gcomp{\cv}{\md} - \gcomp{\cv}{\md}\rVert^2_{L_2[0,1]}$ by
$C_g^2(1-2^{-2\seff})^{-1}2^{-2J_n\seff}$. By $(x+y)^2 \le 2x^2 + 2y^2$,
\begin{equation}
  \label{eq:Dn}
  \Delta^2 \;\le\; \bar{\Delta}_n^2 \;=\; \frac{512\,\penn^2d_n}{\gamma^2}
  + \frac{64\,\mathcal{B}^2_{J_n}}{\alpha'\gamma}
  + \frac{2C_g^2}{1-2^{-2\seff}}\,2^{-2J_n\seff}
\end{equation}
with probability at least
$1 - \alpha - \alpha' - \Prob\{\lmin(\Gramhat) < \gamma/2\} -
\Prob\{\smax^2 > 2B_X^2C_U\}$, the last two terms tending to zero by
Proposition 2(ii) and Lemma S5.3\textcolor{colR1}{; their sum is the
$\omega_n$ of Corollary~\ref{cor:lasso-threshold}(iii), and this paragraph is the proof of that
part}. By Lemma S7.1(iii), on that event
$\widehat{\mathcal{S}}(t) = \mathcal{S}$ for every $t$ with $\bar{\Delta}_n \le t < \nu_{\min} - \bar{\Delta}_n$, a window
that is nonempty whenever $\nu_{\min} > 2\bar{\Delta}_n$. The bound $\bar{\Delta}_n$ depends on the
components only through the constants of Assumption 5, and the two residual
probabilities do not depend on them at all, so everything above holds
uniformly over the class, which is what allows the target to vary with $n$ in
Assumption 6.

\emph{The order of $\bar{\Delta}_n$.} At the resolution of Theorem~\ref{thm:lasso-linear} the three terms of
\eqref{eq:Dn} are at most of order $\rho_n^2$: the first because
$\penn^2d_n \asymp pq2^{J_n}J_n/n$, which is the quantity shown to be of order
$(\log n/n)^{2\seff/(2\seff+1)}$ in the proof of Theorem~\ref{thm:lasso-linear}, and the other two
because $\mathcal{B}^2_{J_n} \asymp 2^{-2J_n\seff} \le
(\log n/n)^{2\seff/(2\seff+1)}$ by the definition of $J_n$. Hence
$\bar{\Delta}_n \le c(\alpha,\alpha')\rho_n$ with $c(\alpha,\alpha') < \infty$. Given $\epsilon > 0$,
take $\alpha = \alpha' = \epsilon/4$ and $n$ large enough that the two
residual probabilities add up to less than $\epsilon/2$; then
$\Prob\{\Delta > c(\epsilon/4,\epsilon/4)\rho_n\} < \epsilon$, that is
$\Delta = O_p(\rho_n)$.

\emph{Part (i).} If $t_n/\rho_n \to \infty$, then
$\Prob(\Delta \le t_n) \to 1$, and on that event Lemma S7.1(i) with
$\bar{\Delta} = \Delta$ gives $\widehat{\mathcal{S}}(t_n) \subseteq \mathcal{S}$. No property of $\nu_{\min,n}$
has been used.

\emph{Part (ii).} If $\mathcal{S}$ is empty, part (ii) is part (i). Otherwise, the
condition $\limsup_n t_n/\nu_{\min,n} < 1$ gives $\eta > 0$ and $n_0$ with
$\nu_{\min,n} - t_n \ge \eta\,\nu_{\min,n}$ for $n \ge n_0$, and Assumption 6 with $\rho_n$ in place of $\rhoG$ gives
$(\nu_{\min,n} - t_n)/\rho_n \ge \eta\,\nu_{\min,n}/\rho_n \to \infty$. Hence
$\Prob(\Delta \le t_n) \to 1$ and $\Prob(\Delta < \nu_{\min,n} - t_n) \to 1$, and
on the intersection of the two events Lemma S7.1(iii) with $\bar{\Delta} = \Delta$ gives
$\widehat{\mathcal{S}}(t_n) = \mathcal{S}$. For the existence, $t_n = (\rho_n\nu_{\min,n})^{1/2}$ has
$t_n/\rho_n = (\nu_{\min,n}/\rho_n)^{1/2} \to \infty$ and
$t_n/\nu_{\min,n} = (\rho_n/\nu_{\min,n})^{1/2} \to 0$. \qed

\textcolor{colR1}{\emph{At the resolution of Theorem~\ref{thm:lasso-main}.} Theorem~\ref{thm:lasso-main} gives
$\sum_{\cv,\md}\lVert\widehat{g}_{\cv\md}-\gcomp{\cv}{\md}\rVert^2_{L_2[0,1]}
= O_p(\tilde\rho_n^{\,2})$ with $\tilde\rho_n = (\log n/n)^{s/(2s+1)}$, hence
$\Delta = O_p(\tilde\rho_n)$, and the bound in its proof depends on the
components only through the constants of Assumption 5, through $C_\tau$ of
Lemma~\ref{lem:lasso-weak} and the design constants, so that this holds uniformly over the
class. The proofs of parts (i) and (ii) above use nothing about $\Delta$
beyond $\Delta = O_p(\rho_n)$ uniformly over the class, together with
Lemma S7.1 and Assumption 6, and they go through verbatim with
$\tilde\rho_n$ in place of $\rho_n$.}
}

%%%%%%%%%%%%%%%%%%%%%%%%%%%%%%%%%%%%%%%%%%%%%%%%%%%%%%%%%%%%%%%%%%%%%%%%%%%%%%
\bibhang=1.7pc
\bibsep=2pt
\fontsize{9}{14pt plus.8pt minus .6pt}\selectfont
\renewcommand\bibname{\large \bf References}
\bibliographystyle{chicago}
\bibliography{references_4}

\end{document}
